% =====================================================================
%  R. Nielsen, FOLKELIGE FOREDRAG efter Indbydelse af Foredragsforeninger i Kolding, Fredericia og Aagaard.
%  Kjøbenhavn: Den Gyldendalske Boghandel (F. Hegel); trykt hos J. Jørgensen & Co., 1875.  164 pp.  Antiqua.
%
%  TRANSCRIPTION -- GENERATED; DO NOT EDIT BY HAND.  Rebuild from the repo root (split.py once; the second --apply places
%  the marks whose keys the first renumbers):
%    python3 ebook-method/folkelige-foredrag/split.py
%    python3 ebook-method/folkelige-foredrag/draft.py --emit --force \
%      && python3 ebook-method/collate/check.py folkelige-foredrag --italics --apply \
%      && python3 ebook-method/collate/check.py folkelige-foredrag --apply \
%      && python3 ebook-method/folkelige-foredrag/draft.py --post
%
%  STATUS: COMPLETE (2026-10-12): title page, the print's Indhold, and the whole text, pp. 1-164: nine lectures and the
%  Isted festival speech, with the book's one note (p. 85).
%
%  EDITION.  The text of record is Columbia University's MICROFILM of the book (master negative 93-81530-4), scanned by the
%  Internet Archive (folkeligeforedra00niel) and served by HathiTrust (nnc2.ark:/13960/t87h2xg7x) as a PDF with two book
%  pages to a frame: ~/bibliotek/Nielsen, Rasmus/1875-fokelige-foredrag.pdf (SCANS.tsv), 96 PDF pages, grey, ~370 ppi.
%  Its text layer is the Internet Archive's ABBYY OCR (the IA djvu text is the same, word for word), so the IA files are no
%  third reading.  split.py cuts each frame at its gutter into a derived one-page-per-page PDF (.parts/split.pdf); every
%  page number below is a printed page.  The only printing; no e-book, no other scan and no earlier transcription found
%  (checked 2026-10-12: the repo, the web, the Internet Archive).
%
%  ACCURACY (2026-10-12).  The words are those of the ABBYY layer (draft.py), collated word by word with an independent
%  tesseract reading of the images (collate/check.py, witness 'ocr'): 1018 disagreements in 41,162 words at the first
%  collation (twice the rate of the Google- and KB-layer books: ABBYY reads the film worse).  check.py's own rules settled
%  51; rules.py settled 281 more, each rule switched on only after a BLIND PILOT of 122 crops (reader A, 26/26 controls)
%  had measured it on this book: tesseract's glyph faults 8/8, aa read an 8/8, accents it adds 8/8, its » and „ read as a
%  comma 8/8 and 8/8, the layer's glyph faults 4/4, 8/8 and 8/8 (lengths within two), ABBYY's „see I“ for „see!“ 8/8, its
%  reject mark ^ 7/8 (narrowed), the printer's space before ; : ? ! 8/8, o/ø by frequency 7/8 (narrowed), tesseract's
%  specks and the note star 4/4, dashes 2/2; and three narrow parts of the layer's word splits and stray apostrophes
%  („gj ælder“, „som'“: 4/4, 1/1, 1/1).  The rest were read at the image BLIND (sheets.py: 577 word crops on 58 sheets,
%  two planted controls a sheet, 114 of 116 controls right -- the two misses crops cut at the margin; ten Sonnet
%  readers), and about 70 spots in the main session from line crops (zoom.py).  STRUCTURE read at the image: the ten
%  lecture heads and first lines (heads.py, reader N), the title page and Indhold; places where the layer put a word on
%  a line of its own (pp. 17, 19) or a line out of order (p. 102); the folio and film marks it read as text (pp. 17, 28).
%  LETTERSPACING (the print's emphasis, \emph): the white-gap measure is noise on grey film, so candidate lines were
%  taken by box WIDTH (spacedwidth.py, ratio 1.13: 138 lines, read blind, 20/20 controls; 58 runs), and twelve wide words
%  the readers left roman were checked at the image (zoom zg: none spaced).  ITALICS by stroke slant: none in the book.
%  SWEEPS after OPEN 0, each spot read at the image: --moved (alignment only), --splits and the split sweep; the
%  W-shortens / T-longer list of reads.py's decisions (found one W that had cut „op-|dage“); the reader-line sweep
%  (linecheck.py: found „Fri-“ lost inside p. 69 and two misprints below); the per-paragraph quote balance (a speck read as
%  “, p. 82; the print's own open quotations, pp. 48, 82); paragraph starts against tesseract's indents (one false break,
%  p. 84); the unknown-word sweep (798 words in no other transcription, 172 not compounds).  SAMPLED PAGE CHECK
%  (pagecheck.py, seed 202610125): pp. 4, 82, 120 and the title page read whole and blind, 703 words: 0 word and 0
%  punctuation errors; one emphasis error (a spaced „Jeg“, p. 82), mended.
%  TRANSLATION as a last reading: the ten translators listed 77 odd spots (ebook-method/folkelige-foredrag/tr/odd-*.tsv),
%  ~40 read at the image (zoom zo, zp): mended by POST patch -- specks both OCRs read as a stop or apostrophe (pp. 58, 67,
%  69), thin word spaces (pp. 50, 118, 122, 150), „absolnte“ and „o. d. 1.“ (both OCRs: pp. 51, 18), the printer's spaces
%  before ; and ? that a W decision had taken from tesseract (pp. 30, 161), two false paragraph breaks at page turns left
%  by the junk removal (pp. 28, 32) -- and four f's whose head did not ink (pp. 52, 72, 143, 157), at first taken for
%  misprints; the rest are the print's.
%  What this cannot see: a misprint both OCRs normalise into a real word -- the sweeps find only non-words and nonsense.
%  PRINTED AS IS (each seen at the image; a % comment at the site): no stop after „Jorden“ (p. 7), „en“ for er (p. 10),
%  „Mennneskelige“ (p. 10), „hiælpe“ (p. 29), „Bekymring.:“ (p. 37) and „er.:“ (p. 52), „de.“ before „Stridende“ (p. 47),
%  „Øvrigked“ (p. 83), „deri.“ (p. 84) and „stor.“ (p. 93), a stop or a comma whose tail did not ink, „bvilken“ (p. 87),
%  „o, s. v.“ (p. 130), „Brystet. og“ (p. 144) and „ind. at“ (p. 149), stops before a lower-case word, „gjennnemtrængt“
%  (p. 153); the accents of stress „sóm“ (p. 94), „Dén“ (p. 27), „dén“ (p. 108); also, unmarked, „Midie“ (p. 145),
%  „iskænkes“ (p. 31), „1848“ for the Constitution of 1849 (p. 154), and the open quotations of pp. 48 („Himlen.““) and
%  82 (Holberg's quotation is not closed).
%  PRINTER'S DEFECTS (letters supplied, a % comment at the site): four f's whose head did not ink, read t by both OCRs --
%  „Folkefriheden“ (p. 52), „fornemme“ (p. 72), „Uforstand“ (p. 143), „for“ (p. 157); „afskrækkende“ (p. 61, a broken f
%  read í); „Skoleforsømmelse“ (p. 83, a broken o read c); a speck over the a of „spøgelseagtigt“ (p. 3) is not an accent.
%  The ø rule (Hans 2026-10-08): the draft's o/ø faults were settled by the collation (oefreq, layerfault).
%  Not transcribed: folios, signature marks, the library's stamp, film marks; the large initials of the lecture openings
%  are set as plain letters; the rule under each heading is drawn by \lecture.  Quotation marks as printed („ “); the
%  layer's .. ., ,. and '' for „ and “, and tesseract's straight quotes, are mapped to them.  The note is marked *).
%
%  PAGE MAP.  split page = printed + 23 (p. 1 = split 24, the right half of frame 12) up to p. 129; the film repeats the
%  frame with pp. 130-131 (frames 77 and 78): the retake, frame 78, is used, so p. 130 = split 155 and printed + 25 from
%  there to p. 164 = split 189.  Title page split 20, Indhold split 22.  Each lecture begins on a new page: I p. 1, II p. 18,
%  III p. 35, IV p. 54, V p. 70, VI p. 86, VII p. 103, VIII p. 120, IX p. 139, the Festtale p. 156.
%
%  TOOLS (ebook-method/folkelige-foredrag/, copied from om-samfundets-vilkaar): split.py (the two-up frames), draft.py,
%  ocr.py, rules.py (with --classify for the pilot), sheets.py + reads.py + READER-BRIEF.md, heads.py, spacedwidth.py +
%  spacedlines.py + spaced.tsv, zoom.py, pagecheck.py + PAGECHECK-BRIEF.md, patches.py (PATCHES, POST), decisions.tsv,
%  and the sweeps splitsweep.py, linecheck.py, parcheck.py, unknown.py.
%
% =====================================================================
\documentclass[12pt,a4paper,openany]{book}
\usepackage[utf8]{inputenc}
\usepackage[T1]{fontenc}
\usepackage{libertinus}
\usepackage{libertinust1math}
\usepackage[danish]{babel}
\usepackage[protrusion=true,expansion=true]{microtype}
\usepackage{setspace}
\usepackage[margin=1.3in]{geometry}
\usepackage{fancyhdr}
\usepackage{hyperref}
\hypersetup{pdftitle={Folkelige Foredrag}, pdfauthor={Rasmus Nielsen}, hidelinks}
\pagestyle{fancy}
\fancyhf{}
\fancyhead[LE,RO]{\thepage}
\fancyhead[C]{\textit{R. Nielsen: Folkelige Foredrag (1875)}}
\setstretch{1.3}
\usepackage{marginnote}
\renewcommand{\thefootnote}{\ifcase\value{footnote}\or*)\or**)\or***)\fi}   % the print marks a page's notes *) and **); the counter restarts at each page
\newcommand{\tocentry}[3]{\noindent\hangindent=2em #1 #2\dotfill\ \textbf{#3}\par}
\newcommand{\opage}[1]{\marginnote{\footnotesize\textit{[#1]}}}
% a lecture opens a new page: its heading ("Første Foredrag,") and place ("holdt i Kolding.", in smaller bold), centred, then a short
% rule; the PDF outline takes each lecture's title from the Indhold (\lecturetitleN: the print has none at the head); the large initial is not reproduced.
\newcounter{lect}
\expandafter\def\csname lecturetitle1\endcsname{Tidernes Vilkaar: et gammelt Grundlovsbrud}
\expandafter\def\csname lecturetitle2\endcsname{Christus som Folkenes Lærer og Fører}
\expandafter\def\csname lecturetitle3\endcsname{Hvad Christus lærer om den historiske Fremtid}
\expandafter\def\csname lecturetitle4\endcsname{Nutiden: Lys og Skygge}
\expandafter\def\csname lecturetitle5\endcsname{Kirke og Religion: Magtsprog og Myndighed}
\expandafter\def\csname lecturetitle6\endcsname{„Gudsdomme“ i Historien}
\expandafter\def\csname lecturetitle7\endcsname{Den danske Folkekirkes Ordning ved Lov}
\expandafter\def\csname lecturetitle8\endcsname{Universitet og Folkehøiskole}
\expandafter\def\csname lecturetitle9\endcsname{Et lille Folk}
\expandafter\def\csname lecturetitle10\endcsname{Til Minde om Istedslaget, en Festtale}
\newcommand{\lecture}[2]{\clearpage\stepcounter{lect}\phantomsection\addcontentsline{toc}{chapter}{#1 \csname lecturetitle\arabic{lect}\endcsname}\begin{center}{\large #1}\\[0.4em]{\normalsize\bfseries #2}\\[0.4em]\rule{3em}{0.4pt}\end{center}\medskip\noindent\ignorespaces}

\begin{document}

% --- front matter, read at the image (front20, front22: reader N, 2026-10-12): split page 20, the title page; 22, Indhold.
% The round stamp on the title page is the library's (Columbia), not print: not transcribed.  The title page's lines are set
% in display faces of different widths (condensed, extended), not letterspaced: no \emph there (zoom zh, page check pc).
\begin{titlepage}
\begin{center}
\vspace*{2em}
{\Huge FOLKELIGE FOREDRAG.}\\[1.5em]
{\small EFTER INDBYDELSE AF FOREDRAGSFORENINGER I KOLDING,}\\[0.3em]
{\small FREDERICIA OG AAGAARD}\\[1.2em]
{\small HOLDTE}\\[0.6em]
{\small AF}\\[0.6em]
{\large R. NIELSEN.}\\[5em]
{\large\bfseries KJØBENHAVN.}\\[0.5em]
FORLAGT AF DEN GYLDENDALSKE BOGHANDEL (F. HEGEL).\\[0.5em]
{\small TRYKT HOS J. JØRGENSEN \& CO.}\\[0.5em]
{\small 1875.}
\end{center}
\end{titlepage}

% Indhold (split page 22), as printed: „Andet Foredrag,“ with a comma, the others with a stop; the dot leaders are set by \dotfill.
\begin{center}{\large Indhold.}\\[0.4em]\rule{3em}{0.4pt}\end{center}
\hfill{\small\textbf{Side}}\par
\tocentry{Første Foredrag.}{Tidernes Vilkaar: et gammelt Grundlovsbrud}{1}
\tocentry{Andet Foredrag,}{Christus som Folkenes Lærer og Fører}{18}
\tocentry{Tredie Foredrag.}{Hvad Christus lærer om den historiske Fremtid}{35}
\tocentry{Fjerde Foredrag.}{Nutiden: Lys og Skygge}{54}
\tocentry{Femte Foredrag.}{Kirke og Religion: Magtsprog og Myndighed}{70}
\tocentry{Sjette Foredrag.}{„Gudsdomme“ i Historien}{86}
\tocentry{Syvende Foredrag.}{Den danske Folkekirkes Ordning ved Lov}{103}
\tocentry{Ottende Foredrag.}{Universitet og Folkehøiskole}{120}
\tocentry{Niende Foredrag.}{Et lille Folk}{139}
\tocentry{}{Til Minde om Istedslaget, en Festtale}{156}
\clearpage


% ===== TEXT (pp. 1--164) =====
% --- p. 1 ---
\setcounter{footnote}{0}%
\lecture{Første Foredrag,} {holdt i Kolding.}
\opage{1}Den nærværende Tid er fuld af Spørgsmaal, thi Verden
er fuld af Tvivl og Uro. Er det end stille paa Overfladen,
saa koger det dog i Dybet; Stilheden er som et
Havblik. De Mægtige tale om Fred, men see! de staae
væbnede indtil Tænderne, de maale hinanden med
Øinene: et rynket Bryn, et Trusels Ord, en graalig Sky
paa Magtens Pande aftegner sig strax som et draget
Sværd, et Varsel, der vækker Angst, og sætter Sind og
Tanker i Røre. Anelserne tynge, Spændingen voxer;
den ene Forklaring tyer hen til den anden, det ene
Spørgsmaal kalder paa det andet: hvorfra kommer det?
hvad sigter det til? hvor fører det hen?

Vi see paa det, som ligger os nærmest; men dette
os Nærmeste viser hen til Noget, som er længere borte,
dette igjen til Noget, som endnu er fjernere. Hvorfra
udgaaer det første Stød? Hvor er Bevægelsernes Udspring,
det store Drivhjul, der holder det Hele i Gang?
Vi bæres af en mægtig Strøm, men hvor munder den
ud? Hvor er vort Maal, vor Ankerplads, vor trygge
Havn, naar Stormene komme?

At give en alsidig, fuldstændig, i alle Maader tilforladelig
Besvarelse af saa store Spørgsmaal, over%
% --- p. 2 ---
\setcounter{footnote}{0}%
\opage{2}gaaer menneskelige Kræfter, ingen Dødelig skal formaste
sig til at raade saa mørke Runer. Dog, uden
Veiledning ere vi ikke. Der er i alt det Uvisse Noget,
som er ubedragelig vist. Historien er en opslagen
Bog; vi have et aabenbaret Sandhedsord til at lyse paa
vor Vei: der er Ledestjerner paa Aandens Himmel.
Enhver, som har seet i den dybe Flod og derhos givet Agt
paa de ledende Stjerner, har en Antydning af Bevægelsernes
Retning og vil fra sit Stade kunne yde,
var det end kun et ringe, Bidrag til Spørgsmaalenes
Besvarelse. Jeg skal bestræbe mig for at yde et Bidrag.
Hvad jeg har at meddele, angaaer ikke ligefrem
Øieblikkets brændende Spørgsmaal, ikke saadanne
praktiske Anliggender, som ellers i Folkeforsamlinger
drøftes og bør drøftes af indsigtsfulde, sagkyndige
Mænd, Mænd, som med uhildet Blik for det virkelige
Hele og for alt det Enkelte, kjende de foreliggende
Midler og de nærmeste Veie, som fører til Maalet.
Udenfor --- men dog ingenlunde ligegyldig for --- de
politiske Bevægelser i Nutiden, er jeg ifølge min Livsretning
nærmest henviist til en beskuende Opfattelse,
en Opstilling af almene Synspunkter, en Iagttagelse af
de store Love, hvorefter Bevægelserne foregaae. Jeg
skal fremsætte min Opfattelse saa jævnt og klart, saa
aabent og frit som det er mig muligt. Hvad jeg da
ved denne Anledning har at udtale, anbefaler jeg til
den ærede Forsamlings velvillige Opmærksomhed og
alvorlige Prøvelse.

Prædikeren siger, at Alt er Forfængelighed, „Forfængeligheders
Forfængelighed. En Slægt gaaer og en
Slægt kommer, men Jorden staaer evindelig. Solen
gaaer op og Solen gaaer ned og iler hen til det Sted,
% --- p. 3 ---
\setcounter{footnote}{0}%
\opage{3}hvorfra den skal gaae op igjen. Alle Bække løbe til
Havet, men Havet bliver ikke fuldt; thi det Sted, hvor
Floderne gaae hen, derfra gaae de atter tilbage. Alle
Ord afmattes; Ingen kan udsige det, Øiet bliver ikke
træt af at see og Øret ikke fyldt af at høre. Hvad
der har været, det vil være, hvad der er skeet, det vil
skee, og der er intet Nyt under Solen.“

Dersom dette var Forklaringen; dersom de skiftende
Tider intet Nyt og Eiendommeligt bragte, men
det Samme og det Samme i uendelige Gjentagelser;
dersom det Nærværende kun var et Gjenfærd fra Forbigangenhedens
Grave, hvor huult og spøgelseagtigt % defect: a speck over the a („spøgelseágtigt“), not an accent
vilde det mangefarvede, rigt bevægede Menneskeliv da
ikke være! Det er Prædikerens sidste Tanke, at \emph{Gudsfrygt
er Viisdommens Begyndelse}; tages Gudsfrygtens
Tanke bort, da er det Hele et Op og Ned.
Til at begrunde denne korte, fyndige Tanke er den
opstillede Betragtning af Kredsløbet at forstaae som
en vidtløftig Indledning; det er den falmede Stormandserfaring,
den mølædte Verdensklogskab, den livstrætte
Viisdom, der taler.

„Hvad der har været, det vil være; hvad der er
skeet, det vil skee;“ det er uimodsigelig sandt, saa
længe Talen kun er om Tilværelsens almindelige Love.
Der er og der har altid været Synd og Forkrænkelighed
med i Slægternes Levnetsløb; der er og der har altid
været en høist betænkelig, her og der oprørende Splid
mellem Mægtige og Ringe, Rige og Fattige, Nydende
og Lidende, Bedrøvede og Glade. Hvad der er skeet,
det vil skee, at de Stærke overvælde de Svage, at der
beiles til Magten, at dens Tyngdepunkt snart er foroven
og snart forneden, snart hos nogle Faa og snart
% --- p. 4 ---
\setcounter{footnote}{0}%
\opage{4}hos Mængden. Det er skeet og det vil skee, at de,
der komme til Magten, misbruge den, enten i den Enkeltes
eller i Mængdens Navn, hvis ingen Modvægt opstilles.
Det er skeet og det vil skee, at Menneskene
misunde hverandre Magtens Brug i samme Grad som
de smigre de Magthavende. Er Magten hos den Enkelte,
og den Enkelte enevældig, saa bliver den Enevældige
hadet som en Djævel og hilset som en Gud.
Er Magten hos Mængden, saa er Fleertalsrøsten Guds
Røst, saa er Klogskaben paa Færde med at udspeide
Hobens Lidenskaber, benytte dens dunkle Instinkter og
bilde de Uvidende ind, at de virkelig have en Villie:
saa smigrer og kryber man for Mængden. Er det kun
dette Syn, hvorved der skal dvæles, da har Prædikeren
Ret: Forfængeligheders Forfængelighed; der er intet
Nyt under Solen.

Men Synet bør ikke standse ved Bevægelsernes
bølgende Overflade, ved den eensformige Gyngen op og
ned eller ved Skyggerne paa den brede Flod, som
forbiilende Skyer kaste; dybere i Bunden staaer et
Billede, Viisdommens Billede, Fingerpeg af et Forsyn,
som styrer Bevægelserne hen til et Maal. Hver Tidsalder
har, naar vi see nærmere til, sine Arbeider, sine
særegne Opgaver og faaer i Forhold hertil sit Præg og
sin eiendommelige Charakteer. Det skeer saaledes, at
een stor Bestræbelse faaer Magt og optager de andre
i sig, een Livstanke hæver sig og behersker de mange
Tanker, een Villiesretning arbeider sig frem og faaer
Overhaand over de mange Villier, i denne Livstanke
og den tilsvarende Villie har Styrelsen indlagt et
Maal.

I Middelalderen var det Kirken, som blev det
% --- p. 5 ---
\setcounter{footnote}{0}%
\opage{5}høieste Styrende. Menneskene kunde vende sig, hvorhen
de vilde, forsøge det med hvadsomhelst, privat
eller offenlig, i Familielivet eller det borgerlige Samfund,
de mødte overalt den kirkelige Myndigheds aarvaagne
Øie: kun hvad der fik Kirkens Velsignelse,
kunde voxe og trives; hvad der blev ramt af dens
Forbandelse, maatte visne og døe. Middelalderen var
Kirkens Tidsalder.

Mellem Reformationen og den franske Revolution
(1517---1789) ligger et Tidsrum, i hvilken Samfundsforholdene
efterhaanden omdannedes paa en saadan
Maade, at den verdslige Magt fik afgjort Overvægt
over den geistlige. Det var Staten med sin Enevælde,
som under Form af det absolute Kongedømme udstrakte
sit Scepter over Høie og Lave. Fra Kronen
udstraalede Livets Lys; Smaaurterne fandt Læ, og de
høie Planter voxte i Thronens Skygge. Det var intet
Paafund af Enkeltmands Vilkaarlighed, men en naturlig
Følge af Tingenes Gang. Staten var Maal, alt Andet
kun Middel. De Store nød deres Forrettigheder, deri
var ingen Uret; det var jo som Statens Støtter, for
Statens Skyld, til Statens Forherligelse. Mængden
trællede og sukkede, deri var ingen Uret; det var jo i
Statens Tjeneste, for Statens Skyld, til det Heles Bedste.
Saaledes er Statens Enevælde modsat Kirkens; Ingen
kan tage feil af de to Tidsaldere.

Men nu vor egen Tidsalder, den Nutid, hvori vi
leve, hvad Charakteer har den, og hvilket Navn vil
den bære? Den er endnu i sin Begyndelse, i sit tidlige
Foraar. Men i det tidlige Foraar er Veirliget
gjerne ustadigt: Solskin vexler med Rusk og Slud.
Mange Sædekorn og mange Slags Frø ligge skjulte i
% --- p. 6 ---
\setcounter{footnote}{0}%
\opage{6}Jorden, og Ingen kan see, hvad det er, før det ret har
spiret og er kommet op. Adskilligt af det, der kommer
først op, er end ikke altid det Bedste, thi hvor Jordbunden
ikke er dyrket og renset, voxer Ukrudtet nok
saa frodig som Kornet. Ikke desto mindre har Tidsalderen
allerede saa faste og eiendommelige Træk, at
vi ikke kunne være i Tvivl om dens Charakteer og
Væsen. Dens Fødsel og Herkomst er ingen Tilfældighed;
den er sin Fortids ægte Barn, født af Tilskikkelsernes
Moderskjød og døbt i Begivenhedernes Strøm. Det var
i 48, et bevæget Aar, at Traaden brast, og den Nyfødte
rørte sig kraftig i Svøbet. Paa en Uveirsdag
blev den holdt over Daaben; sorte, lynsvangre Skyer
formørkede Himlen, men Skyerne skiltes, og Solen brød
frem, da Tiden fik det nye, dog ældgamle Navn: \emph{Folkets
Tid, Folkenes Tid!} Det er store Forventninger,
rige Forhaabninger, der knytte sig til dette
Navn; maatte den unge Tid dog fatte sin Opgaves
Storhed; maatte den holde sit Vaabenskjold uplettet
og blankt og blive sin Sendelse tro; maatte dens Folkedaab
være mere end en Vanddaab, maatte den gjøre
sit Navn Ære!

Der melder sig strax en Vanskelighed, naar Talen
er om Folkets Tid, thi hvorpaa skal den egentlig
kjendes? Naar den høieste Styrelse har sit Sæde i
Kirken, da er det let at sige, hvad Kirken er. Den
er som Staden paa Bjerget, en storartet Foranstaltning
med eiendommelig Lære og Love. Den stiller sig til
Skue som en guddommelig Tingenes Orden; den kjender
sine Venner og den kjender sine Fjender: det kan
nøiagtig angives, hvorvidt dens Herredømme naaer.
Det Samme gjælder om Staten, naar den har faaet sin
% --- p. 7 ---
\setcounter{footnote}{0}%
\opage{7}rette Skikkelse som eneraadende Magt paa Jorden % print: no stop after „Jorden“
Men at paavise nøiagtig, hvor Folket er, og hvad det
er, at gjøre tilbørlig Forskjel mellem dets Venner og
Fjender, er ingenlunde nogen let Opgave. Hvor er
Folket? I hele Befolkningens Mængde. Mængden maa
med, naar Folket skal med; og dog er Mængden ikke
Folket: Folket som Folk staaer høit over Mængden.
Mængden maa aves og styres; men Folket maa styre
sig selv og sin Mængde. Uden folkelig Selvstyrelse
intet sandt Folkeliv; men hvorledes er en folkelig Selvstyrelse
mulig?

Under Kirkens Enevælde kunde Folkelivet ikke
trives; thi Selvstyrelsen var udelukkende i Kirkens
Haand. Kirken alene havde at byde, Folkene maatte
troe og lyde; dermed var Sagen afgjort. I en saadan
Tilstand af blind Underkastelse maatte Folket, idetmindste
deelviis, synke ned til blot at være Mængde.
Ligesaa under Statens Enevælde. Det er i denne Henseende
lærerigt at mærke sig den Lighed, der er imellem
de to Stormagter, Kirken og Staten, hvad Opfattelsen
af Folket angaaer.

Naar Mængden, den store Hob, synker i Knæ for
Paven, troer paa hans Mirakler og gaaer andægtig med
i Processionen, saa er det ikke Mængden, der knæler,
nei, det er Folket, det fromme, troende Folk; og Kirken
er folkelig, den elsker Folket. Men vaagner der i
Folket en dybere religiøs Trang, rører der sig i dets
Sind en Følelse af, at det Indre er Maalet, at det
Guddommelige vanhelliges ved udvortes Gøglerier: da
bliver Kirkens Øie mørkt; det seer da slet ikke Folket,
det seer kun Mængden, den vantro Hob, fordømte
Kjættere, som skulle udryddes.

% --- p. 8 ---
\setcounter{footnote}{0}%
\opage{8}Naar, i Statens Tidsalder, Herskeren viste sig
offenlig, og Hoben raabte Hurra, og Mængden jublede,
da var det ikke Hoben, ikke Mængden, der raabte
Hurra, nei! det var de gode Borgere, det kjære, trofaste
Folk; og Eneherskeren var folkelig, han elskede
jo sit Folk. Men vaagnede der i Folket et Spørgsmaal
om Ret, en Følelse af Undertrykkelse og oprørende
Ulighed, og der hørtes uhyggelig Mumlen, da maatte
Politiet frem. Det var naturligviis ikke Folket, der
mumlede, det var kun Mængden, den raa Hob, og
Mængden, den raa Hob, skulde tugtes.

Var det endda blot saadanne selvstyrende Tugtemestere
som den enevældige Kirke eller den enevældige
Stat, Folket havde imod sig, vilde det for længst være
voxet fra Riset; men det har en Modstander at bekæmpe,
der er langt farligere, fordi han ikke blot er
udvortes, men ogsaa indvortes. Han melder sig ikke
som Fjende, men som Ven og Broder; hans Væsen er
saa naturlig sammenvoxet med Folkets Tilværelse, at
han hvad Øieblik, Leilighed gives, kan træde op og
sige: det er mig, der er Folket. Denne falske Ven,
denne farlige Broder er --- \emph{Mængden.} At ville lade
Mængden styre eller lægge an paa at indføre folkelig
Selvstyrelse efter Mængdens Sind, er det Urimeligste
af Alt, thi Mængden kan ikke styre sig selv: dens
Væsen er at være ustyrlig. Man lægge vel Mærke til,
at Mængden akkurat har samme Forestilling om Ret
og Magt, som det meest oprørende Enkeltmandstyrannie.
Samme Feltraab, samme Løsen; Vi ere stærke, thi vi
ere mange; Løsenet er: \emph{den Stærkeres Ret.}

Det er, under dette Synspunkt, en sørgelig Belæring,
Historien giver. Hver Gang en Folkevillie er i
% --- p. 9 ---
\setcounter{footnote}{0}%
\opage{9}Færd med at danne sig, en Villie, der ligelig omfatter
Alle, Smaa og Store, en Villie, der vil det Gode og
Rette, fordi den i Sandhed vil det ægte Menneskelige,
saa er der strax en Underfundighed med i Spillet, et
listigt Bedrag, som forstyrrer det Hele. Snart er det
en begavet Stormand, som i Folkets Navn faaer Mængden
med sig og bringer Orden paa sin Viis ved at
tage Løvens Part og øve Vold imod Folket. Snart er
det Mængden, som under Afsyngning af Frihedshymner
sønderriver Ordenens Baand og tilsidst, trods Raab og
Skrig paa Lov og Ret, ender med en Uorden, hvoraf
Tyranniet fødes. Det er en Skæbne, som Aartusinder
igjennem har hvilet tungt paa Folk og Slægter. Hvoraf
kan det vel komme? Af gammel Synd, af indre
Splid mellem Villie og Villie, mellem Naturvillie, egenraadig
Driftvillie, og Aandsvillie.

Der er et gammelt Mellemværende mellem Aanden
og Slægten, et Grundlovsbrud fra Arildstid, et Oprør
mod Aandens Love, hvormed Historien begynder.
Verdensbegivenhederne danne et Væv, som skjuler derover;
men det aabenbarede Ord drager Sløret tilside
og viser os, hvad Aanden vil. Der er i Tidernes Morgen
holdt Dom saa vel over Stormandsegoismen som
over Mængden.

I første Mosebogs sjette Kapitel berettes Følgende.
„Og det skete, da Menneskene begyndte at formere sig
paa Jordens Overflade, og der fødtes dem Døttre, da
saae Guds Sønner Menneskenes Døttre, at de vare
skjønne, og de toge sig Hustruer af Alle, hvilke de
valgte. Kæmperne vare paa Jorden i de Dage, og ligeledes
siden, efter at Guds Sønner vare gaaede ind til
% --- p. 10 ---
\setcounter{footnote}{0}%
\opage{10}Menneskenes Døttre. Disse ere de Stærke, som fra
Arilds Tid vare navnkundige Mænd.“

Der en nu dem, som have meent, % print: „en“ (misprint for er)
 at heri var noget
Forblommet og Mythisk, noget Uklart og Hedensk, at
Gudssønnerne (Elohimssønnerne) enten maatte være
Sønner af lavere Guder eller maaskee faldne Engle,
Dæmoner. Disse Grublerier ere overflødige; hvad Ordet
aabenbarer, er tydeligt nok. De Sønner og Døttre,
hvorom her tales, ere af fælles Afstamning, alle, naturlig
forstaaet, Børn af samme Slægt, alle, aandelig forstaaet,
Guds Børn, skabte i Guds Billede. Adam kaldes jo
Guds Søn (Luk. 3, 38); de Mægtige paa Jorden, Fyrster
og Dommere, kaldes Guder. „Gud holder Dom
iblandt Guder“ (Psalm. 82, 1). „Jeg sagde: I ere
Guder, alle tilhobe ere I den Høiestes Børn, men nu
skulle I døe som Mennesker“ (Psalm. 82, 6—7; Joh.
10, 34---35). Den samme Menneskenatur bliver her
fremstillet fra to Sider: i Mandens Styrke sees det
Guddommelige, i Kvindens Skjønhed det Mennneskelige. % print: „Mennneskelige“ (misprint)
Stærke Mænd og skjønne Kvinder, det er Slægten med
dens rige Udstyr, en kraftig Slægt, som ikke Arbeidets
Møisommelighed eller Trældommens Aag endnu har
bøiet.

Det Næste og mere Overraskende er det Blik,
hvormed Jehova seer paa denne Slægt; thi det er ikke
i Mildhed og Naade, men i Fortørnelse og Vrede.
Heri har man da især villet skimte en Rest af Hedenskabet,
en mythisk, endnu halvt afgudisk Opfattelse af
Aandens Væsen. Jehova, siger man, er ifølge den hele
Fortælling aabenbart blevet misundelig og bange. Han
er misundelig paa denne overvættes Fylde af Kraft og
Skjønhed; han er bange for at disse navnkundige Mænd,
% --- p. 11 ---
\setcounter{footnote}{0}%
\opage{11}disse vældige Kæmper skulle --- ligerviis som Titanerne
i den græske Mythe --- forsøge paa at storme hans
Himmel og kuldkaste hans Throne. Det er et lille
Exempel paa Nutidens formeentlig høiere Viden om
det Guddommelige, men for Resten en uhyre Misforstaaelse.
Den Aandens Vrede, hvorom her tales, gjælder
ikke de rige Naturens Gaver, hvormed Skaberen
selv fra Begyndelsen af har udstyret Menneskeslægten;
nei, den gjælder disse Gavers Misbrug. Jehova seer,
at den rigt udstyrede Slægt, der nu skal prøve sine
Kræfter og betræde Frihedens Vei, ikke vil bruge
Gaverne i Aandens Tjeneste; han seer disse Kæmper
give sig hen i selviske Naturtilskyndelser, han seer dem
vildt arbeide paa deres egen Navnkundighed og bruge
den Stærkeres Ret; og Dommerens Stemme høres:
„Ikke skal min Aand herske iblandt Menneskene evindelig,
eftersom han er Kjød, og hans Dage skulle være
eet hundrede og tyve Aar.“

Naar Aandskræfterne hos de Begavede forvandles
til Kjødkræfter, naar en ubændig Egenvillie kun anvender
disse Kræfter til at mætte Rovbegjærligheden,
Ærgjerrighedens Hunger og Nydelsens Lyst, saa vender
Aanden sig med Væmmelse bort fra de Stærke; deres
Grændse er sat, deres Dage ere talte, deres Dom er
fældet. Det er forfærdelige Ord, Aabenbaringens Ord om
Kæmpernes Æt med de store Bedrifter: „Jehova saae,
at Menneskets Ondskab var stor paa Jorden, og alt
hans Hjertes Paafund ikkun ondt hver Dag.“ Aandens
hellige Vrede kan ikke skildres med stærkere Farver
end, naar det siges: „det angrede Jehova, at han havde
dannet Mennesket paa Jorden, og det fortrød ham i
hans Hjerte.“ Striden mellem Aandsvillien og det
% --- p. 12 ---
\setcounter{footnote}{0}%
\opage{12}naturkraftige Menneskes egenraadige Driftvillie er allerede
her paa sit Høieste. Vi kunne være forvissede
om, at hvad Aanden alt ved Historiens Begyndelse har
dømt om den brutale Navnkundighed, der vindes ved
Vold og Uret, det Samme dømmer den endnu den Dag
idag. Ligesom hine første Kæmpeforsøg, ved hvilke
Jorden blev „fordærvet for Guds Ansigt“ og „var fuld
af Vold,“ løste sig op og skyllede bort i Syndfloden,
saaledes har Historien i sit fortsatte Løb været en
Ødelæggelsens Flod for alle Storværker, som ikke vare
byggede paa Sandhed og Ret, og ikke havde Aandens
Indvielse.

Og nu Mængdens Storværk. Stedet findes i første
Mosebogs ellevte Kapitel. „Og hele Jorden var eet
Tungemaal og eet Sprog. Og det skete, der Folkene
droge fra Østen, da fandt de en Dal i Sinears Land,
og de boede der. Og de sagde, den Ene til den Anden,
velan, lader os stryge Teglsteen, og brænde en Brænding;
og Tegl tjente dem som Steen og Jordbeg som
Kalk. Og de sagde: velan! lader os bygge en Stad og et
Taarn, og dets Spidse skal naae til Himlen, og lader
os gjøre os et Navn, at vi ikke skulle adspredes over
Jordens Overflade.“

Der er i denne Begivenhed Noget, som ved første
Øiekast maa vække vor høieste Forundring. Hvad
er det vel for Folk, hvorom her tales? „Noahs Sønners
Slægter;“ altsaa Jehovas, Aandens, eget Folk, thi Jehova
havde jo gjort en Pagt med Noah. Hvad Ondt kan
der da være deri, at disse Mange give sig til at arbeide
i Forening og bygge paa en Stad og et Taarn? Her
er da ingen Kæmper, som slaae ned iblandt Dværge,
ingen Voldsmænd, som forkue de Svage, ingen Slave%
% --- p. 13 ---
\setcounter{footnote}{0}%
\opage{13}fogeder, som pidske de Trælbundne til at samle Straa
og stryge Tegl. Tvertimod! Alt gaaer saa ligelig, saa
flinkt og fredelig, at den uhyre Mængde af travle Hænder
gjør Indtryk af en stor, samdrægtig Arbeiderforening,
i hvilken den Ene uophørlig opmuntrer den Anden med
et: velan, velan! Det er jo et skjønt Syn, en storartet
Anvendelse af Folkekræfter, som arbeide, ikke paa
Adsplittelse, men paa Enighed og broderligt Sammenhold:
hvad kan herimod indvendes? Og dog er Aanden
vred; Jehova stiger ned og forvirrer de Arbeidende
og forstyrrer deres Arbeide. Hvorledes er det muligt,
at Forvirring og Forstyrrelse kan udgaae fra Aanden?
„Og Jehova sagde: see, det er eet Folk og eet Tungemaal
er for Alle og dette have de begyndt at gjøre, og
nu skal Intet kunne formenes dem af hvad de tænke
paa at gjøre.“ Saa dunkle, ja meningsløse som disse
Ord vilde være, dersom vi opfattede dem paa en saa
barnagtig Maade, at den Almægtige virkelig skulde
være tilsinds at ville adsplitte Folkene og svække
Folkevillien af Frygt for det høie Taarn; saa klare, og
forstaaelige ere de, naar vi fatte dem som Udtryk for
Sandheden i den alstyrende Villies Tanke. Taarnet,
hvis Spidse skulde naae til Himlen, var ingenlunde
bestemt til at gjennembryde Himmelhvælvingen; det
var et Storværk, som alene skulde have Betydning paa
Jorden, et fremragende Foreningspunkt, som ved at
trodse Tidernes Omskiftelse skulde være et talende
Vidnesbyrd om, hvad Menneskekræfter formaae.

De arbeidende Folkefærd tænke ved denne Anledning
slet ikke paa enten at fornærme Himlen eller
trodsig udfordre Aanden; deres Tanke er meget mere
den, at de bedst kunne gjøre det Hele selv; saa lade
% --- p. 14 ---
\setcounter{footnote}{0}%
\opage{14}de Himlen være Himmel, og Aanden Aand, og tage fat
paa Arbeidet. Ældgammelt og dog som fra igaar!
Aanden er, om vi saa maae sige, ikke fortørnet paa
sine egne Vegne, fordi den bliver overseet af Menneskene
og ikke nyder den Ære at komme med. Det er
for Menneskenes Skyld, paa Folkets og Folkelivets
Vegne, at Aanden vredes, thi uden Aandens Medvirken
er Folkelivet umuligt.

De travle Folk, der bygge i Sinears Dal, have
givet Afkald paa, hvad der i dybere Forstand kan
minde om Aandens Krav, og staae i Begreb med at
synke ned til en slet og ret Mængde. De have glemt
det Oprindelige, deres Stamfaders Forhold til Jehova
og det Pagtens Tegn, der blev sat i Skyen. Hvad er
vel en Regnbue? Et flygtigt Lysskjær uden Virkelighed,
uden Paalidelighed; saaledes forholder det sig, efter
deres Mening, ogsaa med Aanden. Nei, da er et Taarn
noget ganske anderledes Virkeligt og Varigt; og saa
er det desuden deres egne Hænders Gjerning, det
udødelige Værk, hvorved de „gjøre sig et Navn“. Jehovas
Ord, hans Almagtstegn og hans styrende Viisdom,
det vidste man ikke, hvad var; det var Noget man
slet ikke kunde faae fat paa, det var altsammen flygtigt
og luftigt som Regnbuen. Teglsteen, Leer og
Jordbeg, det vidste man, hvad var, det var Ting af
Værdi, Noget man kunde føle og tage paa. Saaledes
har da denne mylrende Mængde med sin Travlhed, sin
Arbeidsdygtighed og kolossale Stræben fornegtet Menneskelivets
dybeste Forudsætninger, brudt med Aanden,
brudt den hellige Pagt med Jehova.

Dette er Livets og Frihedens Vilkaar: fornegtes
det Oprindelige, faaer Aanden ikke Lov at komme med
% --- p. 15 ---
\setcounter{footnote}{0}%
\opage{15}i levende Tro, i indre Tilegnelse og fri Lydighed, da
samler den Usynlige sin Vredes Sky over den aandløse
Selvraadigheds Myretue, forstyrre dens Værk og
spreder, hvad den har samlet. Det er umuligt, at der
i et menneskeligt Samfund kan være sand Enighed
uden dybere Enighedsgrund; men Enighedens dybere
Grund er en Villiesgrund. Enighedens Grund er i
Erkjendelse af det Gode og Rette; den sande afgjørende
Villie er Aandsvillie. Hvor Grunden og Villien mangler,
er Enigheden indbildt og forbigaaende; Sindene
skifte, Driftvillien yttrer sig i vexlende Udslag af selviske
Tilbøieligheder og Lidenskaber: Uenigheden kommer
pludselig bag paa de Enige, naar de mindst vente det.

Paa Enighedens Grund og Villiens Art skal man
kjende Forskjellen mellem Folket og Mængden. Dersom
den store Plan med det store Taarn i Sinears
Dal virkelig havde havt et sandt, i Menneskevæsenet
dybt rodfæstet Villieskrav, da maatte vi unegtelig indrømme,
at her var mere end en Mængde, at de mange
Folkefærd vare eet Folk, et driftigt, dygtigt, opadstræbende
Folk, som arbeidede paa noget Ophøiet og Ædelt;
den Gud, som af vilkaarligt Lune vilde forstyrre et
ædelt Menneskeværk, maatte vi reentud erklære for en
Afgud. Men de babyloniske Taarnbyggere have ingen
saadan Enighedsgrund; havde de en saadan Grund,
saa maatte de jo nødvendig havt Aanden med, men
det er netop Aanden, der er udelukket. Det er ikke
Folket, det er Mængden, som her har slaaet sig sammen
om at fuldføre et udvortes Overordenligt, som ret
kunde falde i Øinene. Lad os see Sagen naturligt og
spørge: hvorledes skulde Enigheden om et saa uhyre
Fællesforetagende vel kunne tænkes at holde sig? Var
% --- p. 16 ---
\setcounter{footnote}{0}%
\opage{16}der ingen Ledere, saa var der heller ingen Byggeplan;
Enhver maatte gaae efter sit eget Hoved, og det kunde
da umulig føre til Enighed. Var der Ledere, Overbestyrere,
saa maatte disse jo enten anvende Tvang
eller Overtalelse. Tvang kunde de ikke anvende, thi
naar Mængden skal under Tvang, saa giver den sig
øieblikkelig ud for Folket, og saa hedder det strax, at
der øves Vold imod Folket. Lederne have altsaa
maattet indskrænke sig til at anvende Overtalelse. Nu
er jo vist nok Intet lettere end at overtale Mængden
til at gjøre Noget, som den i Grunden har Lyst til,
Noget, hvortil den mener at føle en indre uimodstaaelig
Trang; men at overtale Mængden til at gjøre Noget,
hvortil den ikke har Lyst, det er det Vanskeligste af
Alt. Det er altsaa, kan man tænke sig, da Arbeidet
begyndte at blive haardt og Taarnet skævt i Sinears
Dal, kommet til en alvorlig Opposition imod Overbestyrelsen,
til Forhandlinger saa brogede, brede og
langtrukne, at Sprogene forvirredes, og Talerne tilsidst
ikke forstode hverandre.

Og hvad er det saa Aanden gjør? Den fremskynder
Afgjørelsen og fælder Dommen. Det er dens Villie, at
den i Egenraadighedens Væsen skjulte Splid skal blive
aabenbar; den gjør det vitterligt, af hvad Art et saadant
Sammenhold, en saadan Enighed er. „Velan,
lader os stige ned og forstyrre deres Tungemaal, saa
at den Ene ikke skal forstaae den Andens Tungemaal.
--- Og Jehova adspredte dem over hele Jordens Overflade,
og de lode af at bygge Staden.“

Vi lære heraf, at det tidlig er kommet til et Brud
mellem Aand og Folk. Som Følge af dette Brud
maatte de mange Folkefærd med de mange Tungemaal
% --- p. 17 ---
\setcounter{footnote}{0}%
\opage{17}da vandre hver sine Veie --- til Seire eller Nederlag, Ophøielse
eller Fornedrelse, Hæder eller Vanære. Aandsvillien
er ikke skeet Fyldest; der er Vrede i Aandens
Blik, og dog er Aanden Kjærlighed: der er Lysning af
Aand i Folkenes Naturgrund. Lysskjæret er forskjelligt
i Forhold til de enkelte Folkeslags særlige Begavelse
og ejendommelige Anlæg; men der rører sig noget
Guddommeligt hos dem alle. De hige, de kæmpe, de
famle og snuble; men Bevægelsen gaaer med ustandselig
Fart --- bestandig fremad! Paa tusinde Skær,
af Stormandsvælde, af Ordenstyrannie og den vilde
Mængdes Had til Orden, er Folkeskibet hidtil strandet:
Mindets Kyster ere bedækkede med Vrag. Det
er den gamle Brøde, det gamle Grundlovsbrud, som
atter og atter har hævnet sig. Altid den Stærkeres
Ret! „Men saa maa jo Retfærdigheden selv tilsidst
blive den Stærkere“; det er den inderste Tanke
i Folkenes Sind. Det er dette Syn --- en Regnbue
i Skyen --- hvorefter de hige; det er det Virkelighedsmaal,
hvortil de stræbe hen. Efter tusindaarige
Kampes Forvirringer, Lidelser og Nederlag er omsider
et Vendepunkt indtraadt. Oplysningens Sol er i vore
Dage kommet saa høit paa Himlen, at vi kunne see,
hvor vi ere. Vi kunne endnu ikke see Kampens Udgang,
langt fra! men vi kunne dog see Maalet. Skal
Maalet naaes, da maa der Villie til. Det er ikke Enkeltmandsvillie,
ikke Mængdevillie eller Fleertalsvillie,
hvorom det gjælder; nei, det er Folkevillie. Men den
sande Folkevillie er uadskillelig fra Aandsvillie; skal
Folket seire i den haarde Kamp og naae sit Maal, da
maa det erfare Indvielsens Kraft: det maa døbes med
Aand.

% --- p. 18 ---
\setcounter{footnote}{0}%
\lecture{Andet Foredrag,} {holdt i Fredericia.}
\opage{18}Det læres i den hellige Skrift, at Menneskeslægten
boer paa den ganske Jordens Kreds af eet Blod. Men
det berettes jo ogsaa, hvorledes det gik til, at den ene
Slægt blev til mange Slægter, den ene Folkefamilie til
mange Folkefærd, det ene oprindelige, fælles Sprog til
mange forskjellige Tungemaal. Adsplittelsen og Forskjelligheden er nu saa stor, at Forskerne i vor videnskabelige
Tid maae begynde Undersøgelserne forfra og
forelægge sig selv det Spørgsmaal: er hele Menneskeslægten
virkelig af een Herkomst? Kunne alle disse
hvide, sorte, olivengule og kobberrøde Mennesker antages
at nedstamme fra eet Par Forældre? Uligheden
er mere end udvortes, den indskrænker sig ikke til
Huden og Haaret, til Hjerneskallens og Ansigtets Form,
Øiets Regnbuehinde o. d. l.; den er lige saa stor i det
Indvortes som i det Udvortes. Nogle Folkestammer
ere vidunderlig begavede, andre ikke; nogle udmærke
sig ved en høi Grad af Oplysning og Dannelse, andre
befinde sig endnu paa Raahedens laveste Trin. Selv
iblandt de Begavede er Uligheden stor, Inderne, Ægypterne,
Grækerne, Romerne og Nordboerne have jo havt
% --- p. 19 ---
\setcounter{footnote}{0}%
\opage{19}hver sin Aandsretning, hver sit eiendommelige Syn paa
Livet. Det sees især af deres Guder, thi det er Folkeguder;
disse Folkeguder vare saa at sige, Ihændehavere
af de aandelige Kapitaler, paa hvilke Folkene tærede.

Hvilken Forskjel mellem to saadanne Gudekredse
som den græske og den nordiske; med alt det Beslægtede
er der, hvad Dybde og Inderlighed angaaer,
en betydelig Afstand imellem dem: Odin forstaaer ikke
Zeus, Aphrodite og Freja kunne ikke leve sammen under
eet Tag. Disse Guder ere Folkeeiendommelighedens
oprindelige Bærere; naar de ældes, blegne og synke,
hensvinder Folkelivets Kraft. Dersom nu det vi kalde
Folke-Aand kun var en taaget Indbildning, et Væv af
uklare Rørelser og stærke Drømme: da kunde vi jo vel
forklare os Guderne som de første lyse Skikkelser, i
hvilke det menneskelige Væsen, det med Tanker, Anelser
og halvforstaaede Længsler overfyldte Sind, fandt
en farverig, mangfoldig Afspeiling af sit eget bevægede
Indre; men saa maatte vi da ogsaa med det Samme
nedsætte Folke-Aanderne til forsvindende Tidsaander.
Enten maatte vi da antage, at Folkelivet i Nutiden
skulde kunne trives og blomstre uden Aand, eller
forvente, at den gamle Folketro med de gamle
Guder skulde staae op fra de Døde og fortrænge den
nye Oplysning. Det Første er lige saa urimeligt som
det Sidste. Vi maae da betragte Folke-Aanderne fra
Hedenold som aandeagtige Magter, spredte Glimt og
dunkle Forkyndelser af den ene evige Sandhedsaand,
hvis Lys aldrig har ophørt at skinne i Mørket.

Naar Lyset gaaer igjennem Lag af Skyer, blive
Straalerne brudte og Farverne mange; saaledes ogsaa,
naar det aandelige Lys skal bane sig Vei igjennem
% --- p. 20 ---
\setcounter{footnote}{0}%
\opage{20}Folkenes uligeartede Naturgrund. Vi kunne meget vel
tillægge Folke-Aanderne Gyldighed uden at gjøre dem
til Personer eller stille dem op som selvstændige Væsener.
Sandhedsaanden vil ikke udslette Folkenes Eiendommelighed,
men den vil fjerne det Fremmede, det Fjendtlige,
og efterhaanden som den faaer Lov at raade, adle
det Særegne og forherlige det Eiendommelige. Sandhedens
Aand er personlig i den sande Religion; dens
Menneske-Aabenbarelse er Christus. For at see, hvorledes
Christi Aand forholder sig til Folke-Aand, ville
vi i dette Foredrag betragte Christus og Folket.

Christus og Folket: dette Forhold er ikke det
samme som: Christus og Menigheden. For Menighedens
Tro staaer Faderens Eenbaarne som Gudsaabenbareren
i Kjød og Blod, Forsoneren, den Korsfæstede, der led
og døde for Verdens Synder, opstod af Graven og indgik
til Herlighed; i et saa inderligt Forhold staaer
Christus ikke til Folket som Folk. Det folkelige Øie
seer det Guddommelige i ham fra den menneskelige
Side, betragter ham som en heroisk Skikkelse, en fremragende,
enestaaende Charakteer i Verdenshistorien, og
dvæler ved hans offenlige Liv. I sit offenlige Liv viser
Christus sig som Folkets Mand, dets Fører, Lærer og
Leder. Hvad han lærer Folket, hvad han dømmer om
Folket, hvad han vil gjøre for Folket og hvad han vil,
at Folket skal gjøre, bliver altsaa fra dette Synspunkt
Hovedsagen.

Strax ved sin første Optræden udsender Christus
et Lys, der adskiller Mængden og Folket; Mængden
søger ham, og i Mængden søger han Folket. Hans
Tale paa Bjerget er en Tale til Mængden og dog i
allerhøieste Forstand en Folketale. Her er Noget, som
% --- p. 21 ---
\setcounter{footnote}{0}%
\opage{21}skiller Folketaleren i Galilæa fra Talerne i Athen
og Rom. De store Talere i Verdens Hovedstæder begyndte
saa gjerne med at minde de Forsamlede om
deres naturlige Rettigheder og Fordringer. „Atheniensere! --- Kviriter! Frie Borgere, fribaarne Mænd, I
vide jo, at I selv have Magt til at give Love; I vide
og føle, at det er Eder uværdigt at adlyde andre Love
end dem, I selv have givet.“ Ved slige Taler kom
Lidenskaberne i Røre, og Mængden fik Mod. Forstaaer
sig, der blev ikke ligefrem talt om Mængden, der
blev udtrykkelig, med en fuldt tonende Veltalenheds
hele Klang og Styrke, talt til Folket om Folket, ene
og alene om Folket. Ikke desto mindre var der noget
Betænkeligt bagved. Der var noget Politisk, for ikke
at sige Polisk, i den hele Fremgangsmaade. Det saae
livagtig ud i Mængdens Øine, som om Friheden kom
med et Spring, som sprang den ud af Talerens Hoved,
ligerviis som den voxne, fuldt rustede Minerva fordum
sprang ud af Jupiters Hoved. Saa slog Mængden sig
løs, men Folket var og blev en Fange, det skiftede
kun Ledebaand: det lød paa Frihed, men endte med
slavisk Afhængighed.

Folketaleren i Galilæa gaaer anderledes tilværks.
For at frigjøre Folket, begynder han med at binde
Mængden. „Mener ikke, at jeg er kommen for at opløse
Loven eller Propheterne; jeg er ikke kommen at
opløse, men at fuldkomme. Thi sandelig siger jeg
Eder: indtil Himmelen og Jorden forgaae, skal end
ikke det mindste Bogstav eller en Tøddel forgaae af
Loven, indtil det skeer altsammen.“ Idet han saaledes
taler fra Bjerget, taler han fra Aandens, fra Evighedens
Høie --- ikke om borgerlige Love, der kunne omdannes
% --- p. 22 ---
\setcounter{footnote}{0}%
\opage{22}og skifte, men om den uforanderlige Lov, den hellige
Grundlov, som er skreven i Menneskenes Hjerter, den
Sandhedens, Retfærdighedens og Kjærlighedens Lov,
der bestaaer, om end Himmel og Jord forgaae. Og
han skærper denne Lov og gjør den endnu strengere.
„I have hørt, hvad der er sagt til de Gamle; men jeg
siger Eder“ . . . . Det er ikke nok med en Sum af
udvortes Bud; det kommer væsenlig an paa det Indvortes.
Idet han saaledes indprenter de lyttende Skarer,
hvad det er, der oprindelig fordres, en indre Forandring,
en gjennemgaaende Sindets Omdannelse: ikke
at ville, fordi der er Tvang, men oprigtig at ville, fordi
det er en Hjertets Trang, bliver Strengheden selv til
Mildhed. Saaledes forvandler han Loven fra en Bogstavlov
til en Aandslov, fra en Tvangs- til en Frihedslov,
og hjælper Folket til Frihed.

Man skulde befrygte, at Mængden vilde knurre og
føle sig krænket over en saa dristig Tale af en enkelt
værgeløs Mand --- thi at Christus ikke den Gang kunde
staae for de Manges Øine i det Lys, hvori han nu
staaer for os, er en Selvfølge. Han stod iblandt dem
som En, der alene vilde bedømmes efter sine Ord og
Gjerninger. Men han forstod at gribe Samvittighederne
og forvandle Mængden til Folk. Indtrykket af Talen
er --- betegnende nok --- gjengivet med de Ord: „Folket
forundrede sig saare over hans Lærdom, thi, han
lærte dem som den, der havde Myndighed, ikke som
de Skriftkloge.“ Som den, der havde Myndighed, hvor
besynderligt! Havde de Skriftkloge da ingen Myndighed?
Efter det Udvortes at dømme skulde det jo udelukkende
være dem, der havde Myndigheden; de vare jo Landets
Øvrighed, Lovens Fortolkere og Haandhævere, Folkets
% --- p. 23 ---
\setcounter{footnote}{0}%
\opage{23}beskikkede Formyndere. Men deres Myndighed var kun
et Skin. Folket havde en Fornemmelse af, at disse
Skriftkloge med al deres Ordkløveri og Lovlærdom
uden Livskraft, vare uden dybere Sands for Sandhed,
uden Aandsoverlegenhed. Det hele Bestaaende var et
tomt Udvortes, noget Udlevet, der holdt sig ved Former
og Vedtægter.

For at gjenoprette den sande Myndighed, maa
Christus begynde fra Grunden af med at nedrive den
falske. Han gjør det som Folkets Leder uden endelig
Hensyntagen, uden Omsvøb, uden Skaansel. Han
river Masken af Hykleriet; i hele Folkets Paahør raaber
han Vee over „de kalkede Grave“. Hans Ord
kunde falde som Himlens Dug, husvalende som milde
Lufte, der vifte om Bjergets Kind; men naar Frækheden
saarede ham, og den hellige Vrede brød frem, da var
Talen en Torden, og Lynene knittrede. „Paa Mose
Stol sidde Skriftkloge og Pharisæer. Alt det derfor,
som de sige Eder, at I skulle holde, det holder og
gjører; men gjører ikke efter deres Gjerninger, thi de
sige det vel, men gjøre det ikke. Thi de binde svare
Byrder, vanskelige at bære, og lægge Menneskene dem
paa Skuldrene, men selv ville de ikke røre dem med
en Finger.“

Vi kunne forstaae, at de Skriftkloge og Pharisæerne,
Ordenens, Besindighedens og den gamle Gudfrygtigheds
Mænd, maatte hade den nye Prophet med et grundigt
Had. Gribe ham aabenlyst turde de ikke, fange ham
i Ord kunde de ikke. De lurede da paa Leilighed, og
gave imidlertid deres Følelser Luft i forskjellige Domme.
Snart hed det sig med en lærd Fornemheds paatagne Overlegenhed:
han er en Idiot, et uvidende Menneske, der
% --- p. 24 ---
\setcounter{footnote}{0}%
\opage{24}ikke kjender Skrifterne; snart med brutal Bagvaskelse:
han fraadser og drikker og lever med den gemene
Hob; snart med ufordulgt Harme: han haver Diævelen,
det er en udlært Folkeforfører, et djævelsk Menneske.
Andre, der vilde gaae for at være mere billigttænkende,
meente maaske, at hvad han saadan sagde om sit
Riges Oprettelse, sit Davidskongedømme og sin Throne,
ikke var saa ligefrem at forstaae efter Bogstaven. Men
saa maatte man da høilig beklage, at han fyldte Mængden
med Indbildninger og talede til Folket om Noget,
som langt overgik dets Fatteevne.

Hadets Ord og Forblindelsens Udbrud anfægte os
ikke; men det Spørgsmaal, om Christus som Folkelærer
virkelig kunde forstaaes af Folket, ligger os
paa Sinde. Hvis det Modsatte var Tilfældet, vilde ikke
blot Christi Gjerning være forgjæves; men de, som
endnu i vore Dage befrygte Alt af en folkelig Oplysning
og indstændig advare imod at tale til Folket
om høie Ting, vilde jo have fuldkommen Ret. Vi
kunne trøste os over denne Misforstaaelse. Talerne i
Athen og Rom vidste neppe altid, hvad de gjorde, naar
de talede til Folket og meente Mængden; men Folketaleren
i Galilæa er sig klart bevidst, hvad han gjør,
naar han taler til Mængden og mener Folket: „hvo,
som har Øren at høre med, han høre“.

Christus vidste fuldt vel, at der i den hele Befolkning
var ikke blot ædle, men ogsaa uædle Elementer;
han kjendte Hjerternes Haardhed, de Seende,
som dog ikke see, de Hørende, som dog ikke høre.
Han sammenligner sit Riges Ord med en Udsæd og
taler om Sædemanden: Noget falder ved Veien, Noget
paa Steengrund, Noget iblandt Torne, men der er og%
% --- p. 25 ---
\setcounter{footnote}{0}%
\opage{25}saa Noget, som falder i god Jord og bærer, tredivefold,
tredsindstyvefold, hundredefold. Alt Folket hængte ved
ham og hørte ham, og han vidste godt, at det kunde
forstaae ham, naar det oprigtig vilde. Da ikke alene
de tolv, men siden ogsaa de halvfjerdsindstyve Disciple,
som han havde udsendt, vendte tilbage og glade forkyndte
ham, hvorledes Ordet var blevet modtaget,
„frydede han sig i Aanden og sagde: jeg priser dig
Fader, Himmelens og Jordens Herre, at du haver skjult
disse Ting for de Vise og Forstandige, og aabenbaret
de Umyndige dem.“ Her er for Folkelivets Sag et
Tidernes Vendepunkt. Der er Noget, der overgaaer
det menige Folks Fatteevne, det er i Lærdom og
Videnskab: Stjernernes Gang, Jordklodens lange Udviklingshistorie,
de levende Slægters naturlige Afstamning
fra lavere til høiere i opadgaaende Led; her
er i disse Oplysninger noget Herligt og Stort, Noget,
som det folkelige Samfund ikke kan undvære. Og dog
kan det menige Folk ikke forstaae det, et Beviis for,
at det vel er et Høiere, men ikke det Høieste. Altsaa:
det Høiere kan hele Folket ikke forstaae, men
det kan med Sjælens Inderlighed, i Hjertets Eenfoldighed
forstaae det Høieste. Vidunderlige Viisdom! Det er
denne Livstanke, der kom til Verden med Christus;
det er ham, som staaer inde for en ligelig Forstaaelse
af de store Grundsætninger om Lighed og Frihed, en
Forstaaelse, der væsenlig er forligt med Forskjellen
mellem Vise og Eenfoldige, Lærde og Ulærde. Ved
at sætte denne Livstanke i Kraft har Christus grundlagt
den folkelige Oplysning og gjort det første afgjørende
Skridt til Folkenes Befrielse.

Den Folkeoplysning, Christus har grundlagt, er
% --- p. 26 ---
\setcounter{footnote}{0}%
\opage{26}imidlertid blevet misforstaaet fra to mod hinanden
stridende Synspunkter, et eensidig verdsligt og et lige
saa eensidig gudeligt.

Fra verdslig Side har man vel villet indrømme,
at der er Menneskekundskab, Stemning og dyb Følelse
i Alt, hvad Nazaræeren siger. Hans Ord om Lilierne
paa Marken og Fuglene under Himlen ere egnede til
at skaffe det af Næringssorger og Bekymringer nedtrykte
Sind en velgjørende Lindring. Men, indvender
man, det er ikke nok med Stemninger. Vi leve paa
Jorden i jordiske Forhold. Driftige Folk maae passe
deres Næringsvei; de maae arbeide og spinde, de maae
saae og høste og samle i Lade, de kunne ikke som
Lilierne eller som Fuglene være ubekymrede for den
Dag imorgen. Det skulde altsaa --- mene de Verdenskloge
--- være en Mangel ved Christi Underviisning,
at han ingen Hensyn tog til det Praktiske, det Nærmeste,
det Nyttige. Istedenfor at lade sine Tilhørere
betragte Lilierne paa Marken og Fuglene under Himlen
burde han strax have lært dem at læse, skrive,
regne, sat dem ind i Statshuusholdningen, især i Politiken;
hvilket altsammen maa ansees for nyttigt og nødvendigt.
En stor Misforstaaelse. Hvad her anføres som Mangel,
er netop Fuldkommenhed. Det, som Christus meddeler
Folket, indeholder netop Forudsætning og Betingelse
for alt det Praktiske, alt det Nyttige, hvori
der er Velsignelse.

Der er noget Usundt i den feberhede Iver og hensynsløse
Graadighed, hvormed et verdsligt Sind kan give sig
hen til Vinding og Vinding i Landbrug, i Haandværk,
i Handel og Industri, kort sagt, til Alt, hvad der skaffer
Fordeel, Forskjønnelse og Nydelse af det nærværende
% --- p. 27 ---
\setcounter{footnote}{0}%
\opage{27}Liv. Især er intet farligere for et Folk, der skal lære
at bruge sin Frihed, end at kaste sig voldsomt og
hovedkulds ind i et politisk Røre, menende, at Politik
er det Høieste paa Jorden. Er Sindet ikke fyldt og
bevæget af noget Høiere, vil det lettelig skee, at Rethaveri,
Egenraadighed og urene Lidenskaber forvirre
Synet og rive selv de Bedre hen til en sørgelig Efterligning
af de Konstgreb, der plette Stormandspolitiken,
det Uretfærdighedernes Væv, hvori Trældommen sidder.

Christi Ord er eenfoldigt og dybt: „Søger først
Guds Rige og hans Retfærdighed, saa skulle og alle
de Ting, som I have Behov, tillægges Eder.“ Hvorledes
dette Ord bliver at opfatte i Menigheden, skal
her ikke omhandles. Vi tage det i denne Sammenhæng
som et Ord til Folket; vi skulle overbevise os om, at
det finder en stor, afgjørende praktisk Anvendelse.
Langt fra at lamme et Folks Arbeidskræfter eller gjøre
det ligegyldigt mod de timelige Opgaver, er dette Ord tvertimod
det eneste, der giver fuld Besvarelse af et afgjørende
praktisk Spørgsmaal. Hvad er det et Folk trænger til for
at blive det driftigste og vindskibeligste i Verden, blive
et Folk, der har Sands for alt Ophøiet og Skjønt, et
arbeidsomt Folk, der er i Stand til at bære endogsaa
tunge Byrder? Til Guds Rige --- Aandens Rige.

At være i Aandens Rige og føle sig baaren af en
usynlig Magt er med andre Ord at føle sig opløftet
over det Timelige. Men at føle sig opløftet over det
Timelige er en Betingelse for at kunne sænke sig deri
og arbeide dermed uden Fare for at gaae under. Dén,
der veed med sig selv, at han, om end Ulykken bryder
ind, har en større Kapital i Baghaanden, kan vove det
Yderste, og den, der ikke behøver at være angst for
% --- p. 28 ---
\setcounter{footnote}{0}%
\opage{28}Følgerne, kan benytte Øieblikket. Vi tage Livet som
det er; lad os betragte de Nydende. Hvem nyder vel
Timelighedens Goder allerbedst, enten den, som sætter
sit eneste Maal i Nydelsen, eller den, som midt i Nydelsen
er hævet over Nydelsen og har Aand og Kraft
til at bryde af, hvad Øieblik det skal være? Dog vel
den Sidste; thi den Første er en fortvivlet Sjæl, der
jamrer sig, naar Nydelsens Mulighed hører op, og døer
af Kjedsommelighed, hvis den er kjedsommelig nok til
at holde ud med ham. Lad os betragte de Kæmpende.
Hvem er den Tappreste, enten den, hvem Saar og Død
er den største Ulykke, eller den, som veed sig hævet
over Stridens Udfald og kan falde, om saa skal være?
Den Sidste er midt i Faren opløftet over Faren, og
Faren selv er næsten bange for ham; den Første er
en Stakkel, hvem Faren altid er efter og ligesom har
sin Fornøielse af at kyse, fordi han saa inderlig gjerne
vil løbe fra den. Lad os saa tage Sagen i Alvor,
tænke paa Christi Ord og betragte det timelige Livs
Vilkaar. Hindringer og Forviklinger, Modstand og
Skranker allevegne! ikke blot for de Smaa, men ogsaa
for de Store; dog efter Hverdagsreglen især for de
Smaa.

En fattig Mand har mange Byrder at bære, hvordan
skal han holde det ud? Han mander sig op, han
tager sig sammen, han vil i sin Strid med det Timelige
opløfte sig over det Timelige. Dersom der nu i
hans Sind ikke er Andet end timelig Kraft og timelig
Attraa, hvad bliver saa Enden? Den Timelige overvindes
af det Timelige, taber Modet og lader Hænderne
synke. Men dersom der er i hans Sjæl en Evighedsgnist,
en Funke af Aand, da næres han af usynligt
% --- p. 29 ---
\setcounter{footnote}{0}%
\opage{29}Brød, da holder han ud; hans Mod forvandler sig til
Taalmod, og Taalmodets Mod er uovervindeligt.

Et lille Folk maa i de store Verdensforhold fra
mange Sider føle sig som den fattige Mand. Om det
just ikke lider af Næringssorg, ja selv om det bugner
af Velvære: saa er der dog et Tryk, en Bekymring,
som ikke vil ophøre. Overmagten kan komme med
Vold og Rov, hvorledes skal det lille Folk finde Midler
til at værge om sin Ret, forsvare sin Selvstændighed
og holde Modet oppe? Er der nu i et saadant Folk
ingen Aand, ingen høiere Villieskraft, ingen Evighedstanke;
er Alt kun timeligt, timelig Beregning af timelig
Vinding og timeligt Tab, da vil det skee, at den
Timelige overvindes af det Timelige: Anstrengelsen
mattes, Kræfterne slappes, Mismodet hvisker: hvad kan
det hiælpe? % print: „hiælpe“ (for hjælpe)
Men dersom Aanden er med, en Evighedstanke,
en Villiesreisning, som griber Folkets Sind og
løfter det op over timelige Brydninger, da forandres
Tilværelsens Vilkaar. Alt, siger den kolde Forstand,
gaaer efter Love. Rigtigt. Men --- lad os ikke glemme
det --- der er i Verden to Slags Love, Nødvendighedslove
og Frihedslove, Naturlove og Aandslove. Hvor
inderlig de end hænge sammen, saa ere de dog forskjellige.
At ville bygge paa Friheden og dog troe
paa den blinde Nødvendighed, er en Modsigelse, er: at
ville raade selv for sit Liv og dog give sig ind under
Skæbnen. Christi Ord vil bryde Skæbnens Aag og
stille det kraftigvillende Folk under Forsynet. Det
styrende Forsyn, som lader sin Sol opgaae over Onde
og Gode, forvandler Mod til Taalmod; det faderlige
Forsyn, uden hvis Villie ikke en Spurv falder til Jorden,
% --- p. 30 ---
\setcounter{footnote}{0}%
\opage{30}forvandler Taalmod til Mod: „frygter ikke, I ere bedre
end mange Spurve.“

Den lysskye Gudelighed seer Sagen fra en anden
Side; den vil slet ikke indrømme, at Aanden i Christi
Ord har Noget med Folkelivet at gjøre. Den forstaaer
Formaningen: „Søger først Guds Rige“, som om det
Høieste var naaet, naar man blot kastede det Jordiske
bort og ikke befattede sig med Andet end med at
græde og hyle over Verdens Synd og Daarskab. Vi
tør ikke glemme det hellige Ord: „Salige de, som
sørge“; men den Sorg, hvori der er Salighed, kan umulig
bestaae i at slukke alle Glædens Lys for blot at
see suurt til Verden. Dersom det i sig selv var en
Synd at bruge det Timelige og arbeide med det Verdslige,
da var ikke blot Folkelivet i særlig Forstand,
men Menneskelivet, fra hvilken Side man end betragter
det, med eet Slag gjort umuligt. Den sygelige Fromhed
har aldrig formaaet at gjennemføre sin Livsbetragtning,
har aldrig drevet det til at opgive alt Timeligt.
Den er blevet staaende ved en Halvhed uden at
giøre sig Rede for det Aandeliges Krav paa det Naturlige;
den har villet være med og dog ikke være med i
det fremadskridende Samfund, den har af lutter Ærbarhed
ikke vidst, om den turde leve blandt virkelige
Mennesker. Middelalderens Eremiter og Munke gik
frem med en ganske anden Følgestrenghed; men Historien
lærer os, hvorledes det endte.

Et af To: enten har Christus villet hjælpe ogsaa
Folkelivet frem eller han har ubetinget villet standse
det; han var --- vi tale menneskeligt om ham --- ikke
den Mand, der blev staaende paa Halvveien. Havde
det Sidste ligget i hans Plan, da havde han visselig
% --- p. 31 ---
\setcounter{footnote}{0}%
\opage{31}ikke afviist den fastende Hellighed; han havde brudt
Staven over hele den retlige Samfundsorden, og fremfor
Alt: han havde begyndt med at forbyde Ægteskabet
og banlyse Slægtens Forplantelse. Naar han
nu gjør lige det Modsatte. Naar han nu betragter den
menneskelige Orden, de givne Forhold mellem Fyrster
og Folk, Øvrighed og Undersaatter, Forældre og Børn,
Herrer og Tjenere som Noget, der naturlig følger af
sig selv, ja som en guddommelig Orden; naar han
atter og atter anvender disse Forhold som Underlag
for sine Lignelser om Himmeriges Rige; naar han
rolig og fri gaaer ud og ind i Menneskenes Huse, seer
paa deres Færden og Virken med milde Øine, deeltager
i deres Gæstebud, hjælper de Lidende, trøster de Bedrøvede,
velsigner de smaa Børn og lader Glædens
Viin iskænkes for Brud og Brudgom: hvo tør da sige,
at Christus har brudt Staven over det Timelige, det
Menneskelige, det Folkelige? At velsigne Familielivet
er at velsigne Folkelivet, thi hvad er vel Folkelivet
andet end et vidt forgrenet, mangfoldige Familier omfattende
Familieliv?

Hvad Englene forkyndte Hyrderne, „en stor Glæde,
som skulde vederfares alt Folket“, det har Christus
som Folkes Lærer og Fører i høieste Maal villet gjøre
til Sandhed. At han, der var Sjælenes Frelser, ogsaa
vilde være sit Folks Frelser, har han selv udtrykkelig
viist. Det laae ham tungt paa Sinde, hvorledes det
skulde gaae med det lille Folk, hvis store Time nærmede
sig. Det var et kraftigt, men gjenstridigt Folk,
der tærede paa store Minder. Han saae det adsplittet
i Partier, Pharisæere, Sadducæere, Herodianere, Romersksindede,
Saadanne, som ivrede for det Gamle, og
% --- p. 32 ---
\setcounter{footnote}{0}%
\opage{32}holdt paa det Nationale, Saadanne, som higede efter
Nyt og beundrede det Fremmede; han saae dets Lærere
uden Aand, dets Øvrighed uden Myndighed, den menige
Hob som Faar, adspredte paa Bjergene uden
Hyrde; han saae den mægtige, lurende Fjende, Romerørnen
med de skarpe Kløer, der snart skulde gribe
sit Bytte; han saae det Alt, og han sørgede.

Ved hans Indtog i Jerusalem, da Disciplene jublede
og Mængden raabte Hosianna, var Folkefesten
ingen Glædesfest for ham. „Da han kom nær til og
saae Staden, græd han over den og sagde: Vidste du
dog, ja end paa denne din Dag, hvad der tjener til
din Fred. Men nu er det skjult for dine Øine. Derfor
skulle de Dage komme over dig, at dine Fjender
skulle kaste en Vold op omkring dig, og beleire dig
rundt om og trænge dig allevegne. Og de skulle lægge
dig slet øde og dine Børn i dig, og ikke lade Steen
paa Steen i dig, fordi du ikke kjendte din Besøgelses
Tid.“

Han \emph{vilde}, men de \emph{vilde ikke}, det var hans
Klage. Selv da han gik til Golgatha og bar sit Kors,
og Kvinderne græd over ham, forglemte han sine egne
Lidelser ved Tanken om, hvad der forestod det ulykkelige
Folk. „I Jerusalems, Døttre! græder ikke over
mig, men græder over Eder selv og over Eders Børn.“

Der er Knudepunkter i Folkenes Historie, Vendepunkter,
om hvilke det gjælder: nu eller aldrig. Forsømmes
den beleilige Tid, misbruges de gode Vilkaar,
gribes der feil, gjøres der Overgreb: da maa Slægter
bøde derfor igjennem Aarhundreder. Intet Folk paa
Jorden har kæmpet med et saadant Løvemod, en saa
storartet Fortvivlelsens Kraft for sin Ret og Selv%
% --- p. 33 ---
\setcounter{footnote}{0}%
\opage{33}stændighed som det jødiske Folk paa det Sidste. For
silde, Hjælperen var borte: vee det Folk, som ikke
kjender sin Besøgelses Tid. Mere end halvandet tusinde
Aar, fulde af Lidelser, skuffede Forhaabninger og bittre
Erfaringer, ligge mellem hiin Tid og os. Nu kommer
Besøgelsens Tid igjen; thi Folkene vaagne.

Folketaleren fra Galilæa vil atter frem. Troer
ham, han glemmer ikke sit Alter og svigter ikke
sin Menighed, fordi han endnu en Gang viser sig, som
han menneskelig er, paa Verdenshistoriens Strøm; han
har Kræfter tilovers. Længe nok have Folkenes selvbeskikkede
Tugtemestere, geistlige som verdslige, holdt
ham indespærret som i en lummer Luft indenfor officielle
Kirkemure. Han vil ud i det Frie, under aaben
Himmel; fra Bjerget, blandt Lilierne paa Marken og
Fuglene derude vil han sees af Mængden og høres af
Folket.

Vi kunne ikke af os selv begynde det store Værk;
vor Styrke er for ringe, vort Syn for mat, vort Hjerte
for svagt. Men kommer Han selv, den Vældige i Aandens
Kraft, da vil hans Røst kunne høres af Alle. Han
veed, hvad der boer allerinderst i Mennesket, han veed,
hvad Opoffrelse er, og hvad det er, der kaldes Mod.
Han har Ordet, han er jo Ordet selv, det Sandhedens
Ord, som gaaer til Hjertet, det Retfærdighedens Ord,
som rammer Samvittighederne, det Broderlighedens
Ord, som hævder Frihed og Lighed, det Kjærlighedens
og Forsonlighedens Ord, som jævner al Strid mellem
Smaa og Store.

„Ville I nu?“ det er hans Spørgsmaal til alle
Folkefærd; ville I? --- Han sætter sin Tornekrone i
Pant, at Værket skal faae Fremgang. De kunne høre
% --- p. 34 ---
\setcounter{footnote}{0}%
\opage{34}ham, hver i sit Tungemaal; de kunne alle forstaae
ham --- thi Folket forstaaer det Høieste; de kunne
alle følge ham, hver paa sin Vei, i hvilket Land, de
end boe, og hvilke Farver de end bære. Det Fælles,
han byder, det første Fornødne, er netop skikket til at
fortære det Fjendtlige og nære det Eiendommelige.

Vilde dog vort Folk, vort lille Folk med de store
Minder under det tunge Tryk, modtage ham med aabne
Arme! Vi bære jo hans Bannermærke; den røde Dug
med det hvide Kors, det er hans Seiersfane. Fremad,
frem, hvor Fanen veier; fremad, frem, for Sandhed og
Ret; fremad, frem, hvor han selv fører an, Folkenes
Konge, Folkenes Lærer og Leder.

% --- p. 35 ---
\setcounter{footnote}{0}%
\lecture{Tredie Foredrag,} {holdt i Aagaard.}
\opage{35}For at kjende Menneskelivet, dets Anlæg og Maal og
Veien, som fører til Maalet, maae vi, saa vidt vi formaae,
udforske Tilværelsens Love. Vi maae gjøre
grundigt Bekjendtskab med den virkelige Verden, vi
maae raadføre os med Natur og Historie. Det kommer
imidlertid ikke alene an paa, hvad det er, vi see, eller hvormeget
vi faae at see; det kommer ogsaa, og det fornemmelig,
an paa de Øine, hvormed vi see, det Lys,
hvori vi see, det Stade, hvorpaa vi stille os.

Her er det nu strax indlysende, at Mennesket ikke
blot lever i een, men i flere Verdener. Den udvortes
Natur med dens mange Skabninger, Dannelser og Frembringelser
ligger udbredt for vort Blik; dens Sommer
og Høst, dens Vinter og Vaar taler til os i et eget
Sprog og kalder os til forskjelligt Arbeide. Men vi
have jo ogsaa en indre Natur, en Verden af Følelser
og Tanker, vi bære i os selv, og denne indvortes Verden,
hvor kan den dog være forskjellig. Vi gjøre da
ogsaa den Erfaring, at det fornemmelig er den indvortes
Verden, der spreder sit Lysskær over den udvortes.
Naar vi ere glade, saa smiler Alt i Mark og
% --- p. 36 ---
\setcounter{footnote}{0}%
\opage{36}Skov; men ere vi bedrøvede, saa kan Solen skinne saa
meget den vil, det er som bar den hele Natur paa en
stille Sorg. Det indre Øie kan være meget forskjelligt,
og Forklaringen kommer fra det Indre; den Ene seer
i Naturen et forunderligt Spil af Stoffer og Kræfter,
et Nødvendighedsvæv af blinde Love; den Anden seer
Guds Viisdoms Spor, en Aabenbarelse af Skaberens
Almagt.

Noget Lignende gjentager sig, naar vi betragte
Historien. Menneskeslægtens Levnetsløb skrider frem
uafhængigt af os, thi Enhver af os er kun et forsvindende
Led i Slægten. Men hvert enkelt Menneske har
dog sit eget Levnetsløb, sin egen lille Historie; dersom
de mange smaa Historier skulde fortælles udførlig, oprigtig,
hvor vilde de, med alt det Fælles, være forskjellige.
Herved er da ogsaa det Mærkelige, at Verdenshistorien
og vor egen Historie, den store og den
lille Historie, hvad Forstaaelsen angaaer, væsenlig svare
til hinanden og vise sig for os i en fælles Speiling af
et fælles Lys. Det er ikke Begivenhederne, de udvortes,
endelige Tildragelser, hvorom vi her tale; det er Tydningen,
Dommen, Opfattelsen af det Hele, hvorpaa det
kommer an. Den, der ikke kan finde Guds Forsyn i
sit eget Liv, hvorledes skulde han vel opdage det i
Verdenshistorien.

Vi have mange Historielærere, og vi ere dem taknemmelige;
vi have lærde Forskere, som vide at skjelne
mellem Sandt og Falsk og udbrede Lys over Begivenhedernes
Sammenhæng; vi have Tænkere, som arbeide
paa at udfinde Historiens Love: vi ere dem meget taknemmelige.
Skulde det engang lykkes at opdage den
Grundlov, hvorefter Verdenshistorien bevæger sig, lige
% --- p. 37 ---
\setcounter{footnote}{0}%
\opage{37}som man jo allerede har opdaget den Grundlov, hvorefter
Kloderne bevæge sig i Himmelrummet: da maatte
vi vel tilsidst komme saa vidt, at vi kunde udregne,
hvad der vil skee i Fremtiden. Mon det nogensinde
vil lykkes? Hidtil lyde Dommene om Fremtiden, vi
mene ikke blot om vor egen nærmeste, men om den
historiske Fremtid i det Hele, høist afvigende. Der er
dem, der vente sig en stor Tid, en herlig og lyksalig
Tid, skjøndt de rigtignok ere indbyrdes uenige om,
hvori Herligheden og Lyksaligheden skal bestaae. Andre
see Fremtiden an med mørk Bekymring.: % print: „Bekymring.:“
Hvirvel,
Opløsning overalt, Brud paa gamle Vedtægter, Hensmuldren
af de gamle Samfundsstøtter, hvor ængsteligt,
hvor sørgeligt!

Vanskeligheden med at opdage Historiens Love
ligger deri, at det ikke blot er Nødvendighedslove, som
i den materielle Natur, men ogsaa Villieslove, Frihedslove.
De, som slet ikke kunne faae Øie paa andre
Love end de saakaldte uforanderlige Naturlove, maae
komme til hvilken Forklaring, de ville, deres Forklaring
er altid eensidig og falsk; og det af den eenfoldige
Grund, at Historiens Love ere mere end blotte Naturlove:
de ere Aandslove.

Skulle vi tænke paa en Historielærer, som kan
holde Speilet op og vise os Træk af Fremtidens Ansigt,
da maa det være ham, som i Aandens og Guddommens
Fylde selv er Historiens Vendepunkt; det kan ikke
være nogen Anden end Christus. Men hvad er det,
Christus lærer os om Historien, særlig om den historiske
Fremtid? Synerne forfærde os, vi gribes af Rædsel.
Først Jerusalems Ødelæggelse og saa strax derefter ---
Verdens Undergang; „Strax efter de Dages Trængsel
% --- p. 38 ---
\setcounter{footnote}{0}%
\opage{38}skal Solen formørkes, og Maanen ikke give sit Skin,
og Stjernerne falde af Himlen og Himlens Kræfter
røres. Og da skal Menneskens Søns Tegn aabenbares
i Himmelen; og da skulle alle Jordens Slægter hyle“.
(Math. 24, 29—30).

De, som bygge deres historiske Betragtning paa
rene Naturlove, have let nok ved at befrie sig fra
Forfærdelsen. Naar de hyle, saa hyle de kun til Ære
for „de uforanderlige Love“, og nedlægge desaarsag de
allerbestemteste Indsigelser mod den hele Spaadom.
Deres Beviis er kort og klart: „Det forudsiges, at
Verden skal forgaae \emph{strax efter} Jerusalems Ødelæggelse.
Verden har siden den Tid ikke desto
mindre nu staaet i atten hundrede og fem Aar;
altsaa er Spaadommen gjendreven. Hertil kommer,
at den hele Betragtning af Himmel og Jord er forældet.
Jorden er ikke det Midtpunkt, hvorom Alting
dreier sig, Stjernerne sidde ikke saa løst paa Himmelhvælvingen,
at de skulde falde ned“, o. s. fremdl.
Saalænge man udelukkende følger Naturlærens Anviisning,
ville disse Indvendinger være uigjendrivelige.
Vee dem, som have paataget sig det overmenneskelige
Hverv at forsvare Aabenbaringens Ord om Tegnene i
Sol og Maane med videnskabelige Grunde: deres Fremtid
er mørk, de ere ilde stedte.

Anderledes, naar vi stille os paa Personlighedens
Standpunkt og spørge om Aandslovene. Jorden er da
Midtpunkt, forsaavidt Mennesket, det Væsen, hvori
Aanden boer, er Midtpunkt. Christus, „Verdens Lys“,
er Historiens Midtpunkt. Hvad han har udrettet paa
Jorden, vilde ikke blive forhøiet ved at være udført
paa en større Stjerne, var det end det fælles Tyngde%
% --- p. 39 ---
\setcounter{footnote}{0}%
\opage{39}punkt, hvorom Solsystemerne dreie sig. Enhver, som
har Aand, vil forstaae, hvor falsk det er at anlægge
en Maalestok for udvortes Storhed i Rummet, naar
Talen er om aandelig Storhed. For Christi prophetiske
Syn svinder den hele historiske Fremtid ind til nogle
faa Dage; thi han maaler med Evighedens Maal. Hvad
er vel en Dag imod hundrede tusinde Aar, og hvad er
selv hundrede tusinde Aar mod Evigheden? Derfor
siger Christus: „Strax efter de Dage“. Men hvorledes,
spørger man, kan Solen formørkes, og Stjernerne falde
af Himmelen? --- Det maa vel skee, naar Himlens
Kræfter røres, og Menneskens Søns Tegn aabenbares
i Himlen; det maa vel skee, saa sandt det er Aanden,
som bærer Naturen.

Naturlæren er god og dens Veiledning fortræffelig,
naar vi tage den for, hvad den er, en Opklaring af det
Virkelige, en Belysning af de naturlige Aarsager, en
Paaviisning af de givne Love. Men naar den anvendes
til at bevise, at „de givne Love“ ikke \emph{ere givne}, fordi
der ingen Lovgivning er og ingen Lovgiver: da er
dette et videnskabeligt Overgreb og en Krænkelse af
Aand, thi Aanden er det Oprindelige. For Aanden
selv, den almægtige Eenhed af Viisdom og Villie, er
den timelige Natur under alle Skikkelser i al dens
Varighed med al dens Storhed en forsvindende Størrelse.
Aandens Magt i Materien er evig den samme, enten
Massen kun er et Fnug eller tusinde Gange større end
Solen, thi Aandens Magt er Almagt.

Christus, Aandens, den Almægtiges Aabenbarer,
veed med sig selv, at det inderste Bevægende i ham
er det, som bevæger Himmel og Jord. Fra dette
Høidepunkt maa hans prophetiske Fremtidssyn opfattes.
% --- p. 40 ---
\setcounter{footnote}{0}%
\opage{40}Den uklare Blanding af Sandt og Falsk, af Godt og
Ondt, som forvirrer os, saalænge Endeligheden med de
mange Betingelser, staaer i Forgrunden, forsvinder, naar
Alt føres tilbage til det første Udspring. Historien
viser sig her i overhistorisk Belysning. Ikke nok, at
Folk staaer op imod Folk og Rige imod Rige, at der
er Hunger og Pestilentse og Jordskjælv her og der;
det Inderste i „Veerne“ er, at Aand staaer op imod
Aand, at de skærende Modsætninger blotte sig og støde
sammen uden mæglende Mellemled: Sandhedens og
Løgnens, Lysets og Mørkets, Livets og Dødens Magter
færdes her i den sidste Strid.

Hvor haard end Striden er, saa er dens Udfald
dog ikke tvivlsomt. Seieren er i Grunden allerede
vundet. Christus seer sig i Synet som Seierherre,
skjøndt hans Lidelse og Død endnu forestaaer. „Jeg
saae,“ siger han „paa et andet Sted, Satan falde ned
af Himlen, som et Lyn“; det Ondes Magt er i evig Forstand
en Afmagt, hvormegen Magt det end kan have i
det Timelige. Det er den samme Tanke, som i Johannes'
Aabenbaring (12, 7—12) er fremstillet under følgende
Billede: „Og der blev en Strid i Himlen; Michael og
hans Engle strede mod Dragen, og Dragen stred, og
dens Engle. Men de mægtede Intet; ei heller blev
deres Sted ydermere fundet i Himlen\dots{}. Derfor
fryder Eder, I Himle, og I som boe i dem; Vee dem,
som beboe Jorden og Havet; thi Djævelen er nedsteget
til Eder! Han haver stor Vrede, fordi han veed, at
han kun haver liden Tid.“ Den historiske Kamp mellem
det Gode og det Onde bliver altsaa at opfatte som de
sidste Optrin af et stort Slag, hvori Fjenden har tabt,
men strider endnu Fortvivlelsens Strid. Den sig fort%
% --- p. 41 ---
\setcounter{footnote}{0}%
\opage{41}sættende Historie er den lange Slagdag, der strækker
sig ud i Tiden, at Befrielsen maa fuldendes i Virkeligheden,
at saa mange, som ville, maae deeltage i Kampen
og blive deelagtige i Seieren.

Men --- her reiser sig en ny Indvending --- naar
den hele historiske Fremtid kun skal betegne Verdens
fortsatte Undergang, hvad bliver der saa af vor Forventning
om Menneskerettighedernes Tid, den nye
Folkets og Folkenes Tid, hvoraf vi love os saa meget?
Skal Christus være vor Historielærer, da er det nok
bedst at lade den Slags Indbildinger fare og slutte os
til dem, som klarlig see, at Alt bliver værre og værre.
Kan det Gode kun seire ved, at Verden gaaer under,
saa gjøre vi jo bedst i at arbeide med paa dens Undergang.


Hvad betyder saa al denne Talen om det Folkelige,
det Nye, det Store, det hidtil Useete; med mindre
man da vil forstaae Sagen saaledes, at Folkevækkelse,
Folkeliv og Folke-Aander skulle være de kraftigste
Midler til at fremskynde Verdens Ødelæggelse? --- Den
kolde Spot forraader en dulgt Fortvivlelse. Paa Fremtidens
Vei ligge mange Stationer. „Lærer en Lignelse
af Figentræet: naar Vædsken er kommet i dets Grene,
og Bladene skyde ud, da vide I, at Sommeren er nær.“
Denne Lignelse passer ikke alene paa det allersidste
Gjennembrud, den passer ogsaa paa de mange mellemliggende.
Mange Somre kunne komme, hver med en
Tidens Høst, før vi naae til Evighedens Sommer.

Udviklingsloven er nu denne: En Tidsforvirring
opstaaer, hver Gang et Bestaaende løser sig op, og
noget Nyt skal frem. Det gamle falder hen; det Nye,
som er i Færd med at danne sig, har endnu ikke
% --- p. 42 ---
\setcounter{footnote}{0}%
\opage{42}vundet Skikkelse: af Gjæringer, Forstyrrelser og vild
Uorden udklarer sig saa, vistnok under voldsomme
Anstrengelser, efterhaanden en ny og bedre Tingenes
Orden. Samfundet kommer i Ligevægt, fordeler sine
Arbeider, saaer og planter og bereder sig paa en rig
Høst. Men jo rigeligere man nyder alle timelige Gaver,
jo almindeligere man drømmer om evig Fred og ubegrændset
Velvære, desto føleligere bliver Overraskelsen,
naar den Omvæltning, som en snigende Fordærvelse
længe har forberedt, pludselig indtræder. Der maa da
af nye Gjæringer og nye Kampe danne sig en ny og
bedre Tilstand, altid paa andre og høiere Betingelser;
men enhver, selv den fuldkomneste, Tilstand paa Jorden
er forbigaaende. Hvor inderlig denne Tanke hænger
sammen med den christelige Hovedtanke, kan udtrykkelig
paavises.

I Johannes' Aabenbaring (20, 1 --- 3) møde vi følgende
Syn. „Og jeg saae en Engel fare ned fra Himmelen;
han havde Afgrundens Nøgle og en stor Lænke
i sin Haand. Og han greb Dragen, den gamle Slange,
som er Djævelen og Satanas, og bandt ham for tusinde
Aar, og kastede ham i Afgrunden, og tillukte og forseglede
over ham, at han ikke mere skulde forføre
Folkene, indtil de tusinde Aar fuldendtes; og efter dem
bør det ham at løses en liden Tid.“ Det er det tusindaarige
Rige, hvis Komme bliver bebudet; men hvad
betyder det tusindaarige Rige? Vi tør ikke sammenblande
det Overhistoriske med det Historiske og hengive
os til Sværmerier. Hvorledes Forventningen om
et tusindaarigt Rige skal opfattes i Menigheden, er en
Sag for sig: her staae vi paa folkelig Grund og have
slet og ret med en Aandslov at gjøre. Hver Gang en
% --- p. 43 ---
\setcounter{footnote}{0}%
\opage{43}ny og bedre Tingenes Orden seirer, bliver Mørkets Aand
bogstavelig kuet; den onde Villie lænkebindes, tvinges
ned i Afgrunden, og maa holde sig skjult til en Tid,
da den igjen kommer løs: det er den almene Tanke.
At selv det tusindaarige Rige, hvor herligt og mægtigt
det end skal være paa Grund af „den første Opstandelse“,
ikke desto mindre er beluret af stærke Fjender og udsat
for store Farer, skildres med levende Farver (Joh.
Aab. 20, 7—8). „Naar de tusinde Aar ere fuldendte,
skal Satanas løses af sit Fængsel. Og han skal udgaae
at forføre Folkene fra Jordens fire Hjørner, Gog og
Magog, at forsamle dem til Strid; deres Tal er som
Havets Sand.“

Det er en vigtig Lære, vi heraf kunne uddrage.
Vi tør ikke forgude det Folkelige eller sætte et Folkerige
istedenfor det Himmeriges Rige, i hvilket Christi
Aabenbarelse fuldendes. Ikke desto mindre kan Tusindaarsriget
dog være os et Billede paa, at Christi Aand
vil give sig Afpræg ogsaa i jordiske Tilstande: Sjælenes
Frelser er ogsaa Folkenes Frelser. Man kan nu sætte
det Folkelige saa langt under det Tusindaarige, som
man vil, saa maa dog dette staae fast: skal Sandhed
og Ret vinde Overhaand over Løgn og Uret, da maa
det skee i Kraft af Retfærdighedens og Sandhedens
Aand. Er Timeligheden for ringe og Endeligheden for
ussel til at afgive Stof for den opad stræbende Aand,
hvad har saa Fremtiden at byde? Villiesslaphed og
Sindstomhed, svækkede Nerver og udbrændte Hjerner,
raae Masser og ubrugte Kræfter, Overmættelse og dog
Hunger, Planer og dog Raadvildhed, Væmmelse ved
Alt og dog graadig Higen efter Nydelse, alt dette i et
Mellemhverandre, rørt sammen til de besynderligste
% --- p. 44 ---
\setcounter{footnote}{0}%
\opage{44}Blandinger og Bundfældinger, skal da ikke give Menneskelivet
et høiere Værd, skal ingenlunde svække den
eenfoldige Tro, at det er Christi Aand, der maa komme
med Foraarstegn, at Sommeren kun da er nær, naar
Figentræet faaer Safter i Grene og Blade.

Den historiske Drivkraft kan maales med Spændingen
mellem det Evige og det Timelige. Der voves
meget, naar det gjælder Livet, allermeest dog, naar der
er en stærk Følelse af, at det gjælder det evige Liv.
Det Aandelige kan ikke gjøre Indtryk paa Mennesker,
allermindst paa et heelt Folk, uden at der paa en
eller anden Maade er noget Sandseligt med. Christi
Formaninger ere indklædte i Lignelser; det er lutter
spændende Optrin, afgjørende Øieblikke, der gløde i
Aand med sandselige Farver. „Vaager derfor, thi I
vide ikke, i hvilken Time Eders Herre kommer. Men
dette skulle I vide, at dersom Husbonden vidste, i
hvilken Nattevagt Tyven vilde komme, da vaagede han
og lod ikke bryde ind i sit Huus. Derfor værer I rede,
thi Menneskens Søn kommer paa den Time, som I
ikke mene.“ Det er ingen ørkesløs Venten; alle Hænder
faae travlt: her er Noget at gjøre. Det er ligesom,
naar „en Mand, der drog udenlands, kaldte sine Tjenere,
og overantvortede dem sit Gods; den Ene gav han fem
Talenter, den Anden to, og en Anden eet, hver efter
hans Evne. Da gik den, som havde annammet fem
Talenter, hen og kjøbslog med dem, og vandt andre
fem Talenter. Ligesaa og den, som havde annammet
de to Talenter, ogsaa han vandt andre to. Men den,
som havde annammet det ene, gik bort og grov i
Jorden, og skjulte sin Herres Penge.“ Hvor haard og
streng er ikke Dommen over denne sidste Tjener. Han
% --- p. 45 ---
\setcounter{footnote}{0}%
\opage{45}træder dristig frem, som om han havde en retfærdig
Sag: „Herre, jeg kjendte dig, at du er en haard Mand,
som høster der, hvor du ikke saaede, og samler der,
hvor du ikke spredte. Frygtende herfor gik jeg hen
og skjulte dit Talent i Jorden; see, der haver du, hvad
dit er.“ Ak, det bekommer ham ilde.

Det Timelige maa bruges, thi Aandsgaverne maae
bruges. Det Aandelige er ikke til for det Timeliges
Skyld; det er omvendt det Timelige, der er til for det
Aandeliges Skyld. Kjendemærket paa, at det Timelige
bruges paa den rette Maade, er, at Aandskræfterne
voxe ved Brugen. Der er i alle timelige Livsforhold,
de være private eller offenlige, smaa eller store, Betingelser
for en Udvikling af Aandslivet: haardføre
Mænd med veirbidte Ansigter og barkede Næver kunne
fuldt saa vel have Aand, som de, der bevæge sig i den
fine Dannelses parfumerede Atmosphære. Særlig maa
det fremhæves, hvad Lignelsen indskærper, at Ingen
tør undskylde sig med, at hans Begavelse kun er ringe.
Den Tjener, som fik det ene Talent, fik samme Forpligtelse
til at arbeide med det ene, som den, der fik
de fem, til at arbeide med de fem. Hvad der gjælder
om det enkelte Menneske, gjælder om et heelt Folk:
at begavet Folk har en stor Forpligtelse. Ikke at
bruge Aandens Gaver er at synde imod Aanden; men
Synd imod Aanden straffes dobbelt. Den første Straf
er, at Gaverne tages bort, at den Paagjældende bliver
aandløs. „Tager derfor det Talent fra ham, og giver
det til den, som haver ti Talenter. Thi hver, som
haver, ham skal gives, og han skal have til Overflod,
men hvo som ikke haver, ham skal endog fratages det,
han haver.“ Den næste Straf er, at den Aandløse
% --- p. 46 ---
\setcounter{footnote}{0}%
\opage{46}bliver dømt som Aandsforbryder. „Kaster den udygtige
Tjener ud i det yderste Mørke; der skal være Graad
og Tænders Gnidsel.“

At Forventningen af Verdens Undergang kan i
Historien falde sammen med storartede Omdannelser
af Verdensforholdene og tilsvarende Gjennembrud af
Aandslivet, beroer just paa, at den er sammenknyttet
med Forventningen af „Menneskens Søns“ Gjenkomst.
Den aandelige Lov, som heri er udtalt, viser nemlig,
at ligesom det Timelige kun er til for det Eviges Skyld,
den ydre Natur kun for Aandens Skyld, Verden kun
for Menneskesønnens Skyld, saaledes er det ogsaa
baade her og hisset Menneskelivet og dets Fuldendelse,
hvorom det Hele dreier sig. Hvorledes det Overhistoriske
under de fremadskridende Bevægelser griber
ind i det Historiske, kan sees ved at kaste et Blik
paa de alt tilbagelagte Stationer. Tre Tidspunkter:
Aposteltiden, Korstogenes Tid og Reformationstiden
frembyde sig for Betragtningen.

I Aposteltiden stod Modsætningen mellem Aandsriget
og de aandsfornegtende, men dog saa mægtige
Verdensriger paa sit Høieste. Nogle Fiskere fra Galilæa,
værgeløse Mænd og eenfoldige Kvinder, det var den
hele Udrustning; det var de Vaabendygtige, hvormed
den nye Aand skulde optage Kampen og fuldende hans
Værk, der var død som en Misdæder paa Korset. Mørke
--- og dog lyse Udsigter! Den store Bortgangne havde
hvisket dem Noget i Kammeret; det skulde de prædike
paa Tagene. Han var ikke forblevet i Graven og
havde glemt de Forladte; han havde besøgt dem igjen,
han havde hilset dem som Dødens Overvinder, og lovet
dem, at han snart skulde komme --- veludrustet. Han
% --- p. 47 ---
\setcounter{footnote}{0}%
\opage{47}havde ikke skuffet dem med et forfængeligt Haab, som
var det strax en Seiersfest. Han kjendte Verden, de
Mægtiges Had og Mængdens Blindhed. „De skulle
overantvorde Eder til Trængsel og slaae Eder ihjel;
og I skulle hades af alle Folkene for mit Navns Skyld.
Og da skulle Mange forarges, og forraade hverandre
og hade hverandre. Og mange falske Propheter skulle
opstaae og forføre Mange. Og efterdi Uretfærdighed
vil være blevet mangfoldig, skal Kjærligheden blive
kold i Mange. Men hvo som bliver bestandig indtil
Enden, han skal blive salig.“ Den første Forudsigelse
har fundet rigelig, bogstavelig Opfyldelse, et Pant paa,
at ogsaa den sidste bestandig bliver opfyldt.

„Maran Atha, Herren kommer, han staaer for
Døren og banker!“ Disse Varsler, som idelig hørtes,
virkede elektrisk paa det synkende Mod og vakte de. % print: „de.“ (a stop before „Stridende“)
Stridende, hver Gang de begyndte at falde i Søvn. Det
var besluttet i Himlen, at „Dyret med de syv Hoveder“,
Verdensstaden paa de syv Høie, Romervældens rige
Stad, „Babylon, den store Skjøge, et Fængsel for alle
urene Aander og for alle urene og afskyede Fugle,“
skulde falde. Skildringen af dens Fald er anskuelig
og gribende. „En vældig Engel fra Himlen opløftede
en Steen som en stor Møllesteen og kastede den i
Havet og sagde: saa skal Babylon, den store Stad, nedkastes
med Hast, og ikke findes mere. Og Harpelegeres
og Sangeres og Fløiters og Basuners Lyd skal
ikke høres i dig mere; og ikke nogen Konstner i nogen
Konst findes i dig mere; og Møllens Lyd ikke høres i
dig mere. Og Lampens Lys skal ikke skinne i dig
mere, og Brudgoms og Bruds Røst ikke høres i dig
% --- p. 48 ---
\setcounter{footnote}{0}%
\opage{48}mere; fordi dine Kjøbmænd vare Fyrster paa Jorden,
fordi ved din Trolddom ere alle Folkene bedaarede.“

Historien har her leveret sit meest storartede Beviis
paa, hvad Sandhedsaanden med aandelige Vaaben
formaaer i timelig Strid, naar Himlens Kræfter røres,
og „Menneskens Søns Tegn“ skal aabenbares i Himlen.“

I Middelalderen var Sammenblandingen af det
Himmelske og det Jordiske saa eiendommelig, og Adskillelsen
af Gudeligt og Verdsligt derhos tillige saa
udvortes, at Aand og Aandløshed næsten flød over i
hinanden. Den Uro og Retløshed, som indtraadte med
Adsplittelsen af Karl den Stores Rige; Muhamedanernes
Angreb, Ungarernes Indfald og især Normannernes
Plyndringer, alt dette i Forening med høie Forestillinger
om Klosterhellighed og munkeagtig Afholdenhed paatrykte
den hele jordiske Tilværelse et mørkt Præg af
Forkrænkelighed. Man meente at være i „de sidste
Veer“. Pest, Dyrtid, Formørkelser og Kometer bleve
anseete som Tegn for „Dragens Løsladelse“ og „Antichristens
Ankomst“. Alle fromme Sjæle beredte sig
paa Verdens Undergang, der blev ansat til Aaret 1000.
Da Aaret gik, og Verden blev staaende, rykkede man Tidspunktet
frem til 1030, og Tegnene bleve endnu mangfoldigere.
En Munk i Ungarn saae en Drage, livagtig
skabt som Leviathan, farende frem med Lyn og Torden;
en fransk Biskop lod bekjendtgjøre, at han havde
faaet et Brev fra den tredie Himmel, hvori det blandt
Andet var befalet, „ikke at nyde uden Vand og Brød
hver Fredag;“ o. s. fremdl.

Men saa vildt og broget det end seer ud i hine
halvbarbariske Tilstande, er Aanden dog med i det
mægtige Tryk; det Evige tynger paa Samvittighederne
% --- p. 49 ---
\setcounter{footnote}{0}%
\opage{49}og varsler om en stor Omvæltning og en Dommedag.
Det tunge Tryk betinger et tilsvarende Opsving. De
gaadefulde Tegn fik omsider en historisk Udtydning.
I Pavens Ord om den hellige Grav blev Maalet klaret.
Opstandelsens Grav! nu var Løsenet givet, og en brusende
Flod af løsladte Kræfter fandt Afløb. Korstogenes
Hovedtanke er en folkelig Troestanke; den
rummer Alt, hvad Middelalderen eier af Eiendommeligt
og Stort. Naar de mange frastødende Træk af blind
Fanatisme, vild Grusomhed, snedig Egennytte og raa
Sandselighed have vakt vor Uvillie, saa tiltrækkes vi
igjen af den Fromhedens Inderlighed, den mægtige Begeistring,
den bjergeflyttende Tro, som, trods Overtroens
Tilsætninger, dog er øst af Golgathas Væld og smager
af Kilden. Menneskesønnens Tegn bliver plettet af Forbrydelser,
men Menneskeslægten gjør Kæmpeskridt fremad.

Reformationstiden betegner i det Hele Middelalderens
Opløsning, og dog paa en Maade dens Fortsættelse.
Vel havde Reformatorerne ogsaa et stærkt
Indtryk af Verdens Forfængelighed, og Muligheden af,
at den yderste Dag kunde være nær forhaanden, stod
levende for dem; men det var dog fornemmelig Christi
aandelige Gjenkomst, hvorpaa de haabede. Den hellige
Skrift var det Speil, hvori hans Skikkelse tydelig kunde
sees; saa droge de Skriften frem, og Bibelen blev en
Folkebog. Hvorledes dette siden af Mange er blevet
misforstaaet, som om Meningen var, at man virkelig
kunde læse Christus til sig ved at læse i Bibelen, skal
her ikke omhandles. Hvad vi have for Øie, er de Betingelser,
hvorunder, og den Spændkraft, hvormed det
halvudslukte Aandsliv brød igjennem paa Ny. Christus
kom ikke ned fra Skyerne, men Menneskesønnens Tegn
% --- p. 50 ---
\setcounter{footnote}{0}%
\opage{50}begyndte at aabenbare sig i Menneskenes Indre. Det
var et levende Indre, der blev sat imod et dødt Ydre.
Pavekirkens Udvorteshed blev forkastet, og dens Verdslighed
dømt. Rom var atter „Babylon“, den store
Skjøge, som ved sin Trolddom bedaarede Folkene. Det
er Kampen, en \emph{aandelig}, men dog igjennem en Mangfoldighed
af verdslige Forhold \emph{folkelig} Kamp mellem
den gamle og den nye Kirke, der gav Reformationen
Spændkraft: Luther var i stor Stiil en Folkets Mand,
en folkelig Mand.

For en udvortes Betragtning seer nu den hele
historiske Bevægelse ud, som om Christi Gjenkomst
ikke skulde skee ovenfra, fra Himlen, men meget mere
nedenfra, fra det aandelige Dyb, hvori Livet grunder.
Fra dette Synspunkt kunne de, som ansee Historiens
Love for blotte Naturlove, paa deres Viis forklare sig
det Hele. „Aanden,“ sige de, „er intet Andet end den
menneskelige Naturs eget indre, almindelige Væsen.
Det var dette rene Menneskevæsen, som for Apostlenes
Forestilling havde antaget Kjød og Blod i Mennesket
Jesus Christus. Da nu Aanden som Aand tillige er
det Guddommelige, saa troede hine af Mesterens Liv
og Lære opfyldte Disciple, at Christus som Guds Søn
selv var baade Gud og Menneske. Det Rige, han
havde grundlagt ved at paatage sig Tjenerens Skikkelse,
vilde han da efter sin Ophøielse fuldende og føre til
Seier ved at komme som Herren i Himlens Skyer med
Kraft og megen Herlighed. Men han kom ikke; der
skete noget Andet. Istedenfor det forventede Jerusalem,
den hellige Stad, der skulde dale ned fra Himlen,
reiste sig et jordisk, og dog himmelsk, Jerusalem paa
Babylons, den gamle Skjøges, Ruiner: den katholske
% --- p. 51 ---
\setcounter{footnote}{0}%
\opage{51}Kirke var, med andre Ord, Virkeliggjørelsen af den forventede
Gjenkomst. Dog, Menneske-Aanden forstod
endnu ikke sig selv. Jorden skælvede under voldsomme
Rystelser af et sønderrevet Samfund; man stirrede op
i Skyerne, og troede fuldt og fast, hvad alle Tegn tydede
hen paa, at nu maatte da Himlen aabne sig, og Dommeren
komme med Englene og alle de Hellige. Men han
kom ikke; der skete noget Andet. Træt af at vente
paa det himmelske Jerusalem, gav man sig i Færd med
at erobre det jordiske. Det var en Omsætning af
Phantasie i Virkelighed. Christi Gjenkomst forvandlede
sig til Pavedømmets Opkomst. Livet med de rørige
Kræfter kom indenfra; Begeistringen tændte, det var
som et Flammehav af sandselig Aand og hellige Lidenskaber.
Først ved Reformationen begyndte Menneske-Aanden
at ane, hvorledes det egenlig hænger sammen
med Gjenkomsten. Menneskesønnen er Mennesket selv,
og Menneskets Selv, det sande, for Gud retfærdige
Menneske, er i Hjertets Indre, i Troen. Men ogsaa
Reformatorernes Tro var middelalderlig behæftet med
Forestillingen om et Gudsrige, der stod adskilt fra det
verdslige Liv. Derfor kunde Reformationen med al
den Folkevækkelse og Sædernes Rensning, den umiskjendelig
fremkaldte, ikke føre Sagen igjennem. Dens
officielle Kirkevæsen gik ind under Staten; Staten
gjorde det Hellige til Middel, og Thronerne af Guds
Naade trykkede Folkene ned til Masser. Først da
Masserne reiste sig og sprængte den absolute, sig selv
forgudende Stat, kom Aanden til Bevidsthed om sig
selv, og den overnaturlige Gjenkomst fik en naturlig
Forklaring. Menneskerettighederne: Frihed, Lighed og
Broderlighed, det er Humanitetens Tegn. Dette Tegn
% --- p. 52 ---
\setcounter{footnote}{0}%
\opage{52}er ikke blandt Engle i Himlens luftige Skyer, det er i
Støvskyer midt iblandt Mennesker hernede paa Jorden.
Hvad Tegnet betegner, skal Fremtiden vise, naar Kræfterne
røre sig og der kæmpes for Menneskelighedens
hellige Sag med virkelige, menneskelige Vaaben.“

Denne Forklaring har i Nutidens Øine noget Tiltalende,
fordi den, skjøndt ophøiet, dog svæver saa lavt
og holder sig saa nær til det Virkelige. Der er noget
Sandt i den, men Sandheden er kun halv; naar en
halv Sandhed udgives for en heel, saa have vi Løgnen.
Det er ganske sandt, at Aandens Vei er Inderlighedens
Vei, at det kraftige Indre maa slaae ud og give sig
Form i et virkeligt, nærværende Ydre. Men det er
ikke sandt --- det er i sørgelig Maade usandt --- at
de store Tanker om Frihed, Lighed og Broderlighed,
disse Evighedstanker, der omslutte alle Menneskelighedens
timelige Opgaver, kun skulde være selviske
Naturtanker, undfangne i Forbittrelse og fødte i Oprør;
de saa kaldte nye Tanker ere gamle Tanker: det er
forvrængede Christustanker. Det var ikke for Intet, at
Christus traadte op som de Fattiges Ven, og antog
sig de Ringe. Kunde de nye Tanker med de nye Fordringer
virkelig sættes igjennem ved en Fornegtelse af
Christi Aand, da vilde Revolutionsmændene fra 1789 og
1871 have Ret: jo blodigere, desto bedre. Men Historien
møder som Vidne, og dens Vidnesbyrd er.: % print: „er.:“
at det ikke
blot er Tyranniet, men ogsaa Folkefriheden, % defect: „Folketriheden“ (the head of the f did not ink, as pp. 72, 143, 157): the letter supplied
som under
Menneskerettighedernes røde Fane kan kvæles i Blod.

Vi befinde os med al vor „moderne Viden“ ingenlunde
paa Menneskelighedens, Frihedens og Fremskridtets
lige Vei; vi ere tvertimod ved at slaae ind
paa en Afvei. Vor troesfornegtende Tilstand er et
% --- p. 53 ---
\setcounter{footnote}{0}%
\opage{53}slaaende Modbillede til Middelalderens Overtro. Middelalderen
phantaserede med det Overhistoriske, Nutiden
vil hjælpe sig med det blot Naturhistoriske. Sandheden
er, at det Overhistoriske maa bære det Historiske, at
Aanden fra Himlen maa være med i alle kraftige Gjerninger
paa Jorden. Det er Bevægelsernes Lov; det er
Folkelivets Grundlag; det er i en Hovedsum, hvad
Christus lærer os om den historiske Fremtid.

% --- p. 54 ---
\setcounter{footnote}{0}%
\lecture{Fjerde Foredrag.} {(Fredericia).}
\opage{54}„Ondt og Godt i alle Lande, Roser sig med Torne
blande“, det er en gammel Vise; den passer ikke blot
paa de mange Lande, men ogsaa paa de mange Tidsaldere;
vor egen nærværende Tid har Roser, men ogsaa
Torne; den har stærke Lyssider, men ogsaa betænkelige
Skyggesider.

Vor Tid er Oplysningens Tid; unegteligt! Natten
selv bærer Vidne derom. Vore Gader og Huse, ere
jo saa straalende oplyste, at vi ikke kunne andet end
føle en vis Stolthed ved at tænke paa, hvor mørkt det
har seet ud i „Kong Hans'es Tid“. Og Forstandens
mange Lys brænde lige saa klart. Hvormeget er der
ikke allerede gjort for en forbedret Ungdomsunderviisning;
i hvilken Grad er Sands og Smag for Kundskab
og Dannelse ikke tiltaget i de sidste to Menneskealdere.
Synskredsene udvide sig; Opmærksomheden er
ikke længer indskrænket og fangen af de nærmeste
Hverdagsanliggender; vore Tanker vandre vidt omkring;
fremmede Egne, fremmede Sprog og fremmede Sæder
indbyde til Sammenligning med vort Eget. Vi forstaae,
at det, som foregaaer i fjerne Lande, kan vedkomme
% --- p. 55 ---
\setcounter{footnote}{0}%
\opage{55}os: vi føle Pulsslaget i det store Hele, og Verdensbegivenhederne
tegne sig i Omrids saa klare, som
aldrig før.

Men Skyggen følger med Lyset. Det Mangfoldige,
der adspreder Sindet, kan saa let fortrænge det første
Eenfoldige; og dog er det just det Eenfoldige, der
giver Samling, Inderlighedskraft og besluttet Villie.
Istedenfor hele Tanker, hvori Sjælen kan hvile, udklækker
Kulturen en Mængde livligt spillende, men
ueensartede Forestillinger, en Vrimmel af Tankebrøker,
til hvilke man forgjæves søger en Generalnævner. At
vide Lidt om Alt er ikke det Samme som at vide
Noget tilgavns. Den, der skal være med i Alt, for
ikke at være eensidig, betages af urolig Higen; og den
Nyfigne, der absolut skal have Nys om Alt og en
Mening om Alt, er sjælden i Stand til at fælde en
grundig Dom om Noget. Det Værste er, at det gaaer
ud over Charakteren. Det Menneske faaer aldrig
Charakteer, som ikke er i Stand til at danne sig en
fast Overbeviisning. I en Tidsalder, hvor en given
Livsanskuelse hersker med myndig Eensformighed ligelig
for Alle, seer det jo rigtignok ud, som om hver
Enkelt i Samfundet var født med Overbeviisning; men
under den stigende Oplysning begynder det Eensformige
at løse sig op. Mange Livsbetragtninger,
mange og modstridende Synsmaader stille sig til Skue,
og hvad bliver Følgen? Der er, siger man, Noget i
det Ene; men der er ogsaa Noget i det Andet: Ondt
og Godt i alle Lærdomme! Det Eneste, man faaer ud,
er, at der kan siges Meget for og imod enhver Ting.
Tilsidst bliver Alt uvist; man maa gaae efter Indtryk
og Stemning. Men Indtrykkene forandre sig og Stem%
% --- p. 56 ---
\setcounter{footnote}{0}%
\opage{56}ningerne vexle: Alt beroer paa Omstændighederne.
Den, som for ti Aar siden var ærkeliberal, er nu alvorlig
mørk og stokconservativ: det er Omstændighederne.
Grantræet, som før var Juletræ, staaer nu paa Loftet,
og er „sat i Væxten“: det er ogsaa Omstændighederne.
Hvor uhyggeligt at møde sin første oprindelige Anskuelse
paa Udviklingens Vei og saa ikke at kunne
vedkjende sig den igjen; det er jo næsten lige saa
ærgerligt som at skulle møde sin „første Kjærlighed“
og see, at hun er blevet forlovet med en Anden. At
tvivle om Alt, kritisere Alt og være misfornøiet med
Alt bliver ofte anseet for et ubedrageligt Tegn paa
Overlegenhed; og dog er denne Overlegenhed lige saa
bedragerisk som Rødmen paa den Svindsottiges Kinder.
Den endeløse Tvivlen er en tærende Sygdom og i Oplysningens
Tid er Sygdommen smitsom.

Vor Tid er en Frihedens Tid; vi erkjende det
med Taknemmelighed og sige det med Stolthed. Men
Frihedens Tid er ogsaa Farernes Tid, thi i Friheden
er der mange Slags Muligheder: at omgaaes med Frihed
er at omgaaes med Muligheder, men at omgaaes
med Muligheder er som at omgaaes med Ild og Lys i
et Krudttaarn.

I ældre Tider, ja endnu langt ind i det attende
Aarhundrede, var hvad nu kaldes Frihed en forbuden
Frugt. Den folkelige Frihed var under Laas og Lukke,
dens Tanker laae gjemte i Hemmelighedernes Bog, en
Bog med syv Segl. Vee den, der vovede at bryde et
af disse Segl, hans Dom var fældet. Selv med Religion
og Videnskab, Tro og Fornuft holdt Statsmagten strengelig
Tilsyn; den rugede over Aandens Klenodier som en
Drage over sine Skatte. Naar vi Protestanter korse
% --- p. 57 ---
\setcounter{footnote}{0}%
\opage{57}os ved at høre om alt det, den katholske Kirke har
faaet udrettet ved Hjælp af Censur, Inqvisition, Hexeprocesser
og Kjætterbaal, da bør vi for Alting ikke
glemme at tage vor egen Kirke med. Navnlig er det
i denne Henseende lærerigt at gjøre Bekjendtskab med
Englands Historie ved Tiden før og efter Reformationen.
Efter at Forfølgelserne med blodige Straffe havde udraset,
blev en døsig Orden i Høikirken og et mystisk
Sværmeri hos Secterne den naturlige Tilstand. Til en
Anskueliggjørelse heraf har Maleren Hogarth leveret
et Billede, hvori Satiren overgaaer Konsten. I Forgrunden
seer man Sectaanden drive sit Uvæsen. En
Kvinde i folderige Klæder ligger henkastet paa Gulvet,
og ud fra Folderne vrimler der en Mængde Smaadyr,
Muus, saa vidt jeg mindes. Grupperne omkring hende
ere fulde af Liv, Nogle i stille Henrykkelse, Andre i
en saadan Bestyrtelse, at Haarene staae som Børster
paa deres Hoveder. Det er et Mirakel, der skeer;
den benaadede Kvinde føder Muus. Altsaa: et stort
Musemirakel! Baggrunden er helliget Høikirken. Menigheden
er samlet, Præsten paa Prædikestolen, her er
Orden og Ro. Som befalet, læser Sjælehyrden sin
Prædiken op af Papiret. En Strimmel deraf underretter
os om, hvad det er for en Text, hvis Indhold,
han lægger de Andægtige paa Sinde. „I skulle ikke
mene, at jeg er kommen at bringe Fred paa Jorden;
jeg er ikke kommen at bringe Fred men Sværd.“ Orthodoxiens
„rene Lære“ er imidlertid saa betryggende,
at end ikke Ordet om Ild og Sværd kan forstyrre dens
Fred: Præsten seer søvnig ud, og den hele Menighed
er allerede sødelig indslumret.

Det, hvis Væsen er Frihed, kan ikke holdes under
% --- p. 58 ---
\setcounter{footnote}{0}%
\opage{58}Tvang; det er en Unatur, der hævner sig. Thi naar
der tales om Frihed, er det jo ikke nok med at tænke
paa Frihed i det Udvortes, saasom Næringsfrihed,
Trykkefrihed, en fri Forfatning o. s. fremdl. Den udvortes
Frihed er et uskatteerligt Gode; men den faaer
dog kun sit Værd ved at afgive Betingelser for den
indvortes. Uvidenhed og Overtro ere uforenelige med
indvortes Frihed. Videnskaben er efter sit Væsen fri,
thi Videnskaben vil erkjende Sandheden, og Sandhedens
Erkjendelse er i et Indvortes. At ville ved udvortes
Befalinger foreskrive Tanken, hvad den skal erkjende
for sandt, er det Samme som at ville befale, at To og
Tre ikke maa være Fem, men Fire. Religionen er
efter sit Væsen fri, thi Religion er i allerhøieste Betydning
et indvortes Anliggende. For ret at bruge
den udvortes Frihed, maae vi erhverve den indvortes.
Det Sind, som er betaget enten af Havesyge eller af
Herskesyge eller af hvilkensomhelst anden Syge, er
ufrit. Der er Mange, som have ført det store Ord om
Frisindethed, Sindets Frihed, og dog lidt utrolig af
Herskesyge: de have da ikke undt sig selv det Gode,
de saa gjerne vilde tilsikkre Andre.

Vor Tid er i særlig Forstand en Forhandlingernes
Tid; det følger med Oplysning og Frihed. Meget af
hvad man i tidligere Dage neppe turde tænke ved sig
selv eller hviske en Fortrolig i Øret, træder nu offenlig
frem i Skrift og Tale og drøftes for Alverdens Øine.
Dersom vi nu turde stole paa, at de mange Stemmer,
de mange Røster og de mange Meninger, der lade sig
høre i den stadig fortsatte Samtale, hvormed den
civiliserede Verden daglig underholder sig selv, virkelig
stode i Sandhedens Tjeneste, da havde vi rige
% --- p. 59 ---
\setcounter{footnote}{0}%
\opage{59}Midler til at faae de brændende Spørgsmaal belyste
fra alle Sider og komme til Bunds i de virkelige Forhold.
Men hvorledes er Fremgangsmaaden i det Hele?

Naar man vil prøve en ny Maskine, vurdere dens
Bygningsmateriale, dens Sammensætning og Arbeidskraft,
saa nøies man ikke med at høre paa de Lovtaler, hvormed
Opfinderen selv anbefaler den; lige saa lidt lader
man sig afskrække af deres Indvendinger, som vitterlig
have Interesse af at nedsætte og dadle Maskinen. Man
gaaer forsigtig til Værks, tager Erfaringen og de Sagkyndige
paa Raad, mærker sig nøie, hvad der taler for
og imod, og lader sig overbevise af Grunde. Det er
ikke Stemninger eller uforgribelige, forudfattede Meninger,
som her faae Lov til at gjøre Udslaget: nei,
det er Grunde, vægtige Grunde.

Hvad Forhandlingerne om det store, indviklede
Samfundsmaskineri angaaer, synes det derimod, som
om Nutidens oplyste og frie Nationer foretrak at slaae
ind paa en anden Vei. Anvende vi Sætningen om
Ondt og Godt i alle Lande paa den nu saa rige, alle
Marker befrugtende, alle Enge overrislende, alle Veie bestraalende
Bladlitteratur, da maatte det jo hedde: Ondt og
Godt i alle Blade! Men nei! Naar vi see Partierne
gruppere sig til Høire og Venstre, yderste Høire, yderste
Venstre, høire Centrum, venstre Centrum o. s. v., saa
tager det sig livagtig ud, som om den gamle Sætning
med Eet forandrede sig og betød: lutter Godt i nogle
Blade, lutter Ondt i andre. De, der have forsvoret sig
til Høire, forsikkre, at der slet ingen Roser er i Venstrebladene,
men lutter Torne; til Gjengjæld maae de, der
have forsvoret sig til Venstre, naturligviis bekjendtgjøre,
og det baade med Tegninger og fede Typer, at Høire%
% --- p. 60 ---
\setcounter{footnote}{0}%
\opage{60}bladene ere visne Torne, og at den ægte, sande, uforfalskede
Rosenduft alene maa søges i Venstrebladene.
Saa danner der sig rigtignok en Overbeviisning og en
Mening; men den er af negativ Art. Overbeviisningen
er, at man har en Mening, og Meningen er, at man
ikke vil lade sig overbevise.

Det Skib, der bærer Nutidsforhandlingernes store
Ladning, er ikke dybtgaaende nok; det pløier ikke i
rum Sø, men holder sig helst, hvor der er flade Grunde.
Vi see det, naar vi kaste et Blik paa Forhandlingerne
om Religion og Videnskab. At intet Folk kan leve
uden positiv Religion, at Throner falde og Stater vakle,
naar de hellige Grundpiller tages bort, derom er Statskløgten
i den Grad enig med sig selv, at Personer, som
for eget Vedkommende vitterlig have opgivet al Tro
paa Kirkens Lærdomme, ofte ere de ivrigste Forsvarere
af den bestaaende kirkelige Orden. Religionen bliver,
kort sagt, offenlig anerkjendt for, hvad den er, en aandelig
Stormagt. Videnskabens Ret og Betydning maa
især i en Oplysningens og Frihedens Tid, være hævet
over enhver Tvivl: Videnskaben bliver da ogsaa anerkjendt
for, hvad den er, en aandelig Stormagt. Dersom
nu disse to Stormagter kunde forene sig om at
arbeide sammen til et fælles Maal, da skulde slumrende
Kræfter vaagne og en ny Dag bryde frem over Folk
og Stater; men det er netop Nutidens Kval, at der
imellem Religion og Videnskab har aabnet sig en Kløft
saa dyb og bred som aldrig før.

Det læres af Bibel og Kirke, at Verden er skabt i
sex Dage; men Videnskaben beviser, at der imellem
Tilblivelsen af de første Dyr og Tilblivelsen af de
første Mennesker maa ligge Millioner og Millioner Aar.
% --- p. 61 ---
\setcounter{footnote}{0}%
\opage{61}Det læres i Bibel og Kirke, at Sol, Maane og Stjerner
ere skabte for Jordens og Menneskenes Skyld; men
Videnskaben beviser, at Jorden er et forsvindende
Punkt i det uendelige Himmelrum, o. s. fremdl. Var
det endda blot Afvigelser i Opfattelsen af det Udvortes,
af Forhold i Tid og Rum, det fik at gaae; men Striden
angaaer Tilværelsens inderste Kjærnepunkt, den guddommelige
Magt og Villie. Religionen hviler paa Miraklet;
Videnskaben kan ikke erkjende Miraklet; hvorledes
skulle to saa indbyrdes splidagtige Aandsmagter
kunne virke sammen: hvad den ene bygger op, maa
den anden jo rive ned.

I gamle Dage var Striden en Hemmelighed; Kirken
sørgede for, at ingen videnskabelig Afvigelse fra den
orthodoxe Lære kom til Folkets Kundskab eller at der
dog i al Fald, hvis en dristig Tænker talede for høit,
statueredes et afskrækkende % defect: „aískrækkende“ (the f broken): the letter supplied
Exempel. I vor Tid, Oplysningens,
Frihedens og Offenlighedens Tid, kan Hemmeligheden
ikke bevares: hvert andet Menneske veed
at Naturlæren ikke retter sig efter, hvad der staaer i
Bibelen. Ikke desto mindre bevæge de offenlige Forhandlinger
sig henad Overfladen. Man nedsætter Commissioner
til at undersøge Kirkens og Statens Rettigheder;
men det falder ikke den oplyste Nutid ind at foranstalte
en alvorlig Undersøgelse om, hvordan det egenlig
staaer til med Religionens og Videnskabens Mellemværende.
Midt i Gjæringernes, Tvivlenes, Kritikens
og de urolige Tankers Tid slaaer man sig til Ro med
en Underviisning, hvori det Bibelske og det Naturvidenskabelige
paa den naiveste Maade anklage hinanden.
Man holder Kirke for dem, der ville troe, og
Videnskab for dem, der ville vide; men et afglattet
% --- p. 62 ---
\setcounter{footnote}{0}%
\opage{62}Sammensurium af Tro og Viden er ikke til Opbyggelse
for Folket.

„Slige Opgjørelser,“ indvender man, „lade sig ikke
fremtvinge.“ --- Her er slet ikke Tale om Tvang; det
her fordres, er Regnskab, Redegjørelse for, hvorledes
det forholder sig med de store bevægende Aandsmagter.
Skulde en saadan Redegjørelse være enten farlig eller
umulig, da maatte det staae daarlig til med Samfundets
Samvittighed. Den eneste Indvending, som med Grund
kunde gjøres, skulde da være den, at en saadan Forhandling
vilde blive saa lærd og langtrukken, at den
store Almeenhed ikke kunde følge med. Vi svare:
hvad der væsenlig vedkommer Alle, maa kunne forstaaes
af Alle; --- Folket kan forstaae det Høieste.
Hvad herved sigtes til, er ingen Autoritetserklæring,
enten fra deres Side, der repræsentere Kirkelæren, eller
fra deres, der repræsentere Videnskaben, men slet og
ret en Paaviisning af det Grændseskjæl, hvorved Religion
og Videnskab nødvendig maae skilles og mødes.

Det er fornemmelig tre Punkter, hvorom Undersøgelsen
dreier sig: \emph{Tingenes Oprindelse}, \emph{Menneskets
Adel} og \emph{Samvittighedens Dom}.

Om Tingenes Oprindelse lærer Religionen, at Verden
er \emph{skabt,} at den, med andre Ord, er et Værk af Almagtens
Villie. Dersom nu Videnskaben virkelig kunde
bevise, at Verden ikke er skabt, saa maatte man, om
end paa lemfældig Maade lidt efter lidt see at faae
Religionen afskaffet. Men hvad veed Naturlæren nu
egenlig om Tingenes Oprindelse? Intet. Erfaringslæren
tager Tingene som de ere givne; ved at udforske
deres Kræfter, Sammenhæng og Love opdager den, at
den ene Tilstand er gaaet forud for den anden. Den
% --- p. 63 ---
\setcounter{footnote}{0}%
\opage{63}første Tilstand, den yderste Grændse, hvortil den kan
naae, er, som bekjendt en uhyre „Taage“. Skal Aandslæren
(Philosophien) gjøre sig Rede for Taagens Oprindelse,
da maa den nødvendig forudsætte, at en virksom
Fornuft, en guddommelig Viisdom, en almægtig
formende Villie har været nærværende i Taagen. Men
fra dette Synspunkt kan man da umulig bevise, at
Verden ikke er skabt. Alt i Naturen udvikler sig fra
første Begyndelse af efter strenge Love. Kan Skaberen
forandre sine Love? Miraklet staaer som uudsletteligt
Grændseskjel mellem Religion og Videnskab; men den
foregivne Grundmodsigelse er hævet. Religionen har
paa sin Side Ret, naar den hævder Miraklet; thi det
inderste Oprindelige i Alt er den lovgivende Villies
Frihed. Videnskaben har paa sin Side Ret, naar den
paastaaer, at Lovene ikke kunne brydes. Men saa
sandt Naturlovene ere den lovgivende Friheds egne
Aabenbarelser, saa brydes de jo ikke ved at optages i
Aandslovenes Sammenhæng og underlægges den frie
Almagt.

Menneskets Adel og høie Bestemmelse udtaler
Religionen kort og fyndigt i det Ord: Mennesket er
skabt i Guds Billede. Hvad har Videnskaben herimod
at indvende? Naturhistorien betragter --- og fra sit
Standpunkt med fuld Ret --- Mennesket som et Led i
Dyrerækken. Naturligviis indtager Mennesket den
øverste Plads, Hæderspladsen ovenfor Aber og Flagermuus.


Der gjøres store Anstalter til at bevise, at ligesom
de høiere Dyreformer lidt efter lidt have udviklet
sig af de lavere, saaledes har Mennesket udviklet sig
af en egen Dyreform, sandsynligviis en abeagtig Urform.
% --- p. 64 ---
\setcounter{footnote}{0}%
\opage{64}Darwin fremhæver, at Aberne have Disposition til
Brystsyge, ligesom Mennesket; at de blive vrede, naar
man gjør Nar ad dem, akkurat ligesom Mennesket; at
de kunne lære at ryge Tobak, ligesom Mennesket o. s. v.

Den Forkjærlighed, hvormed mange Naturhistorikere
for Tiden arbeide paa at udslette de sidste Spor af
\emph{oprindelig} Forskjel mellem Dyr og Menneske, er saa
varm, at man kun kan afkjøle den noget ved at henlede
deres Opmærksomhed paa Forholdet imellem
Videnskaben og Mennesket. Videnskaben har veiet
Kloderne, Videnskaben har maalt Lysets Hastighed,
Videnskaben har opdaget Stjerner, der sidde saa høit,
at det Lys, de sende ned til os, maa være Aartusinder
underveis, og det uagtet Lyset tilbagelægger en Vei af
over 40,000 Miil i hvert Sekund. Det er høie Ting.
Naar Naturhistorikeren nu betænker, at Videnskaben er
menneskelig, er et Aandsprodukt af den menneskelige
Hjerne, kan han endog gaae over til den modsatte
Yderlighed og henrives til en saadan Begeistring, at
han sætter Menneskets Hjerne i Guds Sted og forguder
den menneskelige Aand. Hvilken Afstand imellem
Dyret og Mennesket! Hvormange Millioner- Billioner-Trillioner
Aar maa der ikke hengaae, inden Abens
Hjerne bliver saa udviklet, at den kan veie Kloderne,
inden Flagermusene drive det til at udmaale Lysets
Hastighed.

Den hele Allarm med Menneskets dyriske Afstamning
reiser sig ganske simpelt derfra, at en berømt
Naturforsker af paaskjønnelsesværdig Iver for at paavise
Slægternes Sammenhæng og gradvise Overgange
har glemt, at der er Noget til i Naturen som hedder
\emph{Anlæg.} Det er i Kraft af sit guddommelige Anlæg, at
% --- p. 65 ---
\setcounter{footnote}{0}%
\opage{65}Mennesket fremtræder som et Guds Billede paa
Jorden.

Om Samvittighedens Dom lærer Religionen, at den
er en umiddelbar og naturlig Tilkjendegivelse af den
alvidende, hellige og retfærdige Gudsvillie i Menneskets
Indre. Skulde Videnskaben gjøre Indsigelse mod denne
Lære, da maatte den først bevise, at en alvidende,
hellig og retfærdig Gud umulig kan være til; men Opgaver
af denne Natur ligge udenfor Videnskabens
Omraade eller, hvad der siger det Samme, udenfor
dens Rækkeevne. Det er ikke med Samvittigheden
som med et dyrisk Instinkt, at den strax er færdig fra
Naturens Haand og altid holder sig paa samme Trin.
Samvittigheden vækkes og udvikles paa det Grundlag
af religiøse, sædelige Erkjendelser, som Samfundet frembyder,
og betinges i det Hele af Forholdet mellem Religion
og Videnskab. Det virkelige Liv med dets mange Arbeider
og Pligter maa, saa at sige, bestraales af to Sole:
hvad det indre Personlige angaaer, af Religionen, hvad
Verdenslovene og de virkelige Forhold angaaer, af Videnskaben.
Samvittigheden kan gjøre et Menneske eensomt;
men den kan ogsaa hellige og styrke de Baand, som knytte
Menneskene til hverandre: der gives lige saa vel en
offenlig som en privat Samvittighed.

At Samvittigheden neppe kan regnes iblandt de
Helgene, der hidtil have fundet særlig Dyrkelse i
Politiken, tør maaskee ansees for almindelig indrømmet.
Det er mærkværdigt at see, hvorledes Staterne overalt,
hvor deres „Ære“ og „Interesser“ krævede det, have
forstaaet at bevare fuldstændig Uafhængighed lige over
for de sædelige Forpligtelser, som den statskirkelige
Rettroenhed foreskriver. Af deres indbyrdes Forhand%
% --- p. 66 ---
\setcounter{footnote}{0}%
\opage{66}linger, Krige, Erobringer o. s. v. skulde man, dersom
man ikke vidste det i Forveien, vanskelig kunne gjætte
sig til, at de vare --- hvad de jo dog virkelig ere ---
christelige Stater. Saadanne Vaabenmærker som Løver,
Leoparder, Ørne, Dobbelørne o. desl. ere derimod mere
charakteristiske, thi deri er noget Rovdyragtigt.

Oplysningens og Frihedens Tid er allerede naaet
saa vidt, at den offenlige Mening er blevet en Magt:
Politiken, der er saa uafhængig af Samvittigheden, er
ingenlunde uafhængig af den offenlige Mening. Naar
det store Rovdyr vil frem, speider det til alle Sider og
trækker forsigtig Kløerne ind, hvis det mærker at den
offenlige Mening bliver alvorlig vred. Men den offenlige
Mening, et Produkt af Stemninger, Forstandsbetragtninger,
Interesser og Lidenskaber, er med al
sin sunde Sands dog ingen paalidelig Veileder; den er
Øieblikkets Barn, for meget af et Verdensbarn. En
klarøiet Datter af beregnende Politik, vittig og klog,
men af de mange Beilere idelig udsat for allehaande
Fristelser, faaer den offenlige Mening undertiden en
betænkelig Lighed med den Slags Damer, man i Paris
kalder Demi-monde.

Samvittigheden, Samfundssamvittigheden, derimod,
er en paalidelig Veileder. Men Samfundssamvittigheden
danner sig først efterhaanden, som der danner sig et
folkeligt Sind og en folkelig Villie. Den offenlige Mening
kan ikke paa egen Haand faae Bugt med Magthaveriets
store Dyr. Samvittigheden maa med: naar
Dyrets Bevægelser forraade, at det hungrer efter Bytte
og tørster efter Blod, maa Folkevillien sige Nei, og
Folkets Sjæl igjennem Millioner Øine fæste sit Blik
paa Dyret. Det virker: Ørnens Vinge lammes, Leo%
% --- p. 67 ---
\setcounter{footnote}{0}%
\opage{67}parden glemmer sit Spring, det Rovdyragtige bliver
tæmmet.

Den naturvidenskabelige Betragtning af Historiens
Gang er ikke meget trøstelig. Det gaaer med Folkene
--- hedder det --- som med Individerne, de have deres
Aldere. Først en Ungdomstid med gyldne Syner, mægtig
Higen og stærk Tro; dernæst en Manddomstid
med skarp Forstand, sikkre Planer og en velberegnet
Daadskraft; tilsidst kommer Alderdommen med Kritikens
Rynker, Tvivlen om Alt undtagen det Materielle, Hang
til Ro og Nydelse, onde Drømme, Varsler om Opløsning
og Død. De samme Love, hvorefter Oldtidens
Stater sank, gjælde ogsaa for Nutidens. Europas
Folkeslag ere allerede stærkt paa Heldingen; det gaaer
ned ad Bakke med dem alle. Videnskaben gaaer frem,
men Religionen gaaer tilbage. Den videnskabelige
Forskning maa nødvendig opløse det kirkelig Overleverede;
den kan ikke gjøre for det. Den kan desværre
ikke erstatte Tabet. Videnskaben formaaer ikke
at give Folket, hvad det høilig trænger til, en ny Religion.
Almuen og en Deel af de Halvdannede holde
paa den gamle Tro, det bedste de kunne; men den
svinder umærkelig fra dem, den fordamper. Enden
bliver, at de nuhistoriske Folkeslag omsider maae træde
af, og nye Stammer med friske Rødder og unge Skud
udfylde Skuepladsen.

Vi maae ved denne Anledning atter fremhæve, hvad
der fra et naturvidenskabeligt Standpunkt saa let oversees,
at Historiens Love ere Aandslove. Folkene ældes
ikke efter samme Lov som Individerne; deres Alderdom
indtræder kun da, naar Aandskilderne udtørres.
Nutidens Folkeslag leve i denne Henseende under andre
% --- p. 68 ---
\setcounter{footnote}{0}%
\opage{68}Vilkaar end Oldtidens. Der er i Tidernes Vendepunkt
blevet offret et Liv; en Livskapital er sat ind, som
med Renter og Renters Renter er uudtømmelig. Videnskaben
skal ikke bebyrde sig med at give Folkene en
ny Religion; det behøves ikke. Ethvert Folk maa
blive ansvarligt for, hvorledes det har forstaaet at tilegne
sig og anvende den rige Arv.

Istedenfor det mørke Billede af lutter Skygger,
Svaghed og Alderdom kunne vi med mere Grund udkaste
et Foryngelsens Billede, et Fremtidsbillede med
lyse Farver. Vi vide vel, at Intet er fuldkomment paa
Jorden, at den virkelige Tilstand, hvis Træk skulle
svare til Billedet, kun vil naaes i en Tilnærmelse; men
derfor kan der dog være Sandhed i Billedet. Det er
glædeligt at haabe; ja, det er mere end glædeligt, det
er Pligt. Det er Folkenes Foryngelse, Folkelivets Fremvæxt,
hvorpaa vi haabe.

Naar den da kommer, den forventede Tid, da
Folket er blevet en fuldvoxen Mand, der tilfulde har
lært at bruge sine Kræfter, de aandelige saa fuldt som
de materielle, da ville Skranker falde, som nu synes
uoverstigelige. Bliver Samfundet, det store Legeme
med de mange Lemmer, befriet fra alt det, som kaldes
dødt Kjød, og faaer sin rette Førlighed, sunde Safter i
alle Væv, frisk Blod i alle Aarer: da vil der være
Dygtighed til at optage Virkelighedens Kamp, og Mod
til at møde de mulige Farer. Gaaer den sædelige Selvfølelse
op for Alles Sind og Tanke, er Folkesamvittigheden
saa vaagen, at Selvraadigheden maa blive undseelig,
Egensindigheden krybe i Skjul, og Kortsynethed
og anmassende Dumhed vige overalt for bedre Indsigt
og for klare Grunde: da bliver Enigheden stor og
% --- p. 69 ---
\setcounter{footnote}{0}%
\opage{69}Villien stærk; den voxer jo med Forstaaelsen. Kommer
den engang, den forventede Tid, da Arbeidets Byrde
og Arbeidets Løn skulle veies mod hinanden efter Billighedens
Krav paa Retfærdighedens Guldvægt, da ville
mange Hænder arbeide til eet Maal, og de mange
Interesser gribe ind i hverandre som Hjul og Drev i
en vel sammenføiet Maskine. Er det een Sjæl, een
Aand, een Livets Strøm, der driver Værket med de
mange Hjul, da vil Byrden ikke blive Nogen for tung,
og Ingen skydes udenfor: det rummelige Værksted har
Plads og Virkekreds for Smaa og Store. Ja, naar den
kommer, den forventede Tid, det bliver en Sommerens
Tid: lette Skyer og svale Vinde, lyse Nætter og solrige
Dage. Det bliver en Høstens Tid; da bugner Kjernen
i
tunge Ax; da plukkes der modne Frugter af
Frihedens Træ; da har Kundskabsæblet ingen Gift; Naget
er kvalt, Hugormen i det høie Græs er død: da dufte
alle Livets og Sundhedens Roser.

% --- p. 70 ---
\setcounter{footnote}{0}%
\lecture{Femte Foredrag.} {(Kolding).}
\opage{70}Et sundt og kraftigt Folkeliv er umuligt uden tilsvarende
Folke-Aand, Folke-Aand er umulig uden Folketro;
men Folketroen er gold, vanmægtig og intetsigende,
naar den ikke formaaer at bevise sin Kraft i Gjerning.
Hvad vi her kalde Folketro, er ikke Overtro, Almuetro
paa gamle Sagn om Underjordiske, Gjengangere, Troldmænd
og Hexe, en Tro fra Ammestuen; Folketroen er
en Tro, som taaler Lyset, en Tro paa Lysets, Livets
Kraft, en Tro paa Aand. Naar der ellers spørges om
et Folks Tro, mener man sædvanlig dets Troesbekjendelse.
Man kommer da let til at gjøre den Iagttagelse, at
man ikke uden videre tør slutte fra Tro til Liv, at
Troesbekjendelsen som oftest i den Grad peger ud
over Livet, at man maa vælge en anden Synslinie,
naar man vil have det virkelige Liv i Sigte.

Folketroen er mere end en ophøiet Følelse eller
en skjøn Phantasie; den er en Virkelighedstro, en
levende Gjerningstro; som Troen er, saaledes er Livet,
og som Livet er, saaledes er Troen; her gjælder det
ubetinget: viis mig din Tro af dine Gjerninger. Vistnok
er det Guddommelige, hvorpaa der troes, et Over%
% --- p. 71 ---
\setcounter{footnote}{0}%
\opage{71}menneskeligt; men dette Overmenneskelige er efter sin
umiddelbare Nærværelse dog ganske og aldeles menneskeligt:
alle Hovedtræk af Menneskets Natur og Væsen
ere kjendelige i dets Aabenbarelser.

I denne Henseende kunde det da synes, som om
de hedenske Guder, de gamle Folkeguder, passede
langt bedre til Menneskets Natur og Folkeslagenes
Eiendommeligheder end Aandens evige sande Gud, den
Gud, der har aabenbaret sig i Christus. De græske
Guder have græske Ansigter, græske Dyder og Lyder;
de ere græske Mennesker af høiere Orden. Noget Lignende
gjælder i det Hele om Romernes og Nordboernes
Guder. Men just fordi det kun er Menneskeguder,
Guder med menneskelige Feil og Fortrin, bære de
Dødsspiren i Hjertet og kunne med al deres Udødelighed
kun leve til en vis Tid.

Det var Asalæren, Troen paa de høie Valhals Guder,
der holdt Nordens Aand oppe og bar Folkelivet frem
paa stærke Skuldre; men kun til en vis Tid. I 8de
og 9de Aarhundrede efter Christi Fødsel er Forholdet
vendt om; da er det Menneskene, der bære Guderne.
Helten stoler paa sin stærke Arm og sit skarpe Sværd:
han troer paa Thor, fordi han troer paa sin egen
Tapperhed. Hvad betyder vel det Hemmelighedsfulde
i Gudernes Skjæbne, som kun Nornerne kjende, Gaaden,
hvorover Baldur grunder i Veemod, Mimers Brønd,
hvor Odin har sat sit Øie i Pant, Andet end det
dunkle, uoplukkede Dyb i Nordboens eget Indre. Med
alt deres Mod, Selvstændigheds- og Frihedssands føle
de stærke Mænd sig dog bestandig undergivne mørke
Magter; der er en ubøielig Nødvendighed, noget Koglende
og Troldomsagtigt, der hilder dem i sit Garn.
% --- p. 72 ---
\setcounter{footnote}{0}%
\opage{72}De ere tause, indesluttede, ordknappe, naar de udruge
Planer og skride til Handling. Sindet er tynget, det
er Nødvendighedens Aag, det ufrie, naturbundne Væsen.
En Sigurd Fofnersbane kan dræbe Dragen og hæve
Skatten, men Skæbnens Lindorm, den gaadefulde Tilværelses
store Drage, formaaer end ikke Folke-Aanden
selv at dræbe. Saa opkommer Anelsen om en Alfader:
Harald Haarfager troer paa en eneste, almægtig Gud;
saa er der midt i Kampens Larm en stille Længsel
efter Fred, en Søgen efter Lys, et Suk efter Befrielse.

I sin uforlignelige Romance: „Hakon Jarls Død“, har
Øhlenschlæger fremstillet denne Overgang paa sindbilledlig
Maade. „De Nætter ruge saa lange og sorte,
Syvstjernen skimter saa mat“ . . det er en nordisk
Vinternat, Hedenskabets mørke Nat, hvis kolde Aande
vi fornemme. % defect: „tornemme“ (the head of the f did not ink): the letter supplied
Omkring Offerhøien staae Guderne, gamle
og graae, i sprukne Stene: „Vor Tid er forbi, snart
synke vi“. Modbilledet hertil er et Skyggerids af den
gothiske Kirke med de røde Mure. Gjennem Blyvinduet
naaer Maanens Straale til Alterets Crucifix og
lyser med en Forjættelse: „Du seirer vist, du hvide
Christ, snart skal den vilde nordlige Zone neie sig for
din Tornekrone“. Saa begynder Virkeligheden: „Olaf Trygvesen lander i Norge“. Messer synges, Korsbanneret
vaier, „Hellebarderne vildt i Ætheren vælte,
de lyne, de lyne afsted“; det er et folkeligt Sammenstød,
et blodigt Stævne mellem den gamle og den
nye Tro.

Seet fra et evangelisk Synspunkt kunne blinkende
Hellebarder og klirrende Sværd naturligviis ikke bruges
til Beviser for Christendommens Sandhed; men naar
Talen er om en Folketro, saa er Sandhed en Kraft,
% --- p. 73 ---
\setcounter{footnote}{0}%
\opage{73}Livskraft, Heltekraft: viis mig din Tro af dine Gjerninger.
--- Menneskelivet kan ikke bære sig selv og endnu
mindre bære sine egne Guder. Sandhedens hellige
Gud, som bærer Alt, er uendelig ophøiet over det
Menneskelige; men just derfor ogsaa --- som Kjærligheden
selv --- Menneskenes sande Gud, en sand
Menneskegud. Det sande Livets Ord maa indpræntes
i Hjerterne og vinde Sjælene.

Over halvandet hundrede Aar før Olaf Trygvesens
Optræden i Norge havde Ansgar, en Fredens Apostel
fra Klostercellen, forkyndt Ordet hos os. Mild som
Evangeliet, sagtmodig, selvfornegtende som den, der
træder i Frelserens Fodspor, modig-taalmodig som den
af Aanden Grebne, der higer efter Martyrkronen, havde
Ansgar ved sit Ords Varme allerede begyndt at smelte
Isen, optøe de kolde Sind og tænde Alterlys omkring
de Brynieklædte. Saa gik det ad begge Veie, Ordets
og Sværdets, fremad med Folketroens Omdannelse.

Men hvor meget Evangelisk der end er i Ansgars
Lære og personlige Liv, maae vi dog ikke glemme,
hvorfra han kommer. Hans Evangelium er ikke det
oprindelig rene; der er sat en Farve paa det. Det er
Katholicismen, den romerske, endnu ikke fuldt udviklede
Katholicisme, han bringer med sig til Nordens Omvendelse.


Hvor høit Katholicismen har løftet Folketroen,
hvilken Sindsforandring den har bevirket, hvor mange
Lænker den har løst, hvorledes den har skjærmet de
Værgeløse, adlet Kvinden, antaget sig de Fattige og de
Livegne og arbeidet paa at bringe Eenhed, Fred og
Orden i det vildt forvirrede Samfund, skal ihukommes
med Taknemmelighed. Vore gothiske Kirker: Stræbe%
% --- p. 74 ---
\setcounter{footnote}{0}%
\opage{74}pillerne, Spidsbuerne, de høie Hvælvinger og de slanke
Taarne, der pege mod Himlen, ere mere end Mindesmærker
om en egen Bygningskonst, de ere Mindesmærker
om et eget Folkelivets Trin, Vidnesbyrd om
Troens Kraft i forbigangne Tider. Fra disse tause Mure
taler en forstenet Aand om svunden Magt og Vælde.

Men Katholicismen har fra Begyndelsen begaaet
en Synd imod Aanden; der er en Feil, en eneste
Grundfeil, hvorfra alle dens Overgreb, Misgreb, ja selv
dens Aandsforbrydelser have deres Udspring, den ligger
i Opfattelsen af Begrebet: \emph{Myndighed.}

Efter Christi Forbillede skal det hellige Ord forkyndes
med Myndighed; men efter Christi Forbillede
skal da ogsaa denne Myndighed udtrykkelig stamme
fra Aand d. e. fra overbevisende Sandhed, opoffrende
Kjærlighed, personlig Overlegenhed. Hvorledes har
Katholicismen nu opfattet og anvendt Christi Myndighed?
Den har ganske rigtig taget Udgangspunktet
i Aanden. Men hvormeget den end har ivret for at
faae den Hellig-Aands Personlighed sat fast som Dogme,
saa er det dog umiskjendeligt, at den ved Ordningen
af det hele Kirkevæsen, er kommet til at gjøre, hvad
den ikke vilde, idet den i Virkeligheden er kommet til
at gjøre sin Hellig-Aand saa upersonlig som muligt.

Paa Dogmet om Aandens umiddelbare Nærværelse
i den ufeilbare Kirkes ufeilbare Bispeforsamlinger er
den katholske Lære om den geistlige Myndighed begrundet.
Kirkens Lærere og Styrere have da ikke den
guddommelige Myndighed, i Lighed med Christus selv
og hans Apostle, fra det Personlige i deres eget Indre.
Tvertimod! det religiøs-sædelige Indre kan, som hos
andre Dødelige, befinde sig i megen Skrøbelighed;
% --- p. 75 ---
\setcounter{footnote}{0}%
\opage{75}deres Myndighed lider intet Afbræk, fordi de ikke ere
Dydsmønstre. De geistlige Personer skulle ikke bedømmes
efter deres eget personlige Værd, men efter
den Stilling og Rang, de have i Kirken. De ere Kirkens
Organer, dens Tjenere og dog tillige dens Bærere:
deres Myndighed er Kirkesystemets Myndighed, personlig
i det Feilbare, upersonlig i det Ufeilbare. Kirkesystemet
er udvortes og de kirkelige Personer ligesaa.
Kirken kan ikke handle uden igjennem sine Personer;
det hele Kirkesystem maa da i Virkeligheden omfatte
baade Personer og Kirke. Men Personerne kunne jo
feile; for at det Ufeilbare kan befrie sig fra det Feilbare,
maa der skee et Mirakel, et Aandens Mirakel,
hver Gang Bisperne samles. At Bisperne endda ikke
saa sjældent have været samlede, uden at Miraklet er
skeet, er vitterligt, eftersom den ene katholske Bispeforsamling
jo høitidelig har fordømt den anden. Paastanden
om Miraklet og den dertil svarende Myndighed
bliver altsaa en reen udvortes Paastand. Myndigheden
bliver ikke mindre, men mere udvortes endnu, naar
en enkelt Bisp, en enkelt, i og for sig feilbar Person,
skal i Medfør af Aandsmiraklet og Kirkesystemet
omskabes til ufeilbar.

Paa et grundløst Mirakel blev den geistlige Myndighed
begrundet og bragt til Fuldendelse i Middelalderen.
Hvor udvortes Myndigheden var, sees af den
kirkelige Adskillelse mellem Geistlige og Lægfolk. Alle
Geistlige vare myndige; alle Lægfolk d. v. s. Ikke-Geistlige
vare --- i Aanden --- umyndige; dermed var
Sagen afgjort. At nu en saadan Myndighed ingenlunde
var, hvad man fra protestantisk Side undertiden har
villet udgive den for, et vilkaarligt Paafund af geistlig
% --- p. 76 ---
\setcounter{footnote}{0}%
\opage{76}Herskesyge, er klart af den hele kirkelige Tankegang.
Det skal villig indrømmes, at udvortes Myndighed netop
egnede sig for halvbarbariske Forhold i en halvbarbarisk
Tid. Det skal indrømmes, at der hørte
sandselig haandgribelig Myndighed til for at virke opdragende
paa de raae Masser. Men betragte vi det
Opdragende i denne Myndighed nærmere, da viser det
sig ogsaa, at det netop i sin Udvorteshed har sin
Grændse.

Kommer det an paa at frembringe pludselige og
stærke religiøse Indtryk, kue Sindsraaheden, vække
dybere Længsler, tugte Tvivlen og tage Gjenstridigheden
med Magt, da er der vist intet Kirkevæsen, der kan
maale sig med Katholicismen. Det Samme gjælder,
naar Opgaven er at nære en mystisk Inderlighed og
bringe det udvortes Liv med dets utallige Enkeltheder
ind under strenge almindelige Regler. Men længere
kommer man ikke; det bliver i Inderligheden ved det
Mystiske: den ydre Virkelighed maa nøies med det
Fromhedsskjær, som et farvet Lys igjennem Kirkens
Ruder udbreder over „de gode Gjerninger“. Messer
og Helgenbilleder, Chorsange og Røgelser, Klokkeringning
og Processioner, alt Sligt egner sig til at udfylde
en sandselig Phantasie og gjøre det Aandelige
udvortes. At knæle for Guds Moder, bede sin Rosenkrands,
ligge med bare Knæ paa Ærter for at afsone en
Kirkestraf kan for Umyndige være optugtende nok;
men det er en Synd imod Aanden, at Forstaaelsen
mangler, at de Umyndige aldrig naae saa vidt, at de
nogensinde kunne blive myndige.

Ved Billedstridens Begyndelse have vi et Ord af
Gregor den Store (590 --- 604); han siger: „Lad dem
% --- p. 77 ---
\setcounter{footnote}{0}%
\opage{77}læse paa Væggene, hvad de ikke læse i Bøgerne.“
Herimod vilde der Intet være at indvende, dersom vi
turde forstaae Ordene saaledes, at den foreløbige
Læsen paa Væggene snarest muligt skulde fuldstændiggjøres
ved en tænksom Læsning i Bøgerne; men at
opdrage et Folk til religiøs Selvtænkning vilde jo, hvad
Erfaringen tilstrækkelig har lært, være det Samme som
at modarbeide den udvortes Myndighed.

Magtsprog og udvortes Myndighed ere paa det
Nøieste forenede; den udvortes Myndighed vil, hvad
Christus og hans Apostle aldrig have villet, herske ved
Brug af udvortes Midler og udvortes Vaaben. Herfra
den lange forbittrede Kamp mellem Pavemagten og
Fyrstemagten; hiin var i Kirkens Øine som Solen, denne
som Maanen, der maa laane sit Lys fra Solen. Blandt
de mange udvortes Midler er intet saa virksomt som
Penge. Kirken trængte til Penge, altid Penge, flere
Penge. Naar nu ikke blot Penge, men ogsaa Syndsforladelse
er et Udvortes, saa var der, skulde man
mene, jo heller ingen Modsigelse i at ombytte det ene
Udvortes imod det andet, altsaa slet ingen Modsigelse
i at sælge Syndsforladelse for Penge. Men --- mærkeligt
nok --- denne klare Logik kunde Folkene ikke
forstaae. At Paven tog deres Sjæle, deres Samvittighed,
deres indre Frihed og aandelige Selvstændighed, fandt
de sig i med Taalmodighed som gode Christne; men
at han tog deres Penge, det var oprørende. Imod det
Oprørende gjør man Oprør; tilsidst var Maalet fuldt:
Reformationen kom til Gjennembrud.

Paa Rigsdagen i Vorms (1521) saae Verden et
Exempel, som viste, hvad Aandsmyndighed formaaer,
om den end møder i udvortes Ringhed, og fremtræder
% --- p. 78 ---
\setcounter{footnote}{0}%
\opage{78}under en uanseelig Skikkelse. Man forestille sig den
glimrende Forsamling; Keiseren, Fyrsterne, den pavelige
Legat, Cardinalen, Ærkebiskopper, Biskopper,
Abbeder, hvilke Repræsentanter for verdslig Storhed,
for kirkelig Magt og Myndighed! Og saa den ringe
Munk, den Anklagede, paa hvis Forsvar man ikke vilde
høre, men hvis Ulydighed, man strengelig vilde straffe,
dersom han ikke ydmygede sig under Kirkens Ufeilbarhed
og tilbagekaldte Alt, hvad han havde talt og
skrevet. Det var et stort Øieblik, da den Trierske
Official faldt Luther i Talen og forlangte et slet og
ret Svar, om han vilde tilbagekalde eller ikke; thi
Luther svarede: „Med mindre jeg med Vidnesbyrd af
den hellige Skrift eller med offenlige, klare og indlysende
Grunde bliver overbeviist og fangen i min Samvittighed
af Guds Ord, saa kan og vil jeg Intet tilbagekalde,
thi det er hverken sikkert eller raadeligt at gjøre Noget
imod sin Samvittighed.“ Her var en nyfødt Samvittighed
med en ny Myndighed; det Indre brød frem som
Lynet af den mørke Sky, og Lynet tændte.

Hvilke Vækkelser og Forhaabninger, hvilke Strømme
af aandeligt Liv, Reformationen medbragte, er almindeligt
anerkjendt. Luther gjorde Reformationens Sag til
en Folkesag. „Enhver sand Christen er af geistlig
Stand; Kalken i Nadveren tilkommer Folket; det skal
staae Enhver frit for at troe den nyopfundne Forvandlingslære
eller hvilketsomhelst virkeligt Fællesskab
mellem det indviede Brød og det guddommelige Legeme.
Mod hele Betydningen af den udvortes Kirke sættes
Troens alene saliggjørende Almagt. Derfor maa den
christelige Frihed ikke bindes ved nogen Sætning og
vilkaarligt Løfte“ o. s. v. Men den nye Kirke, som
% --- p. 79 ---
\setcounter{footnote}{0}%
\opage{79}efterhaanden dannede sig, turde ikke staae vaabenløs
som Luther paa Rigsdagen. Gjennem den hellige
Skrift gik man tilbage til det Oprindelige: istedenfor
Pavens Myndighed gjaldt alene Christi Myndighed, og
Apostelmenigheden blev den evangeliske Kirkes Forbillede.
Dersom denne hos Luther oprindelige Livstanke
skulde have været gjennemført, saa maatte de
unge lutherske Menigheder ogsaa have fundet sig i at
blive staaende ved Christi og Apostlenes Livsvilkaar;
men frivillig at vælge Lidelsen og Martyriet efter en
saadan Maalestok, dertil var der med al den Troens
Inderlighed og det Opoffrelsens Heltemod, hvorpaa hiin
Tid har saa mange og skjønne Exempler at opvise,
dog ikke Aandskræfter nok. Saa lod man de for aandelig
Myndighed afgjørende Betingelser blive udenfor,
slog af paa Forbilledets Efterfølgelse og holdt sig til
Læren, „den rene Lære“. Trods sin næsten ubændige
Iver for den rene Lære fandt Kirken sig til Rette med
de historiske Omstændigheder og gav sig ind under
Staten. Den indvortes Myndighed gav tabt og faldt i
Armene paa den udvortes: saaledes blev vor lutherske
Kirke en Statskirke.

Dersom Folket her til Lands kunde enten oversee
eller miskjende, hvor meget Ædelt og Godt, vor lutherske
Statskirke i Tidens Løb har fremmet hos os, hvor
meget der saa vel fra gudsfrygtige Kongers som fra rettroende
Geistliges Side er udrettet til Sjælenes Trøst,
til Sædernes Rensning, til stille Fremgang i Aandskultur
og almeen Oplysning: da maatte vort Folk i
sørgelig Maade mangle --- hvad det lykkeligviis ingenlunde
mangler --- historisk Sands. Ikke at tale om,
at det var Statskirken, som med eet Slag slog Hierar%
% --- p. 80 ---
\setcounter{footnote}{0}%
\opage{80}chiet til Jorden, saa kan jo næsten ethvert Skridt
fremad, enhver Foranstaltning, hvorved den nærværende
Frihedstilstand er blevet forberedt, umiddelbart eller
middelbart føres tilbage til Statskirkens religiøse og
borgerlige Indflydelse. Naar et Folk staaer i Begreb
med at danne og fremarbeide noget Nyt, da er det en
Taknemmelighedspligt at mærke sig de Fortjenester,
som det Gamle, der nu skal skriinlægges, virkelig har
havt. Ved at opfylde denne Pligt sikkrer man sig
bedst imod det Overilelsens Feilgreb at bortkaste netop
det, som nødvendig maa bevares.

I eet Punkt har den lutherske Statskirke syndet;
den har i Lighed med Katholicismen sat den udvortes
Myndighed over den indvortes. Luther har med stor
Styrke fremhævet, at ethvert Sakrament, som forvaltes
i Kirken, udtrykkelig maa have Christi Befaling. Da
dette kun gjælder om Daaben og Nadveren, forkastede
han Katholikernes andre fem Sakramenter som selvgjorte.
Den gamle Statskirke sluttede nu saaledes:
have Sakramenterne Christi Befaling, saa maa det fremfor
Alt være den christelige Øvrigheds Pligt, at Christi
Befaling punktlig bliver opfyldt. For at sikkre sig dette,
blev der da føiet en Øvrighedsbefaling til Christi Befaling:
den udvortes Myndighed traadte til for at hjælpe
paa den indvortes. Det var et uhyre Misgreb, for ikke
at bruge det haarde Ord: en Aandsforbrydelse. Hvor
Christus befaler, taaler han ingen Indblanding af andre
Befalende; hans Myndighed krænkes, naar uvedkommende
Hjælpere træde til og ville understøtte den med en
Myndighed af ganske anden Art. Christi Befaling lyder
paa Tro, og Troen er fri: fri Tro, fri Lydighed, fri
Tilegnelse ere det religiøse Livs første, ufravigelige
% --- p. 81 ---
\setcounter{footnote}{0}%
\opage{81}Grundbetingelser. Havde Christus villet sætte sine Befalinger
igjennem ved Anvendelse af Tvang: da havde
han jo Magt, og kunde have brugt Magten. Men at
han ikke vil tillade, at nogensomhelst udvortes Myndighed
under nogetsomhelst Paaskud indblander sig i
det, som efter sin Natur hører ind under Aandsmyndigheden,
har han i Gjerning beviist ved at gaae Lidelsens
Vei og lade sig korsfæste. Luther forstod Christus,
da han erklærede, at det skulde staae Enhver frit for
at troe --- man tænke sig --- „den nyopfundne Forvandlingslære“! Men, desværre, den lutherske Statskirke
forstod ikke Luther.

Ligesom i Katholicismen saaledes har ogsaa i den
protestantiske Statskirke den udvortes Myndighed gravet
sin egen Grav; men, det forstaaer sig, Graven er en
anden, thi her er en anden Graver.

Det attende Aarhundrede har i særegen Forstand
faaet Navn af Oplysningens, navnlig Forstandsoplysningens
Tidsalder. Med den fremadskridende Videnskab blev
Tænkningen frigjort; den undersøgende Forstand kastede
sig paa Alt; den vilde begribe Alt, og forkastede Alt, hvad
den ikke kunde begribe. Med blussende Fakkel forfulgte
Oplysningen Nisser, Trolde, Gjengangere og
Underjordiske, som flygtede alt hvad de kunde; tilsidst
bandt den an med Djævelen selv. Ifølge den orthodoxe
Grundsætning, skulde man tage Fornuften fangen under
Troens Lydighed; men Troens Lydighed var jo det
Samme som Lydighed mod den af Statskirken foreskrevne
Troeslære, en blind Underkastelse under en
udvortes Myndighed. En saadan Aandstrældom vilde
Fornuften ikke længer finde sig i. Ved Forstandens
Hjælp bragte Videnskaben det snart saa vidt, at den
% --- p. 82 ---
\setcounter{footnote}{0}%
\opage{82}selv blev øverste Domstol, og Religionen maatte nu
underkaste sig Fornuftens Myndighed.

Holberg, der maa betragtes som Rationalismens
første Forkæmper hos os, udtaler sig herom blandt
Andet saaledes: „Jeg haver viset, at Orthodoxie selv,
som grunder sig alene paa Andres Myndighed ikke kan
være Gud behagelig. Thi ligesom en stor Kunstner
finder ingen Behag derudi, at hans Arbeide prises af
en Blind eller uskjønsom Mand, som ingen Øien haver;
og en Musicus ikke kan glæde sig ved, at hans Musik
roses af een, der ingen Øren haver, saa kan Gud ingen
Behag finde udi dem, der dømme om hans herlige
Egenskaber, hvorom de ei have umaget sig at faae
mindste Idee. --- Skal en Lovsang eller Dyrkelse behage,
saa maa den gjøres af den, der veed, hvorfor
han dyrker; og skal en Tro kaldes Tro, saa maa den
grunde sig paa Kundskab, saa at den Talemaade: \emph{Jeg
troer}, maa være det samme som: \emph{Jeg veed.}

Forstandsoplysningen indskrænkede sig ikke til at
rokke ved den bestaaende Rettroenhed fra et Standpunkt
udenfor Kirken, den sneg sig ind ad Kirkedøren.
Under Skin af en oplyst, fornuftig Rettroenhed begyndte
Rationalismen med at kritisere Bibelen og endte
med at bortforklare Underet og udvande Christendommen.


Tilsidst indtraadte et Vendepunkt, hvor Striden
mellem Videnskaben og den orthodoxe, officielt anerkjendte
Kirkelære blev saa hæftig, at den bragte Statsmyndigheden
selv i Forlegenhed. Staten havde unegtelig
forpligtet sig til at opretholde den gamle lutherske
Rettroenhed; den havde ladet alle sine Geistlige aflægge
Ed paa de symbolske Bøger. Men paa den anden Side
% --- p. 83 ---
\setcounter{footnote}{0}%
\opage{83}maatte Staten jo ogsaa holde paa Oplysningen og anerkjende
Videnskabens Ret; her var en Korsvei. Det
viste sig, for at nævne et Exempel, da det kongelige
danske Cancelli i Aaret 1834 kom til at brydes med
„de stærke Jyder“. Den gammeldagstroende Almue
kunde nemlig ikke finde sig tilrette med de nye Lærebøger,
men sendte Børnene i Skolen med de gamle.
Det blev betragtet som Skoleforsømmelse, % defect: „Skolefcrsømmelse“ (a broken o): the letter supplied
og Forordningen
af 1814 bragtes i Anvendelse. Mulktering og
Udpantning skulde overbevise det eenfoldige Folk om
de nye Bøgers Fortrinlighed; men det frugtede ikke.
De fattige Bønder fandt sig i Udpantning og Salg af
deres Indbo, deres Faar, Køer o. s. v.; men i Brugen
af Balles Lærebog vilde de ikke finde sig. De indskrænkede
sig end ikke til simpel Vægring, de anførte
Grunde, og gave en Fortegnelse over de Steder, som
stode i Strid med den gamle Tro og havde vakt deres
Forargelse; men den høie Øvrigked kunde naturligviis % print: „Øvrigked“ (misprint for Øvrighed)
ikke indlade sig paa religiøse Beviisgrunde.

Her stod Ret imod Ret. Bønderne havde virkelig
Ret til at fordre den gamle strenge Lutherdom opretholdt;
det var jo netop denne Opretholdelse, hvortil
Statens øverste Myndighed høitidelig havde forpligtet
sig. Men Cancelliet havde ogsaa Ret; det holdt de
autoriserede Lærebøger i Hævd, fordi de vare \emph{autoriserede.}
At give efter for Opsætsighed mod det
Autoriserede var at svække Statens Myndighed. For
at opretholde sin Myndighed brød Staten sin Forpligtelse:
den statskirkelige Myndighed var kommet i
Splid med sig selv.

For at dække denne Splid med et passende Slør,
har man beraabt sig paa Noget, der kaldes det prote%
% --- p. 84 ---
\setcounter{footnote}{0}%
\opage{84}stantiske Princip, og som navnlig skal bestaae deri. at % print: „deri.“ (a stop, or a comma whose tail did not ink)
den lutherske Lære i Aandens og Frihedens Navn
ligger aaben for kritiske Prøvelser og lader sig omforme,
eftersom Skriftfortolkningerne med dertil svarende
Bearbeidelser af Dogmerne vexle, medens den katholske
Kirkelære, forstenet og krystalliseret i gamle Former,
maa værge sig mod den frie Videnskab ved at unddrage
sig dens Kritik og anvende Magtsprog istedenfor
Grunde. Protestantismens Væsen er Udvikling: den
søger Sandheden og følger med Tiden. Katholicismens
Væsen er Stillestaaen, den paastaaer at have Sandheden
og fører en fortvivlet Kamp med Aanderne i den nye
Tid. Heri er noget Sandt, men ogsaa noget Vildledende.


I Sammenligning med den Høihed, Charakteer og
Sluttethed, som udmærker den romerske Katholicisme,
tager det protestantiske Statskirkevæsen i Nutiden sig
kummerligt ud. Katholicismen veed, hvad den vil, paa
hvad Grundlag den staaer og hvilke Vaaben den skal
bruge; den har paa sin Viis endnu Myndighed, Aandsmyndighed.
Det er mere end man kan sige om de
i sørgelig Maade forkrøblede Statskirker. Paa hvilket
Grundlag ville de staae? Nogle beraabe sig paa de
symbolske Bøger, juridisk, juridisk! Men naar de symbolske
Bøgers juridiske Autoritet bogstaveligt, strengt
skal opretholdes af Statsmagten, hvad bliver det saa
til med det protestantiske Princip? Andre beraabe
sig paa Videnskaben: Kritik, Kritik og mere Kritik!
Men dersom det virkelig er den videnskabelige Kritik,
som i høieste Instans skal afgjøre, hvad Religionen er,
og hvormange Stumper af Bibelen, der endnu maae
% --- p. 85 ---
\setcounter{footnote}{0}%
\opage{85}have Lov til at være ægte: hvad har saa Statsmyndigheden
at betyde, og hvad bliver der af Kirken?

Her er en Opgave at løse. Hvorledes dens Løsning
er mulig, særlig for os efter vore Forhold, er et
Spørgsmaal, hvis fulde Besvarelse først vil kunne fremgaae
af meget udførlige Forhandlinger. Hvad jeg ved
denne Anledning har at tilbyde, er en kort og bestemt
Angivelse af de Synspunkter, som efter min Overbeviisning
ere og maae være de afgjørende. Jeg skal
bestræbe mig for at belyse disse Synspunkter i et følgende
Foredrag.\footnote{See 7de Foredrag.}
% --- p. 86 ---
\setcounter{footnote}{0}%
\lecture{Sjette Foredrag.} {(Aagaard.)}
\opage{86}„Mand døer, men aldrig døer Mands Minde“, det er
et gammel-nordisk Ord, det vidner om en høisindet
Livsbetragtning. De, som kun leve for Egennytten,
burde helst ikke nævne dette Ord, thi det truer dem
med en Dom: den lave Egennytte, den raa Egoisme,
kan opnaae meget i Verden, men der er Noget, den
aldrig har naaet og aldrig naaer: et hæderligt Eftermæle.
Hvad der gjælder om den enkelte Mand, gjælder
i videre Betydning om en heel Samtid. Vi vide jo
alle, hvor fristende det er, at see paa den nærmeste
Fordeel. Naar de, som holde paa et Bestaaende, selv
maae erkjende, at den gamle Bygnings Fald er uundgaaelig,
og dog for Fordelens, den nærmeste Fordeels
Skyld klamre sig til de møre Piller og sætte imod
hvert Fremtidsstød den fattige Trøst: „det holder vel
nok, saa længe vi leve“: da giver Sligt intet hæderligt
Eftermæle. Naar et Fremskridtsparti benytter, hvad det
anseer for et gunstigt Øieblik, paa en saa hensynsløs
og lidenskabelig Maade, at det nødvendig maa forvirre
og skade den nærmeste Fremtid, naar Løsenet bliver:
„Nu er det vor Tid, vi maae benytte vor Fordeel, de
% --- p. 87 ---
\setcounter{footnote}{0}%
\opage{87}fjernere Følger blive, som de ville!“ da giver Sligt intet
hæderligt Eftermæle. Efterslægtens Dom bliver en
Straffedom, en Fordømmelsesdom; „Mand døer, men
aldrig døer Mands Minde.“

Fremtidens Dom er mere end en Menneskedom,
det er en Gudsdom. Der gaaer igjennem Historien en
sammenhængende Række af Gudsdomme; det kommer
kun an paa, hvorvidt vi Mennesker ere i Stand til at
opdage dem, og i bvilken % print: „bvilken“ (for hvilken)
Dybde vi formaae at erkjende
dem.

Middelalderen tog Sagen ligefrem efter det gamle
Testamentes Forbillede; David fældede Goliath, det var
en Gudsdom. Naar en Retssag eller en Æressag af
stor Betydning skulde afgjøres, da kom det ikke væsenlig
an paa Vidnerne og de retlige Forhør, nei, Afgjørelsen
laa i noget Andet, f. Ex. i en Tvekamp. Var der
blevet reist krænkende Beskyldninger imod en fornem
Dames Charakteer, greb hendes Ridder til Sværdet; og
hvis han beseirede Modstanderen, var dennes Udsagn
med det Samme fældet som usandt. Det var en Folketro,
at Gud ved et Mirakel vilde lade den Retfærdige seire.
Betragte vi de egenlig saakaldte Gudsdomme, Jernbyrd,
Nadverprøve, Korsprøve, o. s. v., saa viser sig det
Charakteristiske. Man meente, at Brødet paa Alteret,
idet man fremsagde en vis Formular ved Messen, skulde
vende sig om, hvis den Mistænkte var skyldig. Eller
man indsvøbte nogle Stykker Træ, af hvilke et var
betegnet med et Kors, et andet ikke, i Uld og lagde
dem paa Alteret: dersom den Mistænkte nu tog det
Stykke med Korset paa, blev han erklæret for uskyldig,
i modsat Fald for skyldig. Tørke, Pest, Hungersnød,
bleve betragtede som umiddelbare Foranstaltninger af
% --- p. 88 ---
\setcounter{footnote}{0}%
\opage{88}Gud til Synderes Afstraffelse og Retfærdighedens Fremme.
Selv en mægtig Fjendes naturlige Død var et Fingerpeg
fra oven. Da Keiser Frederik den Anden, Hohenstauferen,
var død (1250), udbrød Pave Innocens den Fjerde i en
Henvendelse til Siciliens Geistlighed: „Himlen jubler
og Jorden fryder sig, thi Lynild og Storm, hvormed
Gud den Almægtige saa længe har truet Eders Hoveder,
ere ved denne Mands Død blevne forvandlede til en
forfriskende Zephyr og en befrugtende Dug.“ Endnu
for omtrent hundrede Aar siden vilde den spanske
Geistlighed forbyde Opførelsen af Stykker med lutter
Sang (Operaer), „fordi denne ugudelige Underholdning
var Skyld i den herskende Regnmangel.“ Man har
nu rigtignok beraabt sig paa, at slige religiøse Udlægninger
vare vel skikkede til at vække Samvittighederne
og gjøre Forestillingen om Guds allestedsnærværende
Retfærdighed levende; men det er dog indlysende, at
dersom den Almægtige ved indgribende Mirakler selv
vilde paatage sig at sørge for Retspleien, saa var det
umuligt, at saa megen Uretfærdighed kunde trives i
Verden.

Den nyere Tid er naaet til den modsatte Yderlighed.
Den guddommelige Frihed har fanget sig selv
i Lovene; de uforanderlige Love ere Baand, som den
ikke kan løse: i Sammenhængen af naturlige Aarsager,
Betingelser og Mellemaarsager hersker en ubøielig
Nødvendighed. Det erklæres reent ud, at man ikke vil
høre Noget om „en uansvarlig Gud“; og med Menneskenes
Ansvarlighed har det, for Historiens Domstol,
heller ikke stort at betyde. Betragtes de handlende
Personer enkeltviis, da maa der vel gjøres Forskjel
mellem gode og slette, ædle og uædle Charakterer:
% --- p. 89 ---
\setcounter{footnote}{0}%
\opage{89}Roes og Dadel kunne ikke ganske udelukkes. Men opfatte
vi Heelheden med de talløse Enkeltheder i Sammenhæng,
da overbevises vi om, at Historiens store
Mænd altid have handlet saaledes, som de ifølge Omstændighederne
nødvendig maatte handle. En sædelig
Maalestok for Godt og Ondt, Retfærdigt og Uretfærdigt,
kan ikke anlægges paa Verdensbegivenhederne og de
stærke Aander, den borgerlige Morals Kræmmeralen er
her for lille. Det Store vil hverken roses eller dadles,
hverken elskes eller hades; det vil, med eet Ord, beundres.
En Tiberius og en Nero have lige saa meget
Tiltrækkende for Historikeren som enten en Paulus
eller en Luther. Verdenshistorien er Verdens Dom:
Dommeren er som Skjæbnen selv, og Dommen uden
Samvittighed.

Vilde Nogen driste sig til at møde denne storagtige
Verdensbetragtning med religiøs-sædelige Vaaben i Troens
Eenfoldighed --- som David mødte Goliath, da vilde
man vistnok med et overlegent Smiil minde ham om,
at David er død, og hans Tro et for længst tilbagelagt
Standpunkt. For „en Gudsdom“ frygter man ikke;
man er sikker paa, at Miraklet vil udeblive o. s. v. ---
Hvor mærkværdigt! I Middelalderen var Almagten fri,
men hele Tidsalderen ufri; i vore Dage skal Tidsalderen
være fri, men Almagten selv ufri. Her maa stikke
Noget under.

Hver Gang et stort Maal med en stor Opgave
sættes ind i Historien, er det Almagten selv, den evig
frie Aand, der aabner Mulighedernes Dyb og stiller
Menneskene Opgaven. Opgaven udtales ikke med rene
Ord fra Himlen; det er foreliggende Kjendsgjerninger
og indre Tilskyndelser, hvorigjennem den melder sig.
% --- p. 90 ---
\setcounter{footnote}{0}%
\opage{90}En historisk Opgave kommer ikke nøgen til Verden;
den er iklædt Naturbetingelsernes, Tildragelsernes, Omstændighedernes
folderige Dragt: den har et Slør for
Ansigtet. Menneskevillien vaagner, idet den Tilslørede
viser sig. Det er Frihed, der vækker Frihed, Villie,
der taler til Villie. Efterhaanden som Opgaven nærmere
bestemmes og dens Fordringer træde ud i Virkeligheden,
afsætter den en Mangfoldighed af særskilte,
underordnede Opgaver: her er Arbeide fuldt op i alle
Retninger. Den menneskelige Villie, som vælger, virker
og arbeider, faaer i alle Dele og fra alle Sider, enten
den veed deraf eller ei, med den guddommelige Villie
at gjøre. Vor Frihed er begrændset, vor Villie paa
utallige Maader betinget af Omstændighederne. Men
Omstændighederne farves af Erkjendelsens Lys, og
Hindringerne selv give Villien Spændkraft. Forskjellige
Omstændigheder kunne lede den menneskelige Villie
ind i høist forskjellige Spor; men der er ingen Omstændighed,
som kan forhindre vor Villie fra at ville.

Den Maade, hvorpaa Millioner Villier arbeide sammen
og i Forhold til den fælles Opgave danne en Fællesvillie,
en herskende Almeenvillie, charakteriserer \emph{Tidens}
\emph{Aand.} Af Tidens Aand ere Tidens Børn alle uden
Undtagelse, den Mægtigste saa vel som den Ringeste,
afhængige, og dog i Forhold til Opgaven alle uden
Undtagelse, den Ringeste saa vel som den Mægtigste,
hvad Muligheden angaaer, uafhængige. Enhver kan
ville og vælge paa sin Maade, den Ene det Gode, den
Anden det Slette; men Styrelsen gjør dem alle ansvarlige.
Naar der saa er arbeidet med den satte Opgave
saa længe, at Indholdet er udtømt, og Kræfterne opslidte,
saa gaaer dens Tidsalder i sin Grav og giver
% --- p. 91 ---
\setcounter{footnote}{0}%
\opage{91}Plads for en ny Tid. For saa vidt vi da ved enhver
saadan Overgang, den foregaae stille og langsomt eller
pludselig som i et aandeligt Jordskjælv, agte paa Aandsvilliens
uendelige Overmagt over de endelige, men hver
især ansvarlige Menneskevillier, erkjende vi en Gudsdom
i Historien.

Fra dette Synspunkt falder der en egen Belysning
over den verdenshistoriske Kamp mellem Ret og Uret,
Godt og Ondt, der har opvakt Tvivl hos saa mange
Tænkere. Den styrende Aandsvillie tilskjøder os ikke
Besiddelse af disse eller hine jordiske Eiendele; den
vil i evig Forstand hjælpe Menneskene frem med at
ville og virke, men den tager ikke ligefrem Parti med
det Gode. Det Gode, den historisk vil, lægger den ind
i den store Opgave. Menneskene kunne gribe denne
Opgave fra hvilken Side, de ville; de, der stride under
dens Fane, have i det Hele Styrelsen med sig. Der
kan paa Opgavens Regning begaaes Forbrydelser; ja
Forbrydelserne ville endog kunne være desto større, jo
dybere Opgaven ligger, jo mere historisk Ret, den har.

Bocaccio, en italiensk Digter fra det fjortende Aarhundrede,
har i sin „Decamerone“ en Række Fortællinger,
hvoriblandt følgende. To Kjøbmænd, en
Christen og en Jøde, vare meget gode Venner. Den
christne Kjøbmand havde saa mange Prøver paa sin
jødiske Vens Ærlighed, Trofasthed, Tjenstvillighed og
retsindige Tænkemaade, at han, som selv var en god
Katholik, følte en inderlig Bedrøvelse ved Tanken om,
at en saa redelig Sjæl skulde gaae fortabt. Hvad
gjorde han? Han søgte naturligviis at omvende ham
til Christendommen. Da han imidlertid snart mærkede,
at hans egne Overtalelsesgrunde ikke vilde slaae til,
% --- p. 92 ---
\setcounter{footnote}{0}%
\opage{92}thi Jøden var sikker, naar han forsvarede sin Religion,
tænkte han ved sig selv: „jeg maa have ham hen til Præsten.“
Men det gik ikke. „Til Bispen da.“ Men det
gik heller ikke; og saa vidste den gode Mand tilsidst
ikke, hvad han skulde gjøre. En Dag træder hans
Ven ind til ham og siger: jeg vil gjøre en Reise til
Rom; jeg har Forretninger der, og ønsker at see mig
om i den berømte Stad. --- Herregud, sukkede Katholiken
ved sig selv, kommer Abraham til Rom og seer
det Sodoma og Gomorrha, saa faaer jeg ham da aldrig
omvendt; nu er mit sidste Haab ude. --- Jøden reiste
og blev borte i længere Tid. Da han kom tilbage blev
der talt om Handelsaffairer og allehaande verdslige
Ting; men det Punkt om Religionen turde den christne
Ven ikke røre, han turde end ikke nævne enten Paven
eller Cardinalerne. Hvor forbauset blev han ikke, da
Vennen pludselig taug, betragtede ham med et eget
mildt Blik og derpaa sagde: jeg vil betroe Dig Noget;
jeg har besluttet at blive Christen. --- Er det muligt?
Hvor kunde Du midt i Frækhedens, Gudsbespottelsernes
og Udsvævelsernes Pøl faae gode Tanker om Christendommen? Det er klart, svarede Jøden, den Gud, som
kan taale Sligt, maa være barmhjertigere end min.

Hvad her er sagt som en bidende Spot, kan tages
for fuldt Alvor. Det Høidemaal af Forbrydelser, hvortil
den middelalderlige Kirke omsider drev det, er
tillige et Maal for den religiøse Dybde og sædelige
Storhed, som Opgaven oprindelig indeholdt. Efter
denne Maalestok maae vi erkjende den Gudsdom, hvorunder
den sank.

Kaste vi et Blik paa det Tidsafsnit, der umiddelbart
grændser op til vort eget, og dog ligger afsluttet
% --- p. 93 ---
\setcounter{footnote}{0}%
\opage{93}for os, nemlig Napoleon den Førstes Tidsalder, da see
vi et Exempel paa, hvad en Gudsdom er. Den Opgave,
som her blev stillet, var ikke ligefrem religiøs; og dog
var Aandsvillien med, det var i Sandhed en stor. for % print: „stor.“ (a stop, or a comma whose tail did not ink)
Menneskelighedens (Humanitetens) Sag betydningsfuld
Opgave. Revolutionen brød brutalt med den historiske
Fortid. Den vilde i sin forfærdelige Vrede knække
det gamle Aag, som man knækker en udtørret Green
paa et gammelt Træ; den vilde med blodig Øxe hugge
Træet om, og mærkede ikke, hvor dybe Rødderne vare.
Ved intet Vendepunkt i Historien har Massernes Vildhed
og Mængdens Afmagt, hvor det gjælder en Styrelse,
viist sig i et klarere Lys end just da „Bjerget“ sank.
og „Dalen“ igjen begyndte at hæve sig. Trods forbausende
Opsving, uhyre Anstrengelser af beregnende
Kløgt, fanatisk Iver og et Røvermod, der har sat Alt
paa Spil, saae det dog ud, som skulde „Menneskerettighederne“
kvæles i det Umenneskelige, og Forsvarerne
af den gamle Trældom vende tilbage med tidobbelt
Ret. Da Groa, siger Mythen, sang Tryllesange
over Thor, blev Guden saa glad, at han fortalte hende,
hvorledes han havde baaret hendes Mand, Ørvandel,
over Elivaagerne. Der hørte overmenneskelige Kræfter
til for at bære det nye, i sit Inderste oprørte Samfund
hen over Revolutionens Elivaager. Men Aandsvillien
var med; der hævede sig op af Mulighedernes Dyb en
heroisk Skikkelse, en Magthaver af høiere Orden, et
verdenshistorisk Geni, Napoleon Bonaparte. Han havde
en Villie.

At en saadan Villie umulig kan være et simpelt,
naturnødvendigt Produkt af Omstændighederne, fremlyser
paa det Klareste deraf, at den overalt, til Ver%
% --- p. 94 ---
\setcounter{footnote}{0}%
\opage{94}dens Forbauselse, gjør sig til Omstændighedernes
Herre; den faaer Lykken til at gaae med og lærer de
mødende Tilfælde at bøie sig dybt for den Overlegne.
Saa havde den nye Verden da faaet sig en ny Villie,
en Centrumsvillie. Millioner og Millioner Villier, stærke
og svage, smaa og store, sættes i Bevægelse. Overraskelse,
Misundelse, Beundring, Venskab, Fjendskab,
Haab, Frygt, Had, Begeistring, Frihed, Tvang, lige
meget! alle Tanker, alle Stemninger, alle Villier maae
med: de maae alle dreie sig sóm % print: „sóm“, accented for stress
Planeter om den ene
Sol, alle bevæge sig om den ene Centrumsvillie. Her
er Frihed midt i denne Nødvendighed; her er en Verden
af Villier.

See vi paa Gangen, det tog med Opgavens Løsning,
da glemte Tidsalderens Heros desværre underveis,
at det var „Menneskerettighederne“, hvorfra han var
gaaet ud. Istedenfor at være Aandsvilliens Tjener, blev
han et Skjæbnens Redskab: Keiserørnen slog ned paa
det Hele som et Rov. Naar man betænker, hvorledes
det ene Skridt førte til det andet, hvorledes den ene
Magtudvidelse fordrede den anden, hvorledes det Overordenlige
blev dynget paa det Overordenlige, Ossa paa
Pelion, som i Gudekrigenes Dage, indtil den kolossale
Bygning omsider styrtede sammen, og Mythen om
Titanen, den til Klippen lænkede Prometheus, blev
gjort til Virkelighed paa St. Helene: da er der unegtelig
Stof nok for dem, som fordybe sig i Begivenhedernes
Sammenkjædning og beundre Nødvendigheden. Ikke
desto mindre er denne Nødvendighed fra Først til
Sidst betinget: udadtil af Tilfældenes Sammenstød, indadtil
af Villiens Valg. I Vexelvirkningen af Tilfælde
og Valg ligger den betingede Frihed. Det er Verdens%
% --- p. 95 ---
\setcounter{footnote}{0}%
\opage{95}herskeren, som af selvisk Drift uafbrudt omsætter Aandsvillien
i Skjæbnevillie. Det store Skuespil, den forviklede
Knude, de vidt straalende Glandspunkter, de
rystende Optrin og den tragiske Slutning, Alt er saaledes
gjennemvævet af Frihedens og Nødvendighedens
Traade, at Forsyn og Skjæbne begge gaae ind paa den
Stærkeres Ret og udfordre hinanden til en verdensrystende
Kraftprøve. En saadan Prøve er aldrig seet
paa den gamle Verdens Skueplads. En Alexander,
Hannibal, Scipio, Cæsar ere Exempler paa, hvad Geniet
formaaer i Kraft af ubetvingelig Driftvillie; men Skjæbnen
hersker, Friheden overskygges af Nødvendigheden.
Heltene fra den klassiske Oldtid kunde umulig have
en saadan Bevidsthed om Frihedens Væsen og Vilkaar,
som Revolutionens Søn fra det attende Aarhundrede.
Saa bliver den Mægtige ansvarlig, saa gjør Aandsvillien
sin Overmagt gjældende; saa bliver Dommen
mere end en Verdensdom, det bliver en Gudsdom.

Naar det stundom kan forarge os, at den Stærkeres
Ret lige saa fuldt skal gjælde i Historien som i Naturen,
da tør vi ikke glemme, at det paa Frihedens
Vilkaar har en egen Sammenhæng med det, vi kalde
det Stærkere. Det hedder i Religionens Sprog: bedre
at lide Uret end at gjøre Uret; men Magthaveriets
Mening er: bedre at gjøre Uret end at lide Uret.

Der kan stundom paakomme raae Magthavere en
saa uvilkaarlig Lyst til Ødelæggelse, at de hjemsøge
Lande med Ild og Sværd og afbrænde fredelige Byer,
uden at have anden Hensigt end den, at ville vise
Verden, hvor ussel en Ting det er, man kalder Lov og
Ret. Nero er ikke den eneste, der indtager en udmærket
Plads blandt de kronede Mordbrændere. En
% --- p. 96 ---
\setcounter{footnote}{0}%
\opage{96}saadan Misbrug af Magten er forsvindende i alle civiliserede
Lande. Men dermed ere „Menneskerettighederne“
endnu kun lidet hjulpne. Staternes Stillinger til hverandre
indbyrdes er og maa ifølge Sagens Natur være,
Magtstillinger. Til at haandhæve Retten udkræves der
Magt; den i og for sig sande Sætning: \emph{Ret er Magt},
vendes da om og kommer til at betyde: \emph{Magt er Ret.}
Derved fremkommer et stiltiende anerkjendt Magthaveri,
som mere end alt Andet i Verden undergraver
den sædelige Bevidsthed hos Folk og Stater. Forstaaer
sig, den Magthavende forsømmer aldrig at give
sine Fordringer, om de endog indeholde den meest
himmelskrigende Uret, et Skin af Ret; der er Paaskud
nok til Krig, naar man vil overvælde en svagere Nabo.
Saa skeer der Uret i det Store. Og naar Uretten
er skeet, hvad saa? Saa dømmer man efter Udfaldet.
Hvad der er skeet, det er skeet; den Forurettede maa
finde sig i sin Skjæbne. Lad Uretten være vitterlig
for Alle, saa vitterlig, at det eenfoldigste Menneske
kan føle paa den med Hænderne; lad den Overmægtige
selv indsee, at det virkelig er Uret, og at det kun
koster ham et Ord, at gjøre Uretten god igjen: han
kan dog ikke --- saa svag er han, den Stærke, i dette
Punkt --- han tør dog ikke gjøre det. Og hvorfor
ikke? Fordi det er hans Ære for nær; det vilde jo
see ud som Svaghed. Den Magthaver som frivillig
giver et eneste, om end nok saa lille Stykke, bort af sin
Magt, har opgivet sig selv og forraadt det hele Magthaveri.
At gjøre en Uret god, det er jo at indrømme
dens Tilværelse. „Magt er Ret“; hvorledes skulde en
Magthaver fornedre sig til at indrømme, at han nogensinde
enten havde gjort eller kunde gjøre Uret. At
% --- p. 97 ---
\setcounter{footnote}{0}%
\opage{97}gjøre en enkelt Uret god? Nei, nei! Skal man indlade
sig paa det Enkelte, saa bliver Gjælden for stor;
der staaer en Hær af Creditorer udenfor Døren. Bedre
at gjøre Uret end at lide Uret.

Og nu den offenlige Mening? Politiken er ikke
sentimental, men praktisk. Den offenlige Mening skal
jo netop fra dette Synspunkt optugtes til at forstaae,
at Spørgsmaal om Ret og Uret kun have hjemme i
Statsstyrelsens indre Anliggender, i Civil- og Criminalprocessen,
men ingenlunde vedkomme Magtstillingen.
Det er ingen Ringeagt for Ret og Retfærdighed. Tvertimod!
Jo høiere vi agte Ret og Retfærdighed, med
desto større Høiagtelse og Ærefrygt maae vi jo see op
til den Magthaver, der er stillet saa høit, at han uden
Undseelse, uden Skamrødme, kan vise Verden, at han
har Retten under sin Fod, og lære os, hvad Uret
betyder. --- Det var saa en Forklaring af den Stærkeres
Ret; vi gaae nu over til en anden.

Bedre at lide Uret end at gjøre Uret; det er
Fattigmands Trøst i de Mægtiges Øine, og desværre:
den kan ved at misforstaaes kun alt for let blive fattig
nok. \emph{At lide} Uret bliver sædvanligviis opfattet
saaledes, at Eftertrykket falder paa den bare Liden,
den daadløse Taalen, der i Sløvhed trøster sig med,
at herved er Intet at gjøre. Det er en oprørende Forestilling,
forargelig især, naar den skal gaae under
Navn af Christendom, thi det er en aabenbar Fornegtelse
af Christustanken.

Der har været En i Historien, en Eneste, som i
Aand og Sandhed vidste, hvad det var, at lide Uret.
Millioner have lidt og lide endnu den Dag idag større
eller mindre Uret; der er vel neppe Nogen iblandt os,
% --- p. 98 ---
\setcounter{footnote}{0}%
\opage{98}som hvad dette Punkt angaaer, ikke har Et eller Andet
at beklage sig over. Vi kunne have Ret i Meget, men
have dog altid Uret i Noget. Ingen tør sige, at han
kun led Uret, men aldrig gjorde Uret. Kun den ene
Lidende led Uret, ikke \emph{endskjøndt} han havde Ret i
Meget, men \emph{fordi} han havde Ret i Alt, fordi han aldrig
havde gjort Uret. De mange Lidende maae tilstaae,
at det, som især paadrager dem Lidelsen, er Savnet af
de Goder, man med Uret har berøvet dem. Hvem tør
sige, at han for Retfærdighedens Skyld frivillig har
valgt at lide? Vi elske ikke saa meget Retten selv,
som vi hade dem, der negte os vor Ret. Kunde vi
komme til vor Ret ved at slaae lidt af paa Retfærdigheden,
saa gik vi vel gjerne paa Akkord: det er Modstandernes
Umedgjørlighed, der tvinger os til at lide.
Kun den ene Lidende, han der frivillig led for Retfærdighedens
Skyld, vilde ikke gaae paa Akkord, og
derfor led han indtil det Yderste. Var nu hans Liden
en mat og daadløs Taalen, fik han maaskee intet
Stort udrettet i Verden? Ogsaa han anvendte den
Stærkeres Ret: han var i Sandhed den Stærke.

Dersom det, at lide Uret, kun var et Tegn paa
Svaghed og Afmagt, saa var Christendommen aldrig
blevet til. Christus med sine Vidner gaaer i Borgen
for, at den, der lider Uret, kan være langt stærkere
end den, der gjør Uret. Overføre vi dette paa Folkelivet,
da er det indlysende, at et Folk, som, idet det
maa lide Uret, forstaaer at forvandle Nødvendighedens
Tryk til en Frihedens Opgave, kan, om det end er et
lille Folk, faae en stor Betydning i Historien.

At elske Ret og at hade Uret er den første Betingelse
for at opnaae en bedre Samfundstilstand. Men
% --- p. 99 ---
\setcounter{footnote}{0}%
\opage{99}hos hvilket Folk skal man nu vente at finde den største
Kjærlighed til Ret og det stærkeste Had til Uret?
Dog vel hos det Folk, hvem Tilskikkelserne have lært,
hvad det betyder at lide Uret. Her er ikke Tale om
en mat og daadløs Liden. Gaaer det saa vidt med
den ørkesløse Taalen og Finden sig i Alt, at man tilsidst
forsoner sig med Uretten og tænker i sit stille
Sind: den Stærkere har dog i Grunden Ret; vi vilde
have gjort det Samme som han, hvis vi havde havt
Magten, da er Selvstændighedsgnisten udslukt, og
Folket maa døe. Der er Intet, der i den Grad vækker
den Overmægtiges Foragt som logrende, feig Eftergivenhed
hos dem, hvis Ret han har traadt under
Fødder. Vi lære det hos Romerne, vi lære det af den
franske Cæsars Historie.

Det Folk, der opdrages til at elske Ret ved at
lide Uret, faaer en god Opdragelse. Der føres i alle
Stater, med fri Forfatning, Forhandlinger nok om Love
og Rettigheder, --- et uhyre Lovgivningsarbeide! Forhandlingerne
tage undertiden en saadan Vending, at
de ere nærved at drive det civiliserede Samfund til
Fortvivlelse. Retsmaskineriet arbeider, stønner og puster;
det knager og knirker i Skruer og Hjul. Men jo mere
man stimes om Ret og det Rette, desmere skævt og
kroget kommer det til at see ud; tilsidst bliver det
umuligt at faae et eneste heelstøbt Arbeide frem. Beregnet
Imødekommen, snilde Overeenskomster, strategiske
Kneb og taktiske Manøvrer ville ikke slaae til:
det er en Alt gjennemtrængende Retfærdighedssands,
Kjærlighed til Ret og Had til Uret, som ufravigelig
fordres. En stærk og levende Sands for Ret, er tusinde
Gange mere værd end titusinde retslærde Folianter.
% --- p. 100 ---
\setcounter{footnote}{0}%
\opage{100}Hvor det Første mangler, giver det Sidste kun en
kummerlig Erstatning; men naar det Første er givet,
vil det Andet jo følge som af sig selv.

Dersom Lovgivningsvæsenet udelukkende skulde
være afhængigt af juridisk Indsigt i alle Enkeltheder,
da kunde hele Folket gjennem sine Valgte umulig sikkre
sig nogen virkelig og levende Deeltagelse i Lovgivningen
selv. Men naar det ikke er juridisk Pindehuggeri i
det Smaa, men en altid vaagen Retfærdighedssands, et
sundt og aabent Blik for Retsforholdene i det Store,
hvorpaa det væsenlig kommer an: saa kan Folket
efterhaanden opdrage sig selv til stedse virksommere
Deeltagelse og dybere Indgriben i alle Lovgivningens
Grene. Selvopdragelsen skrider da frem paa følgende
Betingelser.

1) Vælgerne maae drage Omsorg for, at den Valgte
i det Hele kan siges at repræsentere Vælgerkredsens
Stemning; thi Folket har en ubestridelig Ret til at
lægge sin Stemning i Vægtskaalen.

2) Den Valgte er under alle Forhandlinger i den
lovgivende Forsamling forpligtet til at blive sin Charakteer
tro, den Charakteer, hvormed han optraadte
for sine Vælgere; men han er ingenlunde forpligtet til
at stemme med et sluttet Parti i Et og Alt eller
klamre sig fast til en forudfattet, forstenet Mening;
han er tvertimod forpligtet, i Retfærdighedens hellige
Navn paa Ære og Samvittighed forpligtet til at søge
Oplysning, lytte til Sagkyndiges Udtalelser, mærke sig
Modstandernes Grunde, danne sig en selvstændig Overbeviisning
og stemme efter sin Overbeviisning.

3) Den Valgte staaer Vælgerne til Regnskab.
Dersom disse Vælgere een Gang for alle have sat sig
% --- p. 101 ---
\setcounter{footnote}{0}%
\opage{101}i Hovedet, at kun det er Ret, som behager dem, og
alt det, som er dem ubehageligt, Uret: da udgjøre de
kun en Mængde, ikke en organisk Deel af Folket.
Men ere de modtagelige for Grunde og ville oprigtig,
hvad Ret er, kan han jo oplyse dem om Sagens Sammenhæng
og derved give dem Midler i Hænde til at
fælde en fri og selvstændig Dom. Finde de saa, at den
Valgte ikke har truffet det, som efter deres Synsmaade
er det Gavnlige og Rette, saa have de jo fuld Ret til
at vælge en Anden. Der er i ethvert Tilfælde skeet et
Fremskridt paa den bedre Erkjendelses Vei.

Disse Sætninger med de deri liggende Sandheder
indeholde intet Nyt; men det Gamle, som bestandig
maa blive os nyt, er dette: hvorledes staaer det sig
med vor Retfærdighedssands? Vi, som udadtil maae lide
Uret, have vi indadtil lært at elske Ret? Er Fleertalsmagtens
Uretfærdighed i den Grad blevet os modbydelig,
at vi i Forhold til hverandre indbyrdes hellere,
om saa skal være, selv ville lide Uret end forurette
vore Modstandere? Slige Spørgsmaal ere ikke politiske;
det er Samvittighedsspørgsmaal. Men Politiken tør
ikke udelukke Samvittigheden; saa sandt der er et
Folkeliv, saa er der ogsaa, og maa være, en Folkesamvittighed.


Hvem er saa i høiere Forstand den Stærkere, den,
som beviser sin Magt ved at gjøre Uret, eller den, som
har lært at udvikle sin Selvstændighed ved at lide
Uret? Den Første har stillet sig selv under Skjæbnen,
den Sidste staaer under Forsynet. Skjæbnenødvendigheden
breder sig, naar Uret bestandig skal bøde paa
Uret, og det, som er erhvervet ved Vold, besiddes ved
Tvang. Paa brede Grundlag af Urettens Slangeknuder
% --- p. 102 ---
\setcounter{footnote}{0}%
\opage{102}kan der ganske vist reises store Statsbygninger, og det
er dette Storhedsskin, der uophørlig blænder og frister
de Magthavende. Men saa lidt som man kan plukke
Viindruer af Torne eller Figen af Tidsler, saa lidt kan
man høste Frugter af Ret paa Urettens Ager. Dynger
af gammel Uret afsætte Kræftknuder i Statens Legeme.
Den stærke Constitution taaler meget, og Samfundsorganismen
har mange helbredende Udveie; men Urettens
Stoffer maae bortskaffes. Man bilder sig ind, at
Uretten kan glemmes, at Tiden efterhaanden kan
standse dens Følger og sætte den ud af Virksomhed.
Det er en Magthaverindbildning. Urettens Elementer
ere som et Giftstof, langsomt tærende, uophørligt virkende.
Magthaveriet bilder sig ind, at Politikens Konst
med Tiden skal kunne forvandle det dræbende Stof til
et sundt Næringsstof. Umuligt! Kommer der ikke en
lægende Haand, som bortskjærer Urettens Kræftskade, saa
bliver Sygdommen ulægelig: det Onde gaaer over i Blodet.
Det er en Skjæbnenødvendighed, at Urettens Storværk
maa falde, thi det bærer sin Undergangs Spire i sig
selv. Uretten er det ødelæggende, Retten det opholdende
Princip i Historien. Et stort, af Vold og Uret sammenvoxet
Statslegeme kan ved sin Sammenvoxen tilsidst
blive til en eneste stor, forviklet Kræftknude. Skjæbnedyrkerne
beundre den store Kræftknude og lovprise Nødvendigheden.
Men Nødvendigheden staaer under Friheden;
Nødvendigheden er, at Aandsvillien maa fuldbyrdes;
hver Gang vi i Historien see dens Fuldbyrdelse, erkjende
vi en Gudsdom.

% --- p. 103 ---
\setcounter{footnote}{0}%
\lecture{Syvende Foredrag.} {(Kolding).}
\opage{103}Den gamle lutherske Statskirke har i sit Hjertes Oprigtighed,
vistnok fra Begyndelsen af i den allerbedste
Mening, sammensmeltet tvende til Troeslivet hørende,
men dog uligeartede Bestanddele, nemlig: Tro til
\emph{Salighed og Tro til Sædelighed.} Kirken er derved
blevet stillet i et tvetydigt Lys.

Tro til Salighed er Menighedstro, og Kirken betyder
Menigheden; Tro til Sædelighed er Folketro, og
Kirken betyder Folkekirke. Troen under begge Former
har væsenlig samme Udspring; thi Aanden er, oprindelig
seet, den samme: kun Menighedens Aand formaaer at
afgive det evige, fuldgyldige Grundlag for en sand
Folke-Aand.

Men Aanden virker i Menneskelivet paa forskjellig
Maade, med forskjellig Kraft i forskjellig Retning.
„Vort Borgerskab er i Himlen“; det er den Tro, som
Aanden vækker, stadfæster og nærer i Menigheden.
„Vort Borgerskab er paa Jorden“; det er den Tro, som
Aanden vækker, stadfæster og nærer i Folket, en Tro
paa et sædeligt, daadskraftigt Livs Betydning, en folkelig
Tro, som viser sig i Gjerninger. Menighedstroen
% --- p. 104 ---
\setcounter{footnote}{0}%
\opage{104}er det Oprindelige, Folketroen det Afledede. Forskjellen
er iøinefaldende: vi kunne fordybe os i det
Oprindelige med et taaget, usikkert Blik for det Afledede;
og Mennesker i Tusindtal kunne benytte, bruge og bearbeide
det Afledede, uden derfor i dybere Forstand
at bekymre sig om det Oprindelige.

Aanden er Livets Kilde, alt Liv er fra Aanden;
derfor skulle Tilværelsens forskjellige Grundakkorder
bringes til Samklang i en stor, fuldtonende Harmonie
af de mange Sjæle, de mange Aandens selvlydende Instrumenter.
Men hvor forskjellige ere ikke disse Grundakkorder!
Naar Banneret vaier, den røde Dug med
det hvide Kors, og Folket synger i tusindstemmigt
Chor: „Kong Christian stod ved høien Mast“, saa rører
det sig forunderligt i de Manges Bryst; det er Troen
paa, at vi ere et Folk: det er som en Susen af Folke-Aandens
Vinger, vi høre. Men Menighedens Aand
stiller Korsbanneret i et andet Lys. Menighedstroen
seer i den Korsfæstede, den Tornekronede, der har
bøiet sit Hoved, et Varsel om en anden Strid, den sidste
Strid; saa synger den en anden Sang, til dybe Orgeltoner
en Dødssang: „Hvo veed hvor nær mig er min
Ende?“ Det er forskjellige Grundakkorder; det er
modsatte Livets Poler; mon den ene af disse ligefrem
være bestemt til at udslette den anden?

Livets Grundlov er dobbelsidig: at afdøe fra Verden
og dog leve i Verden; vi skulle leve kraftig for at
afdøe, og afdøe for at leve kraftig. Det Dobbelsidige
beroer paa, at det guddommelige Ord peger i to Retninger:
mod det Salige og mod det Sædelige. „Naar
Nogen slaaer dig paa din høire Kind, vend ham ogsaa
den venstre til,“ saa talede Christus, da han lærte fra
% --- p. 105 ---
\setcounter{footnote}{0}%
\opage{105}Bjerget; men da han stod for den Ypperste Præst, og en
af Tjenerne gav ham en Kindhest, da vendte han
ikke den anden Kind til, men han sagde: „haver jeg
talet ilde, saa beviis det, men haver jeg talet ret, hvi
slaaer du mig.“ --- „Elsker Eders Fjender, velsigner
dem, som Eder forbande, gjører dem Godt, som Eder
hade, og beder for dem, som gjøre Eder Skade og
forfølge Eder;“ det er et Salighedsvilkaar, det Sind,
som kan forstaae dette, faaer at vide, at der er Noget
til, som kaldes saligt. Men hvilket Folk vil kunne
forstaae det, som om Sædelighedens Lov forbød at
føre Vaaben mod Statens og Fædrelandets Fjender?
Det Sædelige kan gaae op i det Salige; men der vil
dog Ingen paastaae, at Menneskene ikke kunne være
sædelige uden at være salige. Her er i Aand og Tro
en Forskjel mellem det Salige og det Sædelige. Denne
Forskjel maa i det jordiske Liv, under Endelighedens
Betingelser, udvikle sig til en virkelig Modsætning ---
ikke Modsigelse --- mellem Menighedstroen og Folketroen.


Til Folketroen svarer Folkekirken. Der staaer i
vor Grundlov med rene Ord, at den evangelisk-lutherske
Kirke er Folkekirke. Hermed er ingenlunde meent, at
den evangelisk-lutherske Menighedstro skal paatvinges
Folket ved Lov. Dersom dette var Meningen, saa havde
der staaet et andet Ord: der havde staaet, ikke \emph{Folkekirke},
\emph{men Statskirke.} Det var jo netop, som nylig
anført, det Charakteristiske ved den gamle Statskirke,
at den paa uklar Maade sammensmeltede Menighedstro
og Folketro. Regjeringen holdt urokkelig fast
ved den Grundsætning, at Ingen kunde være en god
Statsborger, en tro og lydig Undersaat, med mindre
% --- p. 106 ---
\setcounter{footnote}{0}%
\opage{106}han var en sand, troende Christen, og at Ingen kunde
være en sand, troende Christen, som ikke bekjendte
sig til den af Staten, i Overeensstemmelse med Skriftens
og den rettroende Kirkes Lære, befalede Tro. Det
var ikke alene for Sædelighedens, den borgerlige Enigheds
og Eensformighedens Skyld i det Udvortes, at
den anvendte saa strenge Forholdsregler imod Anderledestroende,
nei det var udtrykkelig ogsaa for Salighedens
Skyld, for de umyndige Sjæles egen Skyld.
Statsstyreren antog virkelig, at han, hvad det Kirkelige
angik, holdt Undersaatternes Samvittighed i sin
Haand og skulde paa Dommens Dag staae Gud til
Regnskab som øverstbefalende Sjælesørger.

Hvor ganske og aldeles dette statskirkelige Synspunkt
er forladt i Grundloven, fremlyser tydeligt af
§ 76. „Borgerne have Ret til at forene sig i Samfund,
for at dyrke Gud paa den Maade, der stemmer med
deres Overbeviisning, dog at Intet læres eller foretages,
som strider mod Sædeligheden og den offenlige Orden.“

Hver Tøddel og Bogstav i denne Paragraph er en
Protest imod Statskirken. Det er ikke et Mindretal
af Anderledestroende, nogle enkelte religiøse Partier,
som nu een Gang ikke ville sværge paa den lutherske
Kirkes symbolske Bøger, og som desaarsag erholde
Tilladelse til at dyrke Gud paa deres egen Viis, hvorom
her tales. Nei, her tales ikke om Tilladelse, men
om Ret, ikke om Menighedens Medlemmer, men om
Borgerne, alle Borgerne uden Undtagelse. Men naar
alle Borgere have Ret til at dyrke Gud paa den Maade,
som stemmer med deres Overbeviisning, saa har jo
uimodsigelig hele Folket Ret til at gjøre det Samme.
Derimod have Borgerne ingenlunde Ret til at lære eller
% --- p. 107 ---
\setcounter{footnote}{0}%
\opage{107}foretage Noget, som strider imod Sædeligheden. Adskillelsen
mellem Tro til Salighed, den særegne Menighedstro,
og Tro til Sædelighed, den almindelige Folketro,
er altsaa udtrykkelig anerkjendt. Grundlovens
3die Paragraph, som siger: „Den evangelisk-lutherske
Kirke er den danske Folkekirke og understøttes som
saadan af Staten“ staaer hermed i god Overeensstemmelse,
thi her tales ikke om Noget, der \emph{skal}
være, men om Noget, der er. Vil man søge en Rest
af det Statskirkelige deri, at den lutherske Kirke understøttes
af Staten: da maae vi spørge, hvorfor den
alene skal understøttes. Det kan ikke være fordi
Fleertallet i det danske Folk for Tiden ere Lutheranere,
thi saa maatte der have staaet \emph{Fleertalskirke};
men nu staaer der udtrykkelig Folkekirke. Det er en
sand Lykke for Land og Rige, at der staaer Folkekirke;
havde der staaet Fleertalskirke, saa vilde jo
ikke blot Fleertalstro, men ogsaa Fleertalsvillie, Fleertalsret
og Fleertalsmagt med det Samme været høitidelig
anerkjendt.

At den evangelisk-lutherske Kirke er den danske
Folkekirke, vil sige: den evangelisk-lutherske Menighedsaand
og den danske Folke-Aand ere væsenlig beslægtede;
Folketroen og den dertil svarende Opdragelse
bør, i Medfør af vor Folkenaturs Eiendommelighed,
overveiende gaae i luthersk Retning og gjennemsnitlig
antage et evangelisk-luthersk Præg. Hvorledes dette
nærmere bliver at forstaae, vil vise sig, idet vi betragte
det Følgende: „Folkekirkens Forfatning ordnes
ved Lov“.

Vægtige Stemmer hæve sig imod at lade den Lov,
hvorefter Folkekirkens Forfatning skal ordnes, udgaae
% --- p. 108 ---
\setcounter{footnote}{0}%
\opage{108}fra den folkelige Lovgivning. „Den bekjendelsesløse
Rigsdag“, hedder det, „kan ikke give Love for det, som særlig
vedrører den evangelisk-lutherske Troesbekjendelse“.
Vi kunne ikke slutte os til denne Synsmaade, fordi vi
ikke kunne billige dens Forudsætninger. Her forudsættes,
at det er de evangelisk-lutherske Menighedsanliggender,
Læren, Gudstjenestens Indretning, Sakramenternes
Forvaltning, Sjælesørgernes Kaldelse o. s. v.,
om hvis Ordning Grundloven taler. Men at Grundlovens
Tanker gaae i anden Retning, sees af § 76.
Her forudsættes, at Lovene skulle være Tvangslove,
ikke som i Apostlenes Dage aandelige Frihedslove.
Men at ville regjere Sjælesorgens indre Anliggender ved
Anvendelse af Tvangslove er jo ligefrem at træde i den
gamle Statskirkes Fodspor. Her forudsættes endelig
en Lovgivnings- og Regjeringsmyndighed af saa reen
kirkelig Art, at dén indtager en selvstændig Stilling
overfor den verdslige. Men er det allerede urimeligt
at ville indføre to fuldstændig sideordnede Regjeringsmyndigheder,
en kirkelig lige over for en verdslig, i
een og samme Stat, da er det mere end urimeligt, det
er aabenbar Modsigelse at ville unddrage Folket Lovgivningsret,
naar det, som skal ordnes ved Lov, netop
angaaer Folkekirken, Folkets Kirke.

Der er Andre, der mene, at Tiden nu er kommet,
da man endelig bør gjøre Alvor af at adskille Kirken
fra Staten. For saa vidt det igjen er den gamle Forestilling
om en Statskirke, som spøger i Hjernerne, have
vi Intet herimod at indvende. Religion og Stat kunne
begge være tjente med Skilsmissen: Religionen bliver
fri for en trykkende Byrde, og Staten ligesaa. Men
Talen er ikke her om en Statskirke, men om Folke%
% --- p. 109 ---
\setcounter{footnote}{0}%
\opage{109}kirken og dens eiendommelige Stilling til Menigheden.
Det gjælder om, at gjøre sig Forholdet mellem Menighedskirke
og Folkekirke saa klart, som muligt. Menighedskirken
med Sakramenterne og Salighedsordet har
sin bestemte Troesbekjendelse; Folkekirken med Sædelighedsordet
har ingen Sakramenter og bekjender sig kun
til det religiøssædelige Folkelivs almindelige Grundsætninger.
Det er hermed i god Overeensstemmelse,
at det Lys, det Liv og den Inderlighedskraft, hvormed
de store Sædelighedsgrundsætninger fyldes og fremtræde,
kunne give os en Afglands af den Aand, der rører
sig i Menigheden. Naar det nu netop er Folkekirken,
som i sig indeslutter de Aandsmidler og Betingelser,
der danne Grundlaget for Folkets Underviisning
og Opdragelse i det Hele, hvorledes vil man da kunne
forsvare at adskille Folkekirken fra Folkestaten.

Folkekirkens Forfatning kan og bør ordnes ved
Lov; her kan i dette Punkt være Tale om Tvangslove,
statsborgerlige Love. At ville anvende Tvang paa det,
som ifølge sin Natur ikke lader sig tvinge, er et Misgreb;
men at ivre imod al Tvang, hvor Tvangslove
væsenlig høre hjemme, er ogsaa et Misgreb.

Af gammel Skræk for statskirkelig Myndighed have
Adskillige taget Ordet for en næsten ubegrændset Frihed
i Folkekirken. Det skal efter deres Mening staae
i Strid med Frihedsprincipet, at Lærerne i Folkekirken
ansættes af Staten, thi derved blive de jo, siger man,
Medlemmerne paatvungne. Det skal være et Indgreb
i Familiens Rettigheder, at Børn uden Forældrenes
Samtykke nødes til at søge Skolen og gjennemgaae et
lovbefalet Underviisningskursus, der blandt Andet ogsaa
indbefatter Religion o. s. fremd. Mon det Ube%
% --- p. 110 ---
\setcounter{footnote}{0}%
\opage{110}føiede i slige Indsigelser ikke skulde blive klart, naar
man vil begynde med at spørge sig selv, hvad det er
for en Stat, hvad det er for en Kirke og hvad det er
for et Underviisningsvæsen, hvorom her tales? Den
Statsmagt, hvorfra Lovene udgaae, staaer jo ikke enten
udenfor Folket eller over Folket; det er tvertimod ifølge
Grundloven en folkelig Magt; dens Villie er en i Love
udtalt Folkevillie. Det er Folket selv, som igjennem
den folkelige Stat ad Lovgivningens Vei ordner sin
Kirkeforfatning.

Ved den lovbestemte Ordning sikkrer Folket sig,
at Kirken ikke gives til Priis, hverken for Mængdens
eller for Individernes vexlende Luner. Det er ikke en
selvmægtig Stat, som stiller sig over Folket, men det
er Folket, som formedelst Staten stiller sig over Individerne;
det er Folket, som ved sin lovgivende Myndighed
skaffer sig Sikkerhed for, at alle Individer i
den hele Befolkning deelagtiggjøres i den Underviisning
og Opdragelse, som Folketroen kræver.

I de forskjellige Samfundsklasser kan den særlige
Opdragelse og Dannelse naturligviis være meget forskjellig
og naaes ad meget forskjellige Veie; men der
er Noget, som maa være fælles for Alle, og det er det
Folkelige. Alle Børn, Piger og Drenge, de høre forøvrigt
til Familier i hvilkensomhelst Stand og Stilling,
af hvilkensomhelst Troesbekjendelse, skulle ikke alene
opdrages og undervises; men de skulle, hvad det Folkelige
angaaer, udtrykkelig opdrages og undervises paa
den Maade, som Folkevillien foreskriver: de skulle opdrages
og undervises i Folketroen. Det humane Folk
kan paa de almindelige Sædelighedsbetingelser indrømme
Borgerne Ret til „at forene sig i Samfund og
% --- p. 111 ---
\setcounter{footnote}{0}%
\opage{111}dyrke Gud paa den Maade, der stemmer med deres
Overbeviisning;“ men dersom der er et Samfund, hvis
Troesbekjendelse ikke tilsteder en saadan sædelig Opdragelse,
som Folketroen fordrer, da er det dog vel
indlysende, at Medlemmerne af et saadant Samfund
umulig kunne regnes til Folket. De, der paastaae, at
Forældre have Ret til at bestemme, hvorvidt deres
Børn skulle nyde nogen Underviisning og Opdragelse
eller ikke, paastaae med det Samme, at Forældre have
Ret til at bestemme, hvorvidt deres Børn skulle have
Adkomst til at nyde statsborgerlige Rettigheder eller
ikke. De, der indrømme Forældrene Fuldmagt til at
afgjøre, om deres Børn skulle opdrages i Folketroen
eller ikke, maae tillige indrømme dem Ret til at afgjøre,
om Børnene skulle regnes til Folket eller ikke.
Men naar man indrømmer Befolkningens Individer en
saadan Ret lige over for Folkevillien: hvor er saa
Folkets Myndighed, hvad er saa Folket?

At Folkekirken, saaledes som den for Tiden bestaaer,
er et tvetydigt Væsen, kan allertydeligst sees
paa den i mange Henseender vanskelige Stilling, hvori
Ordets Tjenere befinde sig. Christus siger: „Ingen
kan tjene to Herrer“. Men saa længe Kirkens Forfatning
endnu er, som den er, maae Ordets Tjenere
tilstaae for sig selv, at man i Virkeligheden forlanger
det Urimelige af dem, at de skulle tjene to Herrer.
De skulle være Christi Tjenere, med Myndighed af
Christi Myndighed, Aandsmyndighed; og de skulle med
det Samme være Statens Tjenere, thi de ere jo Statens
Embedsmænd og have i denne Egenskab al deres Myndighed
fra Staten. Hvor vanskelig en saadan Dobbeltjeneste
er, lære vi paa ethvert Blad af den gamle
% --- p. 112 ---
\setcounter{footnote}{0}%
\opage{112}Statskirkes Historie. Først naar Modsætningen mellem
Menighedskirken og Folkekirken træder tydelig frem,
vil det dobbelsidige Tjenesteforhold blive befriet fra al
Tvetydighed.

I Menigheden ere Ordets Tjenere ubetinget Christi
Tjenere; for denne Tjeneste staae de som Sjælesørgere
Christus og Menigheden til Regnskab: her gjælder
ingen Indblanding af statsretlig Myndighed. I Folkekirken
ere Ordets Tjenere i religiøs-sædelig Aand Folkelighedens
Tjenere og maae i denne Egenskab være den
folkelige Statsmyndighed undergivne. I begge Tjenester
ere de Arbeidere, og en Arbeider er sin Løn værd;
men Sjælesørgerens Arbeide kan ikke lønnes paa samme
Maade som Folkelærerens.

„Naar I gaae ud,“ siger Christus til sine Disciple,
„da prædiker og siger, at Himmeriges Rige er kommet
nær. Helbreder de Syge, renser de Spedalske, opvækker
de Døde, uddriver Djævle. \emph{I have annammet det}
\emph{for Intet, giver det for Intet}. I skulle ikke have
Guld, ei Sølv, ei Kobber i Eders Bælter; ei Taske til
at reise med, ei heller to Kjortle, ei heller Skoe, ei
heller Stav; thi en Arbeider er sin Føde værd.“ Her
er ikke Tale om almindelig Folkeunderviisning; her
tales udtrykkelig om Forkyndelsen af Salighedsordet,
om den sjælefrelsende Virksomhed. Alt gives for Intet,
thi det gives af Kjærlighed, og Kjærlighed gives altid
for Intet: det er en Vare, som hverken kan sælges eller
kjøbes! Derimod kan og bør det forudsættes, at hvor
Kjærlighedens Ord finder Indgang i Menneskenes Hjerter,
ville de, som annamme de aandelige Goder, være villige
til at afhjælpe den timelige Trang hos dem, som for
Intet meddele dem disse Goder, idet de erkjende, at
% --- p. 113 ---
\setcounter{footnote}{0}%
\opage{113}„en Arbeider er sin Føde værd“. Et saadant Forhold
kan ikke „ordnes ved Lov“; Sjælesørgerens Arbeide er
intet Akkordarbeide: han er forpligtet til at give Alt
for Intet, men saa er han ogsaa berettiget til at modtage
frivillige Gaver.

Anderledes med Lærernes Stilling i Folkekirken.
Folkelæreren udfører et lovbefalet Arbeide mod Vederlag
i en dertil fastsat Løn. Han er forpligtet til at
udføre Arbeidet og han er berettiget til at fordre sin
Løn. Ligesom Folkemyndigheden i det Hele vil vide
at afdæmpe Individernes Vilkaarlighed, saaledes vil den
folkelige Lovgivning ogsaa vide at sikkre det folkelige
Ords Tjenere den Løn, som en saadan Tjeneste er
værd, og stille de religiøs-sædelige Opdragere saa frit,
at de ikke paa upassende Maade blive afhængige af
Individer, som allerhelst skulle føle sig i fri Afhængighed
af det opdragende Ord.

Hvor lidt man endnu har tænkt paa at omdanne
Statskirken til Folkekirke, kan sees af de Bestemmelser,
som fra Lovgivningens Side ere trufne angaaende Oprettelsen
af Frimenigheder. Hvad danner nu Modsætning
til Frimenighederne? De ufrie Menigheder,
Menighederne i den af Statens Midler understøttede
Folkekirke. De ufrie Menigheder høre altsaa til Reglen,
de frie til Undtagelsen; en uhyggelig Undtagelse!
Er der noget Samfund i Verden, som har oprindeligt
Krav paa Frihed, da er det nok Christi Menighed;
dens Frihed er ikke tilranet, den er dyrt betalt: han
kjøbte den med sit Blod. At Folkekirkens Menigheder
ere ufrie, beviser i Kjendsgjerning, at de endnu befinde
sig i den gamle statskirkelige Tilstand.

Det nye Arrangement med Adskillelsen mellem
% --- p. 114 ---
\setcounter{footnote}{0}%
\opage{114}frie og ufrie Menigheder er især charakteristisk ved
den Maade, hvorpaa det bliver forsvaret. At en menneskelig
Øvrighed har taget sig af de guddommelige
Ting og indrettet Kirkeforfatningen saaledes, at Alt
skeer ifølge en allerhøieste Befaling, at Salighedsordet
er lovbefalet, Sakramenterne lovbefalede, de almindelige
Bønner lovbefalede o. s. v., alt dette for Ordenens
Skyld Lovbefalede --- siger man --- gjør aldeles ingen
Indgreb i den christelige Frihed. Og hvorfor ikke?
Fordi ethvert Menneske, som ikke er tilfreds med den
lovbefalede Ordning, jo har Ret til at træde ud og
søge, om han vil, sin Opbyggelse i en Frimenighed
eller et andet ikke lovbefalet Samfund. Saa trøster
man sig med, at Troen er fri, at vi i Virkeligheden
have en fuldstændig gjennemført Religionsfrihed. Der
er i denne Betragtningsmaade endda ikke saa Lidt af
en Selvskuffelse.

For saa vidt her er Tale om Individernes Valgfrihed,
er Friheden tilvisse saa udstrakt, at Ingen har
Noget at beklage sig over. Men man lægge vel Mærke
til, at det ikke er Individernes Frihed, hvorom Talen
er, men Menighedens Frihed, Menighedskirkens egen
Frihed, hvorledes staaer det sig med den? Eller skulde
det vi kalde den christelige Frihed slet ikke vedkomme
Salighedsordet selv og dets Forkyndelse, ikke Sakramenterne
og deres Forvaltning? Kan Menighedens
Aand maaskee allerbedst virke som Frihedens Aand,
naar de Midler, hvorigjennem den virker, alle ere ufrie?
Kan Menigheden høre frit, naar den veed, at Sjælesørgeren
ikke tør tale frit? Kan der udøves nogen
aandelig Tugt (Kirketugt) af Sjælesørger og Menighed
i Forening, naar begges Samvirken øieblikkelig kan blive
% --- p. 115 ---
\setcounter{footnote}{0}%
\opage{115}gjennembrudt af en udenfra indgribende Myndighed?
Disse Spørgsmaal har man slet ikke faaet besvaret,
fordi man, saa vidt vi skjønne, endnu ikke har forelagt
sig dem med klar Bevidsthed.

Frimenighederne, der skulde raade Bod paa den
totale Mangel af aandelig Frihed, hvoraf Statskirken
lider, ere, som det synes, heller ikke i Stand til at
løse Opgaven. Der kan sagtens danne sig Frimenigheder,
som paa Forhaand juridisk forpligte sig til at
lønne deres Sjælesørgere; men der er neppe nogen
Frimenighed, i hvilken Sjælesørgeren væsenlig kan siges
at gjøre Christi Befaling Fyldest: „I have annammet
det for Intet, giver det igjen for Intet.“ De, som i
Et og Alt lægge an paa at fyldestgjøre den apostoliske
Frihed, maae ogsaa i Et og Alt underkaste sig de
apostoliske Livsvilkaar.

For at det, som ikke kan undergives Tvang, og
det, som ifølge sin Natur ikke kan undvære Tvangslove,
kunne komme i det rette Forhold til hinanden,
maa Menighedskirkens og Folkekirkens, Menighedstroens
og Folketroens indbyrdes Vexelvirkning nødvendig
klares. At Menighedskirken kan bestaae for
sig uden Folkekirke, er ved sig selv indlysende: Apostelkirken
var med hele sit folkelige Præg ingen Folkekirke,
den var i Ordets allerdybeste Betydning en Menighedskirke.
Folkekirken kan fra sin Side virke for sig
uden umiddelbar Forbindelse med Menigheden; thi
ihvorvel den for at være, hvad den især skulde være
hos os, gjennemtrængt af evangelisk-luthersk Aand,
maa have et almindeligt Præg af et evangelisk-luthersk
Menighedsliv: saa er den dog uden bestemt Troesbekjendelse og uden Sakramenter. Den væsenlige Ad%
% --- p. 116 ---
\setcounter{footnote}{0}%
\opage{116}skillelse er en Betingelse for, at de kunne virke sammen
i god Forstaaelse og fredelig Eenhed. Menigheden
med Salighedsordet og Sakramenterne indeholder det
Særlige og Oprindelige; det store folkelige Opdragelsessamfund
med Sædelighedsordet, Ungdomsunderviisningen
o. s. v. hjælper sig med det Almindelige og Afledede.
Opfattes nu begges Forening og Vexelvirkning i Folke-Troens
Lys fra den folkelige Side, saa have vi, hvad
Grundloven forstaaer ved den evangelisk-lutherske
Folkekirke.

For at kunne gjøre Fyldest som Ordets Tjenere i
en saadan Folkekirke, maae de unge Mænd, som dertil
føle et indre Kald, gjennemgaae en Forskole, der
svarer til Forberedelsen i begge Retninger. De maae
som vordende Sjælesørgere gjøre sig fortrolige med den
evangelisk-lutherske Menigheds Aand og Lære; de maae
som vordende Folkelærere gjøre sig fortrolige med de
folkelige Tilstande, med Folkets historisk-sædelige Tankegang,
kort: med de Livsbetingelser, hvorunder Folke-Aanden
virker. Som Folkelærere indsættes de af Statsmyndigheden;
som Sjælesørgere vælges de af Menigheden
og erholde den apostoliske Indvielse. Den Strid,
som kunde opstaae i Anledning af det dobbelte Valg
og den dobbelte Indsættelse, er i Principet udlignet
derved, at den religiøs-sædelige Folkelærer jo meget
vel kan virke uden at være Sjælesørger, og at Valget
af en Sjælesørger er saa frit, at Menigheden, hvis den
i et ganske særeget Tilfælde skulde føle Trang dertil,
kan, uden at krænke den beskikkede Folkelærer eller
støde an imod Statsmyndigheden, vælge sig en Anden
til Sjælesørger.

Hvad Tilsynet angaaer, er det en Selvfølge, at og%
% --- p. 117 ---
\setcounter{footnote}{0}%
\opage{117}saa det maa blive dobbelsidigt. Folkelæreren staaer,
som enhver anden Embedsmand, under Tilsyn af
Regjeringsmyndigheden og hans Arbeidsdygtighed vurderes
i Forhold til den Aandskraft, hvormed han i
sædelig-folkelig Retning indvirker paa Sognets Befolkning.
Sjælesørgeren staaer under Tilsyn af Saadanne,
som have Indsigt og Aandsmyndighed til at klare de
religiøse Livsopgaver, udvikle Eiendommelighederne i
den evangelisk-lutherske Lære og paavise Udskeielserne
derfra. Hvorledes en Regjeringsmyndighed vælger sine
Tilsynsmænd, og af hvilke Virkninger Tilsynet kan være,
er en bekjendt Sag. Hvorledes Menighederne i den
apostolisk frie Tid fik deres Tilsynsmænd (Biskopper
og Ældste) valgte, og hvilke Virkninger Aandsmyndigheden
udøvede, er ogsaa bekjendt.

Der hersker hos Mange en Frygt for, at et aandeligt
Tilsyn uden retlig Ordning og juridiske Tvangsmidler
vil blive ubestemt og ufrugtbart, at Forvirring
i Lære og en tilsvarende Forstyrrelse af Troeslivet
ikke vil kunne undgaaes. Det er en utidig Frygt. Vel
maa det indrømmes, at mange af de Afvigelser, hvilke
Statskirken i sin Tid stemplede og straffede som Kjætterier,
kunne komme frem paa Ny. Men ere disse Afvigelser
farlige, kunne de jo belyses; og naar de belyses,
har jo enhver Menighed Ret til at kalde sin
Sjælesørger til Regnskab. Saadanne Kjætterstorme,
som rystede de gamle Statskirker, ville ikke være at
befrygte i den nyere Tid; thi Opfattelsen af \emph{Dogmet}
er i det Hele undergaaet en væsenlig Forandring.

Dogmet er en i bestemt Form udtalt Troessætning.
I ældre Dage, da Rettroenheden var juridisk, holdt
man sig saa nøie, saa smaaligt til hver enkelt Lære%
% --- p. 118 ---
\setcounter{footnote}{0}%
\opage{118}sætnings Form og bogstavelige Udtryk, at Enhver, der
vovede at tale frit og levende ud af et religiøst bevæget
Sind, udsatte sig for at blive beskyldt for Kjætterier,
hvorom han neppe havde drømt. Her var i
Kjætteriernes Opramsen Nok, som kunde opflamme
Lidenskaberne, ægge Spidsfindigheden og forøge Drillerierne.
For at frigjøre sig fra Orthodoxiens døde
Formelvæsen, vilde Pietismen tilegne sig Alt med Følelsen.
Men idet Følelsen, saa at sige, blødte Dogmerne
op, blev Troeslæren geléagtig, og det Hele flød hen i
taageagtig Ubestemthed. Begge Yderligheder undgaaes,
naar vi betænke, at hver enkelt Sætningsform vel er
et bestemt, men tillige eensidigt Udtryk for Troens
Indhold, at Indholdet altid er rigere end den begrændsede
Form, at Udtryksformerne maae vexle, naar Indholdets
Rigdom tilfulde skal udfoldes. Det er ikke
dogmatisk Indifferentisme, men en grundigere Opfattelse
af Troessandhedernes Væsen, hvorpaa vi ved en Afvisning
af hiin utidige Frygt kunne støtte vor Antagelse.


Hvad her er blevet fremsat skal ikke betragtes
som et praktisk Forslag til en Omdannelse af den bestaaende
Kirkeforfatning; det er slet intet Forslag, men
en Belysning, som med Udelukkelse af ethvert Hensyn til
praktiske Vanskeligheder, ligefrem er lagt an paa at opklare
de Grundsætninger, hvorefter Omdannelsen i Tidens Løb
nødvendig maa foregaae. At der i praktisk Henseende
er Vanskeligheder at overvinde, kan ikke gjælde som
Indvending mod den her udtalte Anskuelse.

Ved Bedømmelsen af denne Anskuelse maa der
udelukkende tages Hensyn til Grundsætningerne selv.
Lader det sig godtgjøre, at de opstillede Synspunkter
% --- p. 119 ---
\setcounter{footnote}{0}%
\opage{119}enten ganske eller tildeels ere forfeilede, saa ere naturligviis
Grundsætningerne falske, og det Hele falder
sammen af sig selv. Men skulde det ved nærmere
Prøvelse befindes, at Grundforholdet mellem den særlige
Menighedstro og den almindelige Folketro er blevet
opfattet og bestemt paa en saadan Maade, at de paagjældende Hovedsætninger flyde deraf og svare til Sagens
Natur, da have vi jo, om end i flygtige Omrids, et fast
Grundlag. Paa et fast, med Sandheden stemmende
Grundlag maa Kirkeforfatningen ordnes. Thi der er
dog vel Ingen, der vil paastaae, at der ved et Anliggende,
som med Rette kan kaldes Menneskelivets, og
da ogsaa Folkelivets, høieste og helligste Sandhedsanliggende,
kan gives praktiske Hensyn, som gjøre det
tilladeligt at slaae sig til Ro ved det Halvsande eller
endog ved det Usande.

% --- p. 120 ---
\setcounter{footnote}{0}%
\lecture{Ottende Foredrag.} {(Kolding).}
\opage{120}Vi have talt om Folke-Aand, Folketro og Folkelærere
i Almindelighed, men hvorledes tager alt dette sig ud,
naar det føres over i Virkeligheden: hvad er nu egenlig
Folkelærerens Opgave? hvad er det, der gjør en
Underviisning folkelig?

Til et Folks Oplysning behøves mange Slags Lærere;
lige som vi jo ogsaa have mange Slags Skoler: Almueskoler,
Borgerskoler, Folkehøiskoler, Realskoler, Landbohøiskoler,
Militærskoler, Veterinærskoler, Navigationsskoler,
politechnisk Læreanstalt, Latinskoler, Universitetet,
foruden andre Skoler saasom: Handelsakademier,
Haandværksskoler, Tegneskoler, Syngeskoler og da ogsaa
--- Dandseskoler! Ere nu alle disse Skoler folkelige,
eller er det kun nogle af dem, der ere folkelige,
medens andre ikke ere folkelige, ja enkelte maaskee
endog ufolkelige?

At Befolkningen i det Hele har Trang til mange
forskjellige Skoler, skjønnes jo allerede deraf, at slige
Skoler ere i Gang; at oprette en Skole og holde den i
Gang koster Penge og Arbeide: Sligt gjør man ikke
omsonst. For nu at overveie, hvilke og hvormange af
% --- p. 121 ---
\setcounter{footnote}{0}%
\opage{121}disse Skoler der kunne siges at være til Bedste for
Folket som Folk, tør vi ikke indskrænke os ved snævre
Hensyn til det Allertarveligste eller dog til det i Hverdagsforstand
Gavnlige og Nyttige; vi maae tillige see
paa det Ædlere, det Aandelige, ja, vi kunne, naar
Blikket er frit, saare nødig overtale os til at udelukke
det Fornøielige. Vel kunde det fra et strengt statsøkonomisk
Standpunkt synes, som om navnlig Syngeskolerne
og især Dandseskolerne maatte høre til Luxusartiklerne.
Men nei! Hvad Syngeskolerne i den nyere
Tid have udrettet, vil vistnok Enhver, som har folkelig
Sands, taknemmelig paaskjønne. Digternes Sange glæde
os dog allermeest, naar de synges, det vil da sige:
naar de synges reent og i den rette Tone; thi naar
Tonen er falsk, saa skurrer det i vore Øren. En naturlig
god Stemme vinder meget ved at gjennemgaae en
Syngeskole. Det Samme kan siges om vore naturlige
Bevægelser, især naar vi hoppe, springe og dandse.
Det er ikke nok med den bare Naturlighed, der maa
lidt Skole til. Anstaaer det et oplyst Folk at synge
oplyst, saa klæder det da ogsaa godt at dandse oplyst.

Men, kunde man indvende, Et er en Underviisning,
som i forskjelligt Maal og forskjelligt Øiemed meddeles
de forskjellige Individer i Befolkningen, et Andet er
en folkelig Underviisning. Hertil maa svares, at det,
som i enhver Underviisningsgreen og navnlig i enhver
Skole væsenlig begrunder Adskillelsen mellem Folkeligt
og Ikke-Folkeligt, ikke saa meget afhænger af Gjenstanden
som af den Maade, hvorpaa, og den Aand,
hvori Underviisningen drives. Efter mit Skjøn kan
baade det Folkelige og det Ufolkelige, for saa vidt
der ved Underviisningen kan være Tale om at danne
% --- p. 122 ---
\setcounter{footnote}{0}%
\opage{122}og udvikle en bestemt Tænkemaade, trives i alle Skoler,
i nogle mere umiddelbart, i andre mere middelbart.

Ifølge en Mening, som med Fortrin kalder sig
folkelig, og paastaaer at have det rette Syn paa, hvad
Folket og det Folkelige er, skal Sagen ikke desto mindre
forholde sig anderledes. Der skal da her i Landet
være to store Skoler, af hvilke den ene ligefrem er
folkelig, den anden, mildest talt ikke-folkelig, ja paa
Heldingen til at blive ufolkelig, nemlig: \emph{Folkehøiskolen
og Universitetet}. At nu Universitetet ikke
kan gaae ligefrem op i det Folkelige, saaledes som
Folkehøiskolen, skal villig indrømmes; men hvorfor den
videnskabelige Høiskole absolut skal være ufolkelig,
indsee vi ikke. Man paastaaer, at den har en Skygge,
en falsk Skygge, der ikke længere svarer til den Side,
hvorfra Lyset kommer, en uhyggelig sort Skygge, der
blinder Aanden og formørker Folket; denne skumle, dødningeagtige
Skygge er --- Latinen!

For ikke at overile os med at forsvare Noget,
som ikke lader sig forsvare, ville vi gaae sindig og
grundig til Værks. Vi ville ikke begynde med at undskylde
Universitetet og Latinskolerne, thi det kunde let faae
en Bismag af eensidig Forkjærlighed. Nei, vi ville begynde
omvendt; vi ville begynde med at samle, fremdrage
og udpege alt det Onde, som deels er sagt og
fremdeles kan siges om Latinen. Vi ville ikke lade os
kyse af en Skygge, men dristig gaae den latinske
„Uvætte“ tillivs. Der kan siges Meget imod Latinen,
grumme Meget.

For det Første er Latinen et dødt Sprog. At
elske et dødt Sprog fremfor et levende, især Moders%
% --- p. 123 ---
\setcounter{footnote}{0}%
\opage{123}maalet, at elske det saa høit, at man ikke blot læser
det, men gjør Stiil baade, naar man taler det, og naar
man skriver det, er at elske det Døde. Men at elske
det Døde er kun muligt for Dødfødte eller rettere
„Dødbidere“: „alle Latinere ere Dødbidere“. Hvor
uhyggeligt, naar Livet er lyst og Hjertet glad, at møde
paa sin Vei en dødfødt, tvær Latiner med et Ansigt
som en Grammatik; man er nær ved at miste sin
Barnetro. Den eneste danske Latiner, som virkelig har
Krav paa Aflad og Syndsforladelse, er Saxo, Absalons
Klerk; derfor er da ikke alene hans Historie, men ogsaa
hans Navn blevet omplantet paa Dansk: han hedder
ikke længer Saxo Grammaticus, men Saxe Sagemand;
hvilket er meget folkeligt.

For det Andet er det en bekjendt Sag, at Latinen
fordærver Folkeblodet. Det latinske Blod er oprindelig
blandet med Ulveblod. Romulus og Remus, disse
to Verdens første, store Latinere, have jo, ifølge Historiens
udtrykkelige Vidnesbyrd, diet en Ulvinde.
Adelsstammernes Marv og Safter, Gothernes og Normannernes
navnlig, vilde have holdt sig langt renere
og Tungemaalenes Eiendommelighed staaet os langt
klarere, dersom det gamle Latineri var blevet kvalt i
Folkevandringen; men Ukrudt forgaaer ikke saa let.
Det latinske Kulturblod gjorde Kulturen ulveagtig til
samme Tid som det i Hjerterne afsatte dødt Kjød.

Hos os i Middelalderen var hele Litteraturen paa
faa Undtagelser nær Latin, Psalmerne paa Latin, hele
Gudstjenesten paa Latin. Det Meste af det, Folk lærte,
var Latin, og derfor var det i Grunden godt, at man Ingenting
lærte. Kæmpeviserne, Riimkrønniken og Euphemiaviserne
holdt sig dog lykkeligviis frie for Latin: Folkets gode
% --- p. 124 ---
\setcounter{footnote}{0}%
\opage{124}Natur modstod Forgiftningen. De Geistlige vare jo
rigtignok Latinere, men Smitten bredte sig ikke vidt;
de beholdt, saa at sige, Blodet for sig selv: de maatte
jo ikke gifte sig. Dertil kom, at Smitstoffet heller
ikke var saa kraftigt, som man skulde have befrygtet,
thi det var daarlig Latin, det var Munkelatin. Farligere
saae det ud ved Reformationen, eftersom Latineriet
ogsaa blev reformeret. Da blev det bedærvede Munkemudder
kastet bort, og den ægte, gamle, kraftige Latin
fra Romervældens gyldne Dage med Flid hentet op af
den klassiske Kjælder. Zwingli og Calvin toge dybe
Drag, og Melanchthon selv var ikke fri for at have sig
en lille Ruus. Luther holdt sig ædru; men saa havde
hans Giftermaal, hvor høit vi ellers maae prise det,
nær afstedkommet en stor Ulykke. Man betænke vel,
hvad det vil sige, at Præsterne, som nu paa Grund af
den forbedrede Latin for Alvor fik deres Blod fordærvet
i den sorte Skole, giftede sig og derved gjorde
Fordærvelsen arvelig.

Den berømte Naturforsker Darwin har, hvad det
Arvelige angaaer, opdaget Noget, han kalder Atavisme.
Der fødes i en Familie et Barn med f. Ex. en besynderlig
Vorte paa Næsen. Man undrer sig og spørger:
efter hvem kan den Lille dog have den Vorte? Hos
Ingen af Forældrene findes mindste Antydning af en
saadan Forhøining; Bedsteforældrene saa vel paa fædrene
som paa mødrene Side examineres og befindes ligeledes
at have normale, vortefrie Næser. Og dog er Vorten
arvet. Der er, naar man gaaer længere tilbage, en
Oldefader eller en Oldemoder, der har havt Vorten;
Arveligheden har sprunget nogle Mellemled over: det
er Atavisme.

% --- p. 125 ---
\setcounter{footnote}{0}%
\opage{125}Dette finder nu en skjøn og rig og saare frugtbar
Anvendelse paa Latineriet. Der fødes i en Familie, lad
os antage det, foruden de sunde Børn et Par Individer
med svage Nerver, Hovedpine, Kjertelsyge, Tørhed i
Hjertet og et yderlig sløvt Blik for det Folkelige: hvorfra
stammer det? Fra Latinen! At Moderen ikke kan
Latin og Faderen ikke heller, gjør Intet til Sagen.
Er der ikke en Bedstefader, saa er der en Oldefader
eller en Tip- eller en Tip- Tip-Oldefader, der har forlæst
sig i den sorte Skole. Man skal nok, ved at gaae
tilstrækkelig langt tilbage, tilsidst støde paa Latinen.

Hvor den latinske Indgift har fundet Sted i store
Portioner, bevirker den autokratisk Hjertebanken, Gjæring
i Blodet og en voldsom Politik. I Frankrig, Italien,
Spanien, i det Hele hos de romanske Folkestammer,
kommer Blodet let i Kog, og Partierne tørne hæftig
imod hverandre. Partierne! Ja, her er rigtignok en
Besynderlighed, som ved første Øiekast maa synes at
tale imod vor Antagelse. Skulde det virkelig være det
latinske Blod, der drev saa stærkt paa i Politiken, saa maatte
jo vore Venstremænd, som dog vitterlig ikke ere Latinere,
staae langt tilbage for Høiremændene, thi Høiremændene
ere desværre stive Latinere. Men at Sagen forholder
sig omvendt, kan ikke nægtes; det er ikke Venstremændene,
men tvertimod Høiremændene, som aande
trangt og ere nær ved at seile agter ud. Ikke desto
mindre har Tingen sin Rigtighed, naar vi see nærmere
til, og ikke forglemme, at det kommer an paa Portionerne
og paa Atavismen. I ældre Dage indtog man
Latinen i store Portioner; for nærværende Tid blive
Portionerne bestandig mindre, ja ofte saa smaa, at de
nærme sig til en Lidenhed, der minder om Homøo%
% --- p. 126 ---
\setcounter{footnote}{0}%
\opage{126}pathernes usynlige Kugler. Nu seer man strax, hvordan
det Hele hænger sammen. Høiremændene, som
have faaet Indgifterne i smaa Portioner, udmærke sig
ved en vis Kjølighed, omskuende Forsigtighed, sindig
Tilbageholdenhed, og en egen Missen med Øinene, der er et
umiskjendeligt Tegn paa, at de ofte have Vanskelighed
ved at skimte det Folkelige. Venstremændene, som
ingen Portioner have faaet, ere i Reglen friskere, med
voxende Mod og et stort Øie paa Folket. At noget
Voldsomt undertiden kommer paa dem, at ogsaa de,
som Latinen udtrykker det, stundom synes at lide af
noget Menneskeligt --- humanum quid patiuntur --- er
ikke deres Skyld; det er Latineriet i deres Blod, der
river dem hen. Det er ikke deres egen Latin, nei, det
er Forfædrenes, som efter Atavismens Lov bryder
frem i dem, og det en Gang imellem med saa stærke
Stød, at vi ogsaa ved dette Syn maae sukke og sige:
„al vor Ulykke kommer fra Latinen.“

Det er en stor Bekræftelse paa en Hypotheses
Gyldighed, at det, som ved første Øiekast synes at
tale imod den, ved nærmere Overveielse befindes at
tale for den. Vi kunne da trøstig gaae videre og
forelægge os selv det alvorlige Spørgsmaal, om ikke
den megen Forvirring i Pengevæsenet, de falske Speculationer
og gruelige Falliter, hvorom man læser i
Aviserne, i det mindste for en stor Deel, skulle befindes
at have deres første Udspring og sidste Rod i
Latinen. Horatz siger: \emph{virtus post nummos}, „Pengene
først, saa kan Dyden komme bag efter“; det er
klassisk Latin. Det være som det vil med Falliterne;
saa meget er vist, at det gaaer tilbage med Forlovel%
% --- p. 127 ---
\setcounter{footnote}{0}%
\opage{127}serne, og at dette er Latinens Skyld, kan bevises saa
klart som at To og To er Fire.

Det er en vitterlig Sag, en ved mange og paalidelige
Erfaringer stadfæstet Kjendsgjerning, at den døde
Latin gjør Hjertet koldt; men naar Hjertet er koldt,
hvorledes skal man da kunne elske? Hagbarth og
Signe, Axel og Valborg, Other Ebbesen og Sigrid med
Sløret, slige Forlovelser findes ikke mere i Norden ---
det skulde da være som Undtagelser, ifølge en ganske
særegen oldnordisk Atavisme --- derimod findes Fornuftpartier,
Pengepartier og, hvad der næsten kan
bringe Haarene til at reise sig paa Ens Hoved, Avisfrierier
i Mængde. Hvorfra det stammer, er let at opdage;
det stammer især fra den slemme, mere og mere
overhaandtagende Vane at frie skriftlig; men denne
Vane eller rettere Uvane stammer igjen fra Latinen.
Hvad er det Latinen ringeagter? Det levende Ord.
Hvad er det Latinen høiagter? At gjøre Stiil. Men
at frie skriftlig er jo at gjøre Stiil. Ved at sætte de
to Ting sammen vil man kunne forstaae, at et nordisk
Menneske med et latinsk Hjerte maa lade det Øieblik
i Livet, hvor det levende Ord og i dobbelt Forstand
den varme mundtlige Meddelelse er og bliver det Høieste,
gaae sig koldt, ligegyldigt forbi.

Et koldsindigt, ukjærligt Hjerte kan ikke elske
Folket, et latinsk Hjerte kan umulig nogensinde enten
være eller blive et ægte dansk folkeligt Hjerte; det er
i en Hovedsum alt det Onde, jeg for Øieblikket veed
at sige om Latinen.

Altsaa tilbage til Spørgsmaalet: er Universitetet
ufolkeligt, fordi Oldtidens Tanker og Oldtidens Sprog
udgjøre en væsenlig Bestanddeel af det Grundlag, hvor%
% --- p. 128 ---
\setcounter{footnote}{0}%
\opage{128}paa det er bygget? Kunne omfattende Undersøgelser
og indtrængende Sagkundskab antages at svække den
sunde Menneskeforstand og omtaage den naturlige
Sands? Er det maaskee en Selvfølge, at vedholdende
Øvelse i sammenhængende Tænkning, mangesidig Opfattelse
af Tilværelsens Vilkaar, grundig Forstaaelse af
Nutidens Forbindelse med Fortiden, omhyggelig Skjelnen
mellem det, som vides, og det, som ikke vides, o. s. v.
skulde være til Hinder for en rigtig Bedømmelse af
det, der foregaaer for Alles Øine og viser sig for det
umiddelbare Syn? Den Folkelærer, som enten ligefrem
eller ad Omveie udbreder slige falske Forestillinger
iblandt Folket, fører Folket vild. At der desværre kan
gives Videnskabsdyrkere eller rettere lærde Sysselmænd,
som sysle med det Forbigangne i den Grad, at de tilsidst
forglemme det Nærværende, er unegteligt; men
det er Folkelærerens Pligt at gjøre opmærksom paa,
at Skylden er hos Personerne, ikke hos Videnskaben.
De, som paastaae, at den strenge Videnskab ikke har
Noget med det virkelige Liv at gjøre, udmærke sig
just ikke ved klart Syn paa det Virkelige. Det er
jo netop den strenge Videnskab, der har lært os at
beherske Naturen ved Kjendskab til Naturens Kræfter
og Love. Ikke alene Dampskibene, Jernveiene og
Telegraphtraadene, men endog vore Agerdyrknings- og
Arbeidsredskaber, Stoffernes Blanding til Bedste for
Jorden, Planternes og Dyrenes Behandling o. s. v.
minde os jo hver Dag og i hver Time paa Dagen ved
hvert Skridt vi gjøre, om Videnskabens vidunderlige
Indgriben i vor naturlige Tilværelse. Vil Nogen yttre
sig med den uforgribelige Mening, at det kun er det
Nyttige i Videnskaben, ikke Videnskaben selv, et Folk
% --- p. 129 ---
\setcounter{footnote}{0}%
\opage{129}maa sætte Priis paa, da er det Folkelærerens Pligt at
oplyse ham om hans Vildfarelse. Det er dog vel lige
saa umuligt at drage Nytte af Videnskaben, naar man
unddrager dens Dyrkere de fornødne Midler, som det
er umuligt at have Nytte af et Huusdyr, naar man
forsømmer dets Røgt og Pleie. Udstykning, Kniberi
og Kvaklen med Adskillelse mellem det Theoretiske og
det Praktiske i Videnskaben har ingensteds hjemme:
enhver Videnskab er, ifølge sin Natur, theoretisk.

Derimod gives der en Forskjel, en væsenlig og
gjennemgribende Forskjel mellem sande og falske, om
man vil, forlorne Videnskaber. Ligesom man i det
virkelige Liv har adskillige Ting, der kaldes forlorne
--- forlorne Krøller, forlorne Tænder, forlorne Hofter
o. dsl. --- saaledes findes der endnu i Videnskabernes
Kreds Adskilligt, der maa kaldes Forlorent. At
adspørge de Døde, (Spiritismen), at gjøre Mirakler
eller, hvad der omtrent er det Samme, at forklare
Mirakler, udgrunde de guddommelige Mysterier og begribe
den hellige Treenighed, Sligt hører ind under de
forlorne Videnskaber. Det er Folkelærerens Pligt at
oplyse Folket om, at der gives noget Forlorent, som
kalder sig Theologie, en Humbugsvidenskab, som svækker
og undergraver den oprindelige Sands for det i Religionen
Oprindelige og derved skader det Folkelige.

Den inderlige Sammenhæng mellem det Videnskabelige
og det Folkelige kommer især tilsyne, naar
Spørgsmaalet er om vort eget Folks Andeel i Civilisationens
og Videnskabens Fremme. Vi have i denne
Henseende ikke alene holdt Skridt med Tiden; men vi
have et endog forholdsviis betydeligt Antal Mænd,
som staae paa øverste Trin og blive af alle oplyste
% --- p. 130 ---
\setcounter{footnote}{0}%
\opage{130}Nationer ærede som Videnskabens ypperste Repræsentanter:
en Niels Steensen, en Tycho Brahe, en Ole
Rømer, en Hans Christian Ørsted, det er Navne, som
have Noget at betyde i Videnskaben. Der udkommer
neppe et større Værk over Videnskabernes Historie, være
sig i Europa eller i Amerika, uden at disse Navne
findes deri og lyse som Stjerner af første Rang;
det er Navne, som kaste Glands over Folk og Fædreland.
Sligt bør fremhæves, ikke for at give den nationale
Forfængelighed Næring, men for at vække den
folkelige Selvfølelse og indskærpe en dertil svarende
Forpligtelse.

I Oplysningens Tidsalder er det af Vigtighed, at
et Folk holder Synskredsen aaben til alle Sider, ikke
lukker sig inde med sine egne Forestillinger og forgaber
sig i det umiddelbart Hjemmegjorte; men det er
tilvisse ogsaa af Vigtighed, at et Folk, især et lille
Folk, ikke fortaber sig i tomme Efterligninger af det
Udenlandske, at det har Tillid til sig selv, lærer at
stole paa sine aandelige Førere og tilegne sig det
Fremmede saaledes, at det derved erholder Betingelser
for en rigere, sikkrere og alsidigere Udvikling af sit
Eget. Det er dets egne fremragende Repræsentanter,
Digtere, Konstnere, Tænkere, Talere, Naturforskere,
Sprogforskere, Historikere o, s. v., % print: „o, s. v.“
paa hvilke et Folk
skal kjende sin Begavelse. Det er Folkelærerens Pligt
at gjøre opmærksom paa, at disse Begavede med deres
Intelligens ikke staae udenfor Folket, men høre i allerstrængeste
Forstand Folket til, ere Folkets egne Organer.
Deres Viden er Folkets Viden, deres Intelligens
er Folkets Intelligens; det er med deres Øine, at Folket
seer, hvad det ellers ikke kunde see; det er i
% --- p. 131 ---
\setcounter{footnote}{0}%
\opage{131}deres Tanker, at det gjenkjender sine egne Tanker og
tilfulde kan arbeide med i Tankernes Verden. Ved
dem er Videnskaben Folkets Videnskab, dens klare
Øie, Folkets Øie. Derfor skal man lære den menige
Mand, at de, der nedsætte Intelligensen, haane det
intelligente Folk. Det skal siges høit og frit, at Enhver,
som angriber den høieste videnskabelige Læreanstalt
i Landet, er, om han ellers veed, hvad han
gjør, i Færd med at slaae et Øie ud paa Folket og
stjæle en Juveel af Folkelivets Krone.

Og Et endnu. Der er hos det menige Folk en
Følelse af, at Videnskaben er verdslig, at den nærer
Vantroen og leder til Ringeagt for Religionen. Nogle
ville heri see et godt Tegn, et Varsel om Befrielse fra
gamle trykkende Fordomme, en Forjættelse om, at
Menneske-Aanden omsider bliver myndig og kommer
selvstændig til sig selv. Men de Fleste ere dog enige
om, at det voxende, mere og mere om sig gribende „Fritænkeri“
er farligt, at man desaarsag ikke kan Andet
end nære en vis Mistro til Videnskaben. Det er Folkelærerens
Pligt at lyse op i dette Virvar og gjøre det
paa en saadan Maade, at det menige Folk kan følge
med og danne sig en selvstændig Dom. Vistnok vilde
det være upraktisk og usundt at belæmre Folket med
videnskabelige Forhandlinger og vikle det ind i lærde
Stridigheder, men det behøves heller ikke. Der er
Noget, som ligefrem kan gjøres indlysende for den
sunde Sands, og det er tilstrækkeligt.

Religion og Videnskab have hver sit Øiemed og
sine særlige Opgaver. Religionen henvender sig til
Menneskets Samvittighed; den vil vække os til \emph{Selverkjendelse},
udvikle og klare vor \emph{Gudserkjendelse}.
% --- p. 132 ---
\setcounter{footnote}{0}%
\opage{132}Videnskaben henvender sig til Menneskets Bevidsthed
om de synlige Ting, om deres Væsen og Sammenhæng;
den vil vække os til \emph{almindelig Selvtænkning},
udvikle og klare vor \emph{Verdenserkjendelse}. Forsøger
man nu i Religionens Navn at ville foregribe Verdenserkjendelsen
og hovmesterere Videnskaben, saa gjør
man Overgreb; forsøger man i Videnskabens Navn at
ville foregribe Gudserkjendelsen og hovmesterere Religionen,
saa gjør man ogsaa Overgreb. Her maa gjøres
Forskjel mellem Religionen selv og en falsk Anvendelse
af det Religiøse, og ligesaa imellem Videnskaben selv
og en falsk Anvendelse af det Videnskabelige. Det
kommer herved fornemmelig an paa en tilstrækkelig
Adskillelse mellem den religiøse Sandhed og den
religiøse Forestillingsmaade, mellem den videnskabelige
Sandhed og den videnskabelige Forestillingsmaade.

Ifølge Bibelens Lære er Verden skabt i sex Dage.
Den Religiøse, som ikke kjender Videnskaben, holder
sig i al Troskyldighed til den bibelske Forestillingsmaade,
og gjør vel deri; thi hvert Træk i den storstilede
Fremstilling anskueliggjør den skabende Almagt,
den skabende Viisdom, den skabende Villies Øiemed.
Den videnskabelige Forsker tager den virkelige Natur,
som den er; han seer bort fra den bibelske Forestillingsmaade,
og gjør vel deri; thi han maa jo see
med uhildet Blik paa Tingene. Saa kommer han da
til det Resultat, at Alt udvikler sig naturligt og desaarsag
langsomt; at der til Jordens Afkjøling maa
være hengaaet Millioner Aar, inden der kunde fremkomme
levende Væsener o. s. v. Forvexler man nu
Forestillingsmaaden med Sandheden, er Forvirringen
grændseløs. De Religiøse holde fast paa Forestillingen
% --- p. 133 ---
\setcounter{footnote}{0}%
\opage{133}om de sex Dage, fordi de holde fast paa Forestillingen
om Verdens Skabelse. De Videnskabelige forkaste
Forestillingen om Verdens Skabelse, fordi de forkaste
Forestillingen om de sex Dage. Fra begge Sider glemmer
man Sandheden. Den videnskabelige Sandhed, at en
Naturudvikling kan udkræve Millioner Aar, staaer ikke
i Strid med den religiøse Sandhed, at den virkelige,
sig udviklende Natur har sin Oprindelse fra den skabende
Aand.

Der er en Række Grændsebestemmelser angaaende
Religion og Videnskab, som for Folketroens, Folkelivets
og Folke-Aandens Skyld nødvendig maae staae
klart for hele Folkets Bevidsthed. Det er ikke de forviklede
Undersøgelser, men de tydelig fremspringende
Hovedpunkter, hvorpaa det kommer an. Ogsaa fra
denne Side vil det kunne sees, hvilken Betydning en
Læreanstalt som Universitetet kan have for Folkelivet
i det Hele. Den Oplysning, hvorpaa det menige Folk
har Krav, og hvortil det høiligen trænger, dersom det
virkelig skal komme med og gjøre sin Stemme gjældende
i Samfundets høieste Anliggender, er af en ganske
egen Art. En eenfoldig Katechismusunderviisning i
Forbindelse med nogle tarvelige Kundskaber og en
Deel praktiske Klogskabsregler strækker ikke til; Folkets
Øie maa lidt efter lidt vænne sig til at see de
virkelige Verdensforhold i Nutidens Lys; Folket maa,
med andre Ord, lære Kulturens Vilkaar at kjende. For
at dette kan skee, er det ikke nok, at en Folkelærer
sætter sig ind i Folkets Tankegang og taler fatteligt
og forstaaeligt. Hans Tale være nok saa jævn, fattelig
og eenfoldig: den maa dog i det Hele som i
det Enkelte have Farve og Skjær af det Sandhedslys,
% --- p. 134 ---
\setcounter{footnote}{0}%
\opage{134}som udgaaer fra de klare Tankers Rige; Folkelæreren
maa være en videnskabelig alsidig dannet Mand, han
maa selv være et levende Bindeled imellem det Folkelige
og det Videnskabelige.

Vende vi os nu til Folkehøiskolerne, jeg mener de
grundtvigske, da have vi her et storstilet, ja enestaaende
Exempel paa, hvor oprindeligt, klart og dybt indtrængende
et aandeligt Syn paa Folkets Tarv og det Folkelige
selv kan være. Den „Livsoplysning“, for hvilken
Grundtvig stred og som han omsider fik sat i Kraft,
er i alle Hovedtræk, saa vidt vi skjønne, saa træffende
anlagt og eiendommelig ordnet, at hermed i Sandhed
kan siges at være gjort en ny Begyndelse. Øret aabnes,
Sindet vækkes; Mythologien, Poesien, Sagaen og Historien
oplukke deres Gjemmer og tilbyde deres rige Skatte,
ikke for at belæsse Hukommelsen og eensidig øve Forstanden,
men for at oplade den indre Sands, belyse
Menneskelivet og afgive Maalestok for hvad der er
smaat og hvad der er stort i Tilværelsen. Det er her
ikke Opgaven enten at lovprise eller forsvare de grundtvigske
Høiskoler; de forsvare sig selv. Hvor megen
Famlen og Uklarhed der end kan hæfte ved en saadan
første, ny Begyndelse, saa er det dog en Kjendsgjerning,
at disse Skoler i Landets forskjellige Egne --- ikke
mindst i denne Egn --- have udøvet en saa vækkende,
saa dannende og oplivende Indflydelse, at ingen andre
Underviisningsanstalter kunne opvise noget Tilsvarende.

Men just fordi vi sætte et saa stort Haab til
Folkehøiskolernes stadige Fremgang og voxende Indflydelse,
just fordi vi ansee det for ubodelig Skade, om
disse Skoler skulde forfeile deres Opgave og skuffe
Forhaabningerne, just derfor er det vor Pligt at gjøre
% --- p. 135 ---
\setcounter{footnote}{0}%
\opage{135}opmærksom paa, at ogsaa disse Skoler med al deres
Folkelighed kunne lede Folket vild og saaledes blive
ufolkelige.

At forelske sig i vor hedenske Oldtid og vor katholskromantiske
Middelalder paa en saadan Maade, at man
derved forsømmer at klare sig Forstaaelsen af Nutidens
Vilkaar, er usundt. Et daadskraftigt Folk kan ikke
leve af bare Minder, om de end ere nok saa herlige,
nok saa store; endnu mindre kan et oplyst Folk ville
gjøre Forsøg paa at skrue Nutiden tilbage til Oldtiden.
Det har for det menige Folk noget Tiltalende at see
tilbage paa den store Bondetid. Thorgny Lagmands
Tale til Oluf Skotkonning bliver undertiden dragen
frem til et Exempel paa, hvorledes et frit Folk har at
forholde sig lige over for en fri Konge. Dette Billede,
saa charakteristisk det end er for hiin Tid, er blændende
og vildledende, naar det stilles op som et Forbillede
for den nærværende Tid. Man lægge vel Mærke
til, at Thorgny's dristige Ord netop ere betegnende
for den udfordrende Selvtillid, hvormed Bondefolket
meente at kunne stole paa de gamle Kræfter og de
gamle Rettigheder, uden Anelse om, at de to Kulturmagter,
Kirken og Staten, som vare bestemte til at
føre Historien videre, ogsaa vare bestemte til at kue
den trodsige Almue, ja trykke den ned i en næsten
ottehundredeaarig Trældom. Thorgny Lagmand seer
ikke, hvad det er for to forunderlig mørke Uveirsskyer,
der alt drage op og trække sig sammen over hans
Hoved.

Vi prise Valdemarernes Tid; vi see op til Biskop
Absalon og ære ham som en Rigets og Folkets Mand.
Men det er jo vitterligt, at Valdemarernes Tid var
% --- p. 136 ---
\setcounter{footnote}{0}%
\opage{136}ingen Bondetid: Biskop Absalon var ikke Bøndernes
Ven; og dog var han, efter Datidens Vilkaar, i Ordets
sande Betydning, Folkets Ven. Hvor lidt „den menige
Mand“ paa Absalons Tid var skikket til at føre det
store Ord om Krig og Fred, skjønnes blandt Andet af
de Betænkeligheder og Omsvøb, der forhindrede Kong
Valdemar fra at udføre sit første Forsæt: strax at angribe
Venderne. Absalons Svar i denne Anledning er
bekjendt; han var i sit stille Sind forbittret over, at
Kongen havde været saa uforstandig at raadføre sig
med den menige Mand om Sligt. Saxo er ganske af
Absalons Mening; „thi,“ siger han, „Krig skal besluttes
med faa vise Mænds Raad og fuldføres med den menige
Mands Hjælp og Styrke.“

Valdemarernes Historie lærer os, at Adelsvælden
lige saa lidt som Hierarchiet har været et Produkt af
Tilfældet eller af Storætternes vilkaarlige Herskesyge.
Det menige Folks Undertrykkelse og stadige Synken
var dybt begrundet i Forholdenes Natur, en sørgelig
men nødvendig Følge af, at den uoplyste og desaarsag
umyndige Almue ikke formaaede at fatte de høie
Kulturmagter, sikkre sin Selvstændighed ved at benytte
de rette Midler og følge Civilisationens Gang.

At give det Udseende af, at Menigmandssyn og
Menigmandsforstand, overladt til sig selv, er i Stand
til at magte Kulturforholdene, organisere Samfundet
efter sit eget Hoved og selv udelukkende staae for
Styret, er at bedrage den menige Mand og berede ham
bittre Skuffelser. Skal det menige Folk raade med og
give sin Stemme Vægt, da maa det efterhaanden lære
at gjennemskue den Modstand, som det hidtil har havt
i sit Ubekjendtskab med de Aandsmagter, hvis Ind%
% --- p. 137 ---
\setcounter{footnote}{0}%
\opage{137}flydelse det kun altfor let har meent at kunne oversee.
Vi ønske den store Bondetid tilbage, men, vel at
mærke, ikke paa Oldtidens, nei, paa Nutidens Vilkaar.

Der er et Stykke Fædrelandshistorie, som slet ikke
seer folkeligt ud; det gaaer fra Enevoldsmagtens Indførelse
til Stavnsbaandets Løsning. Man kan ikke fortælle
det som man fortæller en gammel Saga eller
gjengiver et Optrin fra Saxo eller fra Snorre; og dog
indeholder det Afsnit, hvis rette Belysning er af høieste
Vigtighed, naar det menige Folk skal faae Øiet op for
Tilblivelsen af de Samfundstilstande, hvori vi leve.
Stemninger, Billedsprog og livlige Skildringer ere af
ubetalelig Værdi for folkelig Vækkelse; men det er i
Virkeligheden ikke nok med den bare Vækkelse. Nutidsforholdene
ere indviklede; der hører Indsigt og
Sagkundskab til for at oprede Vævet og løse Knuderne.
At gaae ind i Enkeltheder og forsøge en virkelig Løsning,
kan naturligviis ikke være Folkelærerens Opgave.
Men det kan og maa være hans Opgave at gjøre det
menige Folk opmærksomt paa, at Knuderne nu en
Gang ere der, og at de lige saa lidt lade sig opløse i et
Billedsprog som de lade sig overhugge med en Buløxe.

Skal det lykkes i vort Land, hvad man overalt
tilstræber, men endnu ingensteds har naaet, at gjennemføre
en virkelig Folkeopdragelse, Folkeoplysning og
Folkedannelse, da maa der være et Omraade, hvor
Universitetet og Folkehøiskolen trolig arbeide sammen.
En saadan Samarbeiden vil ikke udslette Forskjelligheden
af de to Skolers Charakteer; men den vil --- i
folkelig Forstand --- udslette det gamle uhyggelige
Grændseskjel mellem Lærde og Ulærde, Menige og
Ikke-Menige, Intelligente og Ikke-Intelligente, som saarer
% --- p. 138 ---
\setcounter{footnote}{0}%
\opage{138}og adsplitter, den vil bringe Lyset fra oven og Kræfterne
fra neden i en saadan Vexelvirkning, at de aandelige
Næringssafter komme i Omløb, at vi med alle en
alsidig Oplysnings Farvetoner, med alle Afskygninger i
alle Overgange fra det Dunkle til det Klare, fra det
Eenfoldige til det i Kundskab og Indsigt Mangfoldige
udgjøre et i Sandheds og Virkeligheds Erkjendelse
enigt, oplyst Folk.

At Vexelvirkningen vil være frugtbar ikke alene
for Folkehøiskolen men tilvisse ogsaa for Universitetet,
er en Overbeviisning, som paa iøinefaldende Maade har
stadfæstet sig for mig ved den aandelige Opvakthed og
uhildede Imødekommen, jeg har fundet i denne Kreds.
Jeg takker den ærede Forsamling, Tilhørere og Tilhørerinder,
for velvillig Opmærksomhed. Jeg skal ikke
glemme, at dersom det her Udtalte er, hvad jeg haaber,
blevet forstaaet, som det er blevet hørt, da skyldes
det vel for en god Deel de Folkehøiskoler, som i en
Række af Aar have virket i denne Egn. Den Tak, jeg
bringer Dem, er vel begrundet; jeg er taknemmelig for
det Indtryk, jeg her har modtaget af aabent Øre, af
opvakt Sands og fremfor Alt af et frisk, fremblomstrende
Folkeliv.

% --- p. 139 ---
\setcounter{footnote}{0}%
\lecture{Niende Foredrag.} {(Fredericia).}
\opage{139}Vi ere et lille Folk. „Ja desværre,“ svarer man med
et Suk, „vi ere for smaa til at gjøre vor Ret gjældende.
Havde vi en Armee paa 800,000 Mand og en Flaade,
som svarede dertil, hvor vilde det lyse op i Folkerettens
gamle og nye Papirer. Alene den moralske
Virkning paa det diplomatiske Corps skulde være opbyggelig
at see paa. Den Rimelighed og Forekommenhed,
den redelige Villie og Hjertens Oprigtighed, hvormed
man da fra alle Sider --- for Fredens og Enighedens
og den gode Forstaaelses Skyld --- vilde tilbyde
sine bona officia, maatte jo gjøre et rørende Indtryk
paa enhver Menneskeven; rastløs Iver med at faae
sønderrevne Traktater helede, besvorne Løfter opfyldte,
forglemte Erstatninger ydede o. s. v. skulde sikkerlig
give det nu tvivlende Europa et lige saa slaaende som
overraskende Beviis paa, at der endnu er Ærlighed,
Ordholdenhed og Redelighed til i den politiske Verden.
Men, ak desværre, vi ere et lille Folk! At vi nu just
skulle mangle det Kvantum Jern og Bly, der er saa
nødvendigt for at sætte Retfærdigheden i Bevægelse,
er det dog ikke, som hvilede der en Forbandelse over os.“
% --- p. 140 ---
\setcounter{footnote}{0}%
\opage{140}Denne Synsmaade er ikke alene trykkende, den er
tillige eensidig og falsk. Er det allerede betænkeligt,
naar en ellers vel begavet Person er blevet saa vred
paa Skjæbnen, at han ikke længer gider bruge de
Kræfter, han har, da bliver Sagen endnu mere betænkelig,
dersom et vel begavet Folk kunde falde paa
at give sig hen til en lignende Fortvivlelse. Der kan
hvile tunge Byrder paa de smaa Folk; men derfor tør
det dog ikke oversees, at de, der hvile paa de store,
kunne forholdsviis være lige saa tunge. Lad os i denne
Sammenhæng betragte vore Vilkaar fra den lyse Side
og betænke, hvor megen Velsignelse der kan følge med
det: at være et lille Folk.

Et lille Folk, hvor modigt og tappert det end er,
kan, navnlig i vore Dage, ikke vente at høste saadanne
Laurbær i Krigen, som et stort Folk; men kan det
derfor Intet gjøre for at forsvare sin Selvstændighed?
Det franske Folk er i Ordets ædlere Betydning tilvisse
et stort Folk, et Folk, der har glimret og i Fremtiden
sikkerlig ogsaa vil glimre paa Krigens Skueplads; men
har det desuagtet ikke maattet prøve, hvad Nederlag
og Ydmygelse er, fuldt saa vel som vi? Der er for
det Selviske i Menneskenaturen unegtelig noget Tilfredsstillende
i at sætte sig ud over Moralens Fordomme
og gjøre store Erobringer; der er, naar Lidenskaberne
hidses, noget endnu mere Tilfredsstillende i
at lade Gjengjældelsens Sværd lyne og tage Hævn over
lidte Nederlag. Slige Tilfredsstillelser ere alene de
store Folk forbeholdne. Maatte vi nu tilstaae for os
selv, at de uhyre Menneskeslagterier, hvormed Herskesygen
og Hævnen gjensidig bespise hinanden, høre til
det Ædleste og Bedste, for hvilket det er værd at leve,
% --- p. 141 ---
\setcounter{footnote}{0}%
\opage{141}saa havde vi unegtelig Grund til Klage; men dersom
Menneskelighedens Sag ogsaa kan fremmes paa fredelig
Maade, da har et lille Folk i denne Henseende jo Intet
at beklage sig over. Et stort Folk kan offre Menneskeliv
i Tusindviis uden just derfor at sætte sin Tilværelse
paa Spil, det har jo Masser nok at gjøre til
Kanonføde; men et lille Folk maa sætte Livet ind, naar
dets Selvstændighed trues. For et lille Folk faaer altsaa
Forsvarskrigen en ganske anden sædelig, personlig
Betydning end enten Erobringskrigen eller Hævnkrigen
faaer for et stort. I dette Sædelige, dette Personlige
er der en Velsignelse.

En tapper, men ellers fredelig Riddersmand spørger
ikke om, hvor stor Sandsynligheden er for at seire,
naar hans Ære og Person angribes; han vover Livet,
thi Ære og Liv ere hos ham uadskillelige. Han gnider
ikke paa en Skilling, naar Spørgsmaalet er om et
Sværd og en tilforladelig Rustning, thi Sværd og Rustning
udgjøre en væsenlig Deel af hans virkelige Personlighed.
Han anseer ikke den Tid og Flid, han har
maattet anvende paa sine Vaabenøvelser, for spildt;
han antager ikke, at Noget derved er tabt for hans
fredelige Liv og Sysler, thi han veed, at et Liv i Fred
og Ro er uden Værd, naar Selvstændigheden mangler.

Et lille Folk maa være som en fredelskende, men
tapper og veludrustet Riddersmand. Det spørger ikke,
hvor stor Sandsynligheden er for at seire, naar dets
Ære og Personlighed angribes. Det veed med sig selv,
at den sande Selvstændighed staaer uafhængig af Seier
og Nederlag. Saalænge Befolkningens Fleertal kun
følte sig som blindthen lydige Undersaatter under en
selvmægtig Stat, kunde det falde naturligt, at Borgerne
% --- p. 142 ---
\setcounter{footnote}{0}%
\opage{142}lode Staten om Forsvaret, betragtede Krig som en
reen Ulykke og Fred for enhver Priis, som et ubetinget
Gode. Kunde Staten ikke forsvare sig, hvad var
det saa Andet end, at man fik en ny Regjering; man
kunde altid være vis paa at beholde de gamle Byrder
og Plagerier. Det var den politiske Betragtning i
Fablen om Hesten og Æslet, som gjentog sig med en
lille Ændring. Anderledes, naar Staten er folkelig, er
Folkets Stat; en saadan Stat er et Udtryk af Folkets
virkelige Personlighed, og Personligheden maa forsvares.


Et stort og mægtigt Folk kan snart være krigerisk,
snart fredelig stemt; men det har, paa Grund af det
Masseagtige, vanskeligt ved at faae de modsatte Stemninger
til at gaae sammen til Eenhed. Hos et lille
Folk, som kun kan føre Krig for Fredens Skyld, maa
Folke-Aanden til enhver Tid i lige Grad være baade
krigerisk og fredelig. Det lille Folks Selvstændighed
er ikke, som det stores, en Gang imellem udsat for
Fare og saa igjen betrygget; nei, det lille Folk er,
væsenlig seet, altid i Fare, just fordi det er lille. Det
maa altid være aarvaagent, altid i Spænding. Men
netop i denne Aarvaagenhed, denne Spænding er der
en Velsignelse. Fredens Gjerninger og Forberedelserne
til Forsvar paa Liv og Død bør gaae Haand i Haand
med hverandre.

Det er ikke Overspændthed, ikke opskruet Ridderlighed,
som herved anbefales. Den opskruede Ridderlighed
er tvertimod paa deres Side, som uden Sands
for Selvforsvarets virkelige Betingelser, synge om Oldtidsbersærker
og gamle Helte, menende, at det egenlig
ikke kommer an paa Vaaben og Vaabenøvelser, men
% --- p. 143 ---
\setcounter{footnote}{0}%
\opage{143}alene paa „Ordet og Aanden.“ Det er Uforstand, % defect: „Utorstand“ (the head of the f did not ink, as pp. 52, 72, 157): the letter supplied
det
er huul Begeistring. At sende uøvede eller dog kun
halvøvede Soldater i Marken er jo, reent ud sagt, at
sende Fædrelandets Sønner blindt i Døden. Nei, saa
hellere slutte sig til dem, der tage Sagen materielt.
Det er dog altid Sands for Realitet, at arbeide paa
Befolkningens Ernæring; man kan jo tilsidst --- af Agtelse
for det Reelle --- give Folkeselvstændigheden en god Dag
og sætte Grødgryden høiere end Æren.

For en alsidig Gjennemførelse af Frihedens og
Lighedens Grundsætninger stiller den historisk nedarvede
Samfundsorden unegtelig store Hindringer. Men
at skrive de stedfindende Misforhold udelukkende paa
Statsindretningernes Regning og mene at kunne komme
Uligheden tillivs ved at kulegrave det Bestaaende,
er en uhyre Vildfarelse. Uligheden er mere end vilkaarlig;
den er dybt rodfæstet i de menneskelige Naturforhold.
Skulde den udslettes, da maatte det skee ved
en saa unaturlig Tvang og en saa ubarmhjertig Samfundsdespotisme,
at spartansk Strenghed vilde være
Børneværk i Sammenligning dermed. Hvor skarpt
Samfundsklasserne med den høist ulige Fordeling af
Goder og Rettigheder kunne staae imod hverandre,
endog paa Dannelsestrin, der mere tilhøre Naturen end
Kulturen, viser Historien i mangfoldige Exempler. Vi
skulle ikke her udbrede os over Kastevæsenet i Indien
eller i det gamle Ægypten, men holde os til vor egen
Fortid.

Det gammelnordiske Digt \emph{Rigsmaal} beskriver os
de forskjellige Stænders Oprindelse. „Guden Heimdal
havde engang under Navnet Rig begivet sig paa Vandring.
Han kom ind til et Par gamle fattige Folk,
% --- p. 144 ---
\setcounter{footnote}{0}%
\opage{144}Ae og Edda; der var usselt hos dem, Brødet var bagt
i Asken og fuldt af Saaer. Nogen Tid efter fødte
Konen en Søn, som blev overøst med Vand og kaldet
\emph{Træl}. Han var mørk af Farve, Huden paa hans
Hænder var grov, Knoerne tykke, Fingrene plumpe,
hans Ansigtsfarve gusten, Ryggen kroget og Hælene
lange. Han voxte op og lærte at bruge sine Kræfter.
Engang kom der en Kvinde til Huset med tung Gang,
solbrændte Arme og Vabler under Fødderne. Hun hed
\emph{Thy} d. e. Trælkvinde, og giftede sig med Træl. Deres
Sønner, som hed Røgter, Kludrer, Plump, Dovning,
Nøler, Krogryg og Slæber, beskjæftigede sig med at
sætte Hegn, gjøde Marker, vogte Sviin og Gjeder og
skjære Tørv. Døttrene hed Sorte, Klunte, Krognæse,
Krumryg og Langbeen. Af disse er Trællenes Æt oprunden.
Rig vandrede videre og kom til Bondens
Gaard, til et Par lykkeligere Folk, Ave og Amma. En
Ild var antændt midt paa Gulvet. Manden sad og
tilsnittede Træ til en Væv; hans Skjæg var redet og
Haaret klippet i Panden, Skjorten var snæver og aaben
i Brystet. Huusmoderen snoede Tenen og tilberedede
Garnet til Væven; hendes Haar var snoet i en Top,
Trøien stod aaben i Brystet. og hun havde en Krave % print: „Brystet.“ (a stop before a lower-case word)
om Halsen og Knapper om Skuldrene. Konen fik en
Søn, der blev svøbt i Lærred, overøst med Vand og
kaldet \emph{Karl} d. e. Fribaaren. Det var en blomstrende
Dreng lys af Haar og med tindrende Øine. Han voxte
til, lærte at tømre, smede, og kjøre Ploven. En Pige,
hvis Navn var \emph{Flink}, kom agende til Gaarden i Gedeskinds
Kjole med Nøgleknippet ved Siden. De giftede
sig og levede fornøiet sammen. Deres Sønner vare
Mand, Svend, Selveier, Bonde, Huusbonde, Langskjæg
% --- p. 145 ---
\setcounter{footnote}{0}%
\opage{145}og Bravmand. Døttrenes Navne vare Ordsnild, Brud,
Sneehvide, Blufærdig, Stadskjær, Myndig og Kvinde.
Rig kom derpaa til en Bygning, hvis Dør vendte mod
Syd; i Døren, som var lukket, sad en Ring, og Halm
var strøet paa Gulvet. Manden og Huusmoderen sade
og betragtede hinanden og legede med Fingrene. Manden
syslede med Vaaben, og Konen bar Ringsmykke
paa Brystet; hendes Pande var klarere og hendes Hals
hvidere end den nyfaldne Snee. Ni Maaneder efter
Rigs Besøg fødte hun en Søn, der blev svøbt i Silke,
overøst med Vand og kaldet \emph{Jarl}; han havde lyst Haar,
og hans Øine tindrede som en Slanges. Han lærte
efterhaanden som han blev stor, at skyde med Bue,
ride, svinge Sværdet og tumle sig i Vandet. Og Rig
kom til ham, lærte ham at kjende Runer og bød ham
at vinde Lande. Saa drog \emph{Jarl} ud, færdedes i Krig
og vandt atten Gaarde som sit Eie. \emph{Erna} (den Raske),
Herses Datter, som var smal i Midie og hvidere og
stoltere end andre Møer, blev hans Hustru, og herlige
Sønner smykkede deres Ægteskab.“

Et Samfund med saa indgroede, faste, naturbundne
Uligheder er, hvad Exemplet med Trællene angaaer,
ikke langt fra det Umenneskelige: skulde noget Lignende
ikke, om end under andre Betingelser, kunne
siges med Hensyn paa det i vore Dage stadig voxende
Proletariat? I ethvert Samfund, det være nok saa
demokratisk, vil der altid gives en Ulighed, som enten
allerede er eller dog staaer paa Spring til at blive oprørende:
mod Ulighedens Overgreb maa Folket uafladelig
være paa sin Post. Men i Ulighedens Væsen er der
en Hemmelighed.

Ligheden bliver umenneskelig, naar man vil af%
% --- p. 146 ---
\setcounter{footnote}{0}%
\opage{146}skaffe Uligheden; Uligheden bliver umenneskelig, naar
man ikke vil gjennemføre Ligheden: \emph{det er paa Ulighedens
Grundlag, at Ligheden maa gjennemføres}. Samfundets Kvaler hidrøre fra, at denne simple
Sætning bliver misforstaaet fra to modsatte Sider.

De, som have Interesse af at opretholde Uligheden,
holde sig eensidig til det Naturnødvendige med alle
deraf flydende Vilkaarligheder. At Samfundsstigen har
mange Trin, at der paa Grund af Arbeidets Deling nødvendig
maa gives forskjellige, høiere og lavere, Arbeidsklasser,
at Eiendomsrettens Ukrænkelighed, Formuens
Arvelighed, forskjellig Dygtighed, forskjellig Lykke o. s. v.,
maa frembringe Uligheder, som ingen menneskelig
Haand formaaer at udviske, om alt dette vide Forsvarerne
af det Bestaaende fortræffelig Besked. Men
naar saa Sagen faaer det Udseende, at Uligheden uden
videre er det Solide, det rette Fornuftige, hvorved man
trygt kan blive staaende; naar det atter og atter forglemmes,
at Uligheden, fastholdt for sig, er umenneskelig,
at Uligheden kun er til, for at Ligheden i stedse
rigere Maal kan virkeliggjøres: da bliver Anskuelsen
ufornuftig, og det Parti, som strider under dens Fane,
umedgjørligt.

De, som lide under Ulighedens Tryk og desaarsag
ere interesserede i at gjøre Ligheden gjældende, have
vægtige Grunde at anføre. „Kan det kaldes en Frihedens
og Lighedens Tid, naar man tilsikkrer de Besiddende
at lade den utrættelige Kapital arbeide for sig, voxe,
uopslidelig under Arbeidets Fremgang, medens tusinde
og tusinde virkelige Arbeidere maae slide sig op uden
at naae det Nødtørftige, fordi det økonomiske Spørgsmaal
ikke er om Arbeideren, men om Arbeidet, fordi
% --- p. 147 ---
\setcounter{footnote}{0}%
\opage{147}Kapitalen, dette Alt blodsugende Upersonlige, er blevet
Arbeide, og Arbeidet Kapital? En Hest kan man vurdere
og behandle efter dens Arbeide, thi Arbeidet er
ikke for Hestens, men Hesten for Arbeidets Skyld.
Skal den offenlige Moral maaskee forherliges ved den
statsøkonomiske Profit: at sætte et Menneske, der arbeider,
paa lige Trin med en Hest? Kan det kaldes
Agtelse for Menneskets Personlighed, naar golde Abstractioner,
almene Frembringelser af virkelige Menneskers
Anstrængelse i deres Ansigts Sved, --- thi Andet
er jo Kapital og Arbeide ikke --- skulle have Ret til
at bruge de arbeidende Personer som levende Maskiner?
Er det maaskee den Forret, Loven tilkjender det Levende
fremfor det Livløse, at Kapitalen er forpligtet til at
reparere de døde Maskiner, efterhaanden som de slides,
medens de levende kunne kastes bort, naar de ikke
længer ere i Stand til at reparere sig selv?“ Her er
Spørgsmaal, som naar de fremsættes i praktisk Form
med Sigte paa det Enkelte, vel kunne sætte Statsøkonomer
i Forlegenhed og bringe det civiliserede
Samfund med sine „Menneskerettigheder“ til at rødme.
Vi kunde skjænke Lighedens Forkæmpere vor udeelte
Sympathie, dersom det blot var muligt at faae den
Vildfarelse udryddet, som svækker og hæmmer den
hele Bestræbelse. Grundvildfarelsen ligger i, at man,
istedenfor at erkjende Ulighedens Betydning som Grundlag
for Ligheden, vil have Uligheden selv rykket op
med Rode. En sørgelig Vildfarelse! Den Lighed, der
ikke formaaer at optage Ulighedselementet i sig, er og
bliver utopisk.

Hos de Dyr, der leve i Selskab, har Naturens
Herre sørget for, at Individerne af forskjellige Klasser,
% --- p. 148 ---
\setcounter{footnote}{0}%
\opage{148}hvor saadanne findes, nøiagtig udføre det til hver Klasse
svarende Arbeide og i Overeensstemmelse hermed erholde
deres Andeel af Samfundets hele Erhverv. Myrer,
Termiter, Bier o. s. v. afgive Exempler paa, hvorledes
Uligheden paa naturlig Maade lader sig forene med
Ligheden. Foreningens Vilkaar er at hver Klasse af
Individer har af Naturen et saa eiendommeligt Præg,
at den indenfra drives til at udføre sit eiendommelige
Arbeide og desaarsag befinder sig i en saadan Afhængighed
af det hele Samfund, at den umulig kan løsrive
sig fra det. Det er en saadan Afhængighed af særegne
Naturanlæg, hvortil Skildringen i det anførte
„Rigsmaal“ peger hen. Trællen, Bonden, Jarlen have
alle deres Oprindelse fra Gudens Besøg; men deres
Naturbegavelse er saa forskjellig, at man endog vil kunne
see det ligefrem paa deres Ydre, til hvilken Klasse enhver
iblandt dem oprindelig hører. Denne Forklaring
er grundfalsk; Uligheden faaer her en saa forfærdelig
Førstefødselsret, at Ligheden umulig kan komme frem.
Vel er der lavere og høiere Naturbegavelse, men den
gjælder Individerne, ikke Klasserne; Historien har viist,
at der lige saa vel kan gives kløgtige Trælle som dumme
Jarler. Fra et andet, mere ophøiet Stade har man
forsvaret Uligheden ved at beraabe sig paa Guds udtrykkelige
Villie. Det skal da, ifølge den officielle
Religion, være ganske afgjort, at Gud selv har villet
en saadan Stændernes Adskillelse, at en Stand er sat
i Høihed, en anden i Ringhed, at nogle Mennesker ere
af Forsynet bestemte til Slid og Armod, andre til
Nydelse og Rigdom. Frugterne af denne officielle Fortolkning
have viist sig i forskjellige Lande. Det forsømte
Proletariat, hvis Uvidenhed om det Aandelige
% --- p. 149 ---
\setcounter{footnote}{0}%
\opage{149}paa mange Steder er grændseløs, har, forstaaeligt nok,
sammenblandet det ydre Kirkevæsen med Religionen
selv og ladet sig det bilde ind. % print: „ind.“ (a stop before a lower-case word)
at Troen paa en levende
personlig Gud var det Samme som Troen paa en overmenneskelig
Despot, der tog Parti med de Store og
hjalp dem med at underkue de Smaa. Denne Indbildning
maa fra Grunden af udryddes, dersom det i
nogen Grad skal lykkes at gjennemføre Ligheden.

Det for Christusreligionen Eiendommelige er, at
den lægger Vægt baade paa Uligheden og paa Ligheden;
den lader ikke blot den stedfindende Modsætning mellem
Høie og Lave, Mægtige og Ringe, Herrer og Tjenere,
Rige og Fattige bestaae, den gaaer videre og henfører
al denne Ulighed til en guddommelig Orden. Men det
maa for Alting ikke oversees, at den netop betragter
den forudsatte Ulighed som en Betingelse for at gjennemføre
Ligheden. Den taler til Herre og Tjener, til
Rig og Fattig, og hvad siger den? „Husker paa, at
I ere lige for Gud, at I ere Brødre.“ Paa dette Vilkaar
har Tjeneren lige saa vel Forpligtelse mod Herren,
som Herren imod Tjeneren, den Fattige lige saa vel
Forpligtelse mod den Rige, som den Rige mod den
Fattige. Men Religionens Aand kjender Verden og veed,
hvor haard Tilværelsen er; just derfor lægger den to
Gange, tre Gange Vægt paa den Riges Forpligtelse
mod den Fattige, to Gange, tre Gange Vægt paa Herrens
Forpligtelse mod Tjeneren, paa Arbeidsgiverens
Forpligtelse imod Arbeideren. Vidste Du fattige Arbeider,
hvo Han er, der har talt og endnu fremdeles
taler din Sag for Gud og Mennesker og med en saadan
Myndighed, at dit personlige Værd, din ufortabelige Ret
skal staae urokkelig fast, at Du i Aand og Væsen kan
% --- p. 150 ---
\setcounter{footnote}{0}%
\opage{150}være Ligemand med Jordens Rige og Mægtige: i Sandhed,
Du glemte ham aldrig!

Med den religiøs-sædelige Erkjendelse af Menneskets
Værd maa den politisk-borgerlige Udligning af de Misforhold,
som Uligheden indfører, uafbrudt tilstræbes.
Det bør fra begge Sider, lige saa vel fra deres, der
føle sig trykkede, som deres, der ere begunstigede,
indrømmes, at Ulighedsknuderne ikke lade sig løse,
som man løser en Sløife. Der maa fra den ene Side
være Villighed til at bringe Offere og paa den anden
Besindighed nok til at give Tid og tage til Takke med
Noget, indtil der kan opnaaes Mere. Hovedsagen er,
at der for Bestræbelserne i det Hele bliver lagt en
Plan, som begge Parter maae erkjende for hensigtssvarende;
at der arbeides sammenhængende Skridt for
Skridt, saa at det Foregaaende forbereder det Efterfølgende;
at der ad den borgerlige Lovgivnings Vei
virkelig udrettes Noget. Her, om ellers nogensteds,
kan det da sees, hvor inderlig Sindelag og Handling
høre sammen. Mangler Sindelaget, saa ere alle praktiske
Forslag forgjæves; mangler Handlingen, saa kan
man ræsonnere og moralisere saa meget, man vil: det
fører til Intet.

See vi os nu omkring efter de gunstigste Betingelser
for en virkelig Løsning af Lighedens Opgaver,
mon det da ikke skulde vise sig, at et lille Folk har
væsenlige Fordele fremfor et stort? Hos et stort Folk
ville de forskjellige Samfundsklasser ikke alene holde
sig adskilte; men hver Klasse vil endog paa Grund af
Massernes Storhed, Individernes Mængde, udgjøre en
egen Verden for sig. To Klasser af samme Befolkning
kunne i Levesæt, Omgangssprog, Sæder og Vaner staae
% --- p. 151 ---
\setcounter{footnote}{0}%
\opage{151}hinanden fjernere, end to Personer af forskjellig Nationalitet,
naar de befinde sig paa samme Dannelsestrin.
Vistnok ville Blandinger, Overgange og Sammensmeltninger
finde Sted paa mange Maader; men Ulighedselementerne
med det Fremmedartede, det Fjernende
og Frastødende dukke op under saa mange
Former og holde sig i en saa haard Umiddelbarhed,
at den ethisk-praktiske Gjennemarbeidelse bliver, om
ikke umulig, saa dog overordenlig vanskelig. „Der er,“
siger en aandrig fransk Forfatter, „to Folkefærd i
Frankrig: Provindsen og Paris; det ene spiser, sover,
gaber, lytter til; det andet tænker, vover, vaager og
taler. Det første drages med af det andet som en
Snegl af en Sommerfugl, idet det afvexlende mores og
foruroliges af sin Førers Lunefuldhed og Dristighed.
Det er da Føreren, vi maae see hen til.“ De sidste
Aars Begivenheder have i sørgelig Maade viist, at den
moralske, politiske og sociale Splid mellem de særlige
Samfundsklasser i det store Land ingenlunde indskrænker
sig til den her angivne Modsætning mellem
Provinds og Hovedstad; den grener sig saa vidt og
skjærer saa dybt i Samfundets Rødder, at det vist har
lange Udsigter, inden Ligheden i Menneskerettighedernes
Fædreland bliver forligt med Uligheden og de stridende
Parter forsonede.

Hos os ere Forholdene anderledes gunstige, at sige,
naar Folkevillien alvorlig vil benytte dem. Det er en
Kjendsgjerning, at jo mindre et Selskab er, desto større
Vægt kan der lægges paa dets enkelte Medlemmer.
Netop den Omstændighed, at Samfundsklasserne hos
os ere for smaa til at udgjøre egne afsluttede Verdener,
maa begunstige en Tilnærmelse, en saa at sige per%
% --- p. 152 ---
\setcounter{footnote}{0}%
\opage{152}sonlig Forstaaelse imellem dem. Et Individ af hundrede
er, i Egenskab af Individ, ti Gange saa meget
værd som et Individ af tusinde, og ti tusinde Gange
saa meget værd som et Individ af en Million, ifølge
almindelig Masseberegning. Men hvad er en Masse?
Ordet Masse er netop betegnende for Materien, det
Døde, det Uorganiske, Livløse. Materien bestaaer af
udeelbare Dele, hvis samlede Mængde udgjør Massen.
Jo flere Dele, desto større Masse, jo større Masse,
desto større Tiltrækningskraft. Kloderne i Rummet
virke i Forhold til deres Masser, idet de paavirke hverandre
efter Tiltrækningens Lov: i denne Region ere de
store Masser Stormagter. Overføres denne Betragtning
paa de menneskelige Samfund, saa blive Menneskeindividerne
Massedele, og Loven gjentager sig: jo flere
Individer, desto større Masse, jo større Masse, desto
større Tiltrækningskraft. Men hvo seer nu ikke det
Farlige, det for Folkelivet Fordærvelige i, at den døde
Tiltrækningslov anvendes paa personlige Væsener. Jo
mindre Selvstændighedsfølelse der er hos de enkelte
Individer, desto stærkere Tilbøielighed have de til at
flokke sig i Mængde og klumpe sig sammen i Masse.
Paa disse Betingelser vil det store Folk altid udøve
en tyngende, trykkende Indflydelse paa det lille. Ikke
alene i Krig, ogsaa i Fred. En fjantet Mode, en tankefattig,
men slibrig glimrende, nervepirrende Roman,
f. Ex., gaaer ud i Paris, og hvad skeer? Det nye
Phænomen bliver ikke bedømt efter sit indre Værd;
langt fra! Om ogsaa Kritiken gjør sin Pligt, det hjælper
ikke. Det er forslugne Masser, der holde Vægtskaalen,
Masseindividerne, der fælde Dommen. Romanen kommer
ud i syv Oplag; den store Nation, hedder det, har godkjendt
% --- p. 153 ---
\setcounter{footnote}{0}%
\opage{153}dens kulturhistoriske Betydning. Hvad følger heraf?
„At en Oversættelse naturligviis er nødvendig for at
berige vor Litteratur; vi tør ikke staae tilbage“: det
lille Folk er ganske imponeret.

Et lille Folk kan ikke slaae sig igjennem ved
Hjælp af Masser; det er af Forsynet just blevet stillet
under saadanne Vilkaar, at den sunde Sands maa gjøre
det opmærksomt paa, hvilken Forskjel der er mellem
Folk og Masse. Massen som Masse er uorganisk og
for saa vidt altid raa; i Folket som Folk er Alt organiseret,
levende, personligt: jo mere Massen breder
sig, desto mere indskrænkes Folket; jo mere Folket
udvider sig, desto mere indskrænkes Massen; eller:
jo mindre Masse, desto mere Folk, jo mere Masse,
desto mindre Folk. Men at være sat under saadanne
Vilkaar, at der for Aandsvillien netop er den allerstærkeste
Opfordring til at faae Individerne udviklede
til personlig Selvstændighed, faae Massen organiseret
og forvandlet til et levende Folk, faae det Hele gjennnemtrængt, % print: „gjennnem-“ (misprint)
beaandet og bevæget af een Tanke, een Villie,
eet Sind, deri er der da ingen Forbandelse; nei tvertimod,
den største Velsignelse.

Vort Fædreland, som i tidligere Dage stundom har
været Undergangen nær, men dog paa vidunderlig
Maade er blevet frelst, har navnlig i dette Aarhundrede
lidt en haard Medfart. Under alt dette have vi dog
beholdt et Gode, som er uskatteerligt. Der er i Historien
to store Love, som, efterdi de ere Aands- og Frihedslove,
under Endelighedens Brydninger saa let komme
i Strid med hinanden, men dog bestandig maae bringes
til Overensstemmelse, det er: \emph{Fortsættelsens}
Lov, en Lov for den jævne Udvikling, og Frem%
% --- p. 154 ---
\setcounter{footnote}{0}%
\opage{154}\emph{skridtsloven}, en Lov for de nye Ideers Gjennembrud.
I ethvert afgjørende Vendepunkt viser der sig
en Spænding mellem disse to Love. I Frankrig, i
„den store Nation“ har Gjennembrudsloven faaet saa
afgjørende Overvægt, at Legitimitetens Grundlag er
undergravet. Enhver Samfundsidee, som er i Stand til
at opflamme Mængden og samle Masserne omkring sig,
har, efterdi den historiske Fortsættelseslov er brudt, i
Grunden samme Ret, som det for et Øieblik Bestaaende.
Statsskibet er hullet i Kjølen og maa bestandig holdes
paa Pumperne. Hvorledes kan Folke-Eenhed komme
til Kræfter i et Land, hvor tre Dynastier og to Republiker,
samtidig i fuld Virksomhed, kæmpe med hverandre
paa Liv og Død, idet enhver af disse Magter ---
efter Omstændighederne --- anseer sig for fuldt berettiget
til at gribe Regjeringens Tøiler. Vi skulle
ikke hovmesterere „den store Nation“; vi skulle med
Beundring erkjende, at den, selv under sine Nederlag,
har udfoldet forbausende Kræfter; vi skulle --- for
„Menneskerettighedernes“ Skyld --- troe og haabe, at
det Folk, som under forfærdelige Kriser og Udskeielser
har formaaet at holde Frihedsbanneret oppe, vil fortsætte
sin Gang og naae sit Maal til Hæder for sig
selv og til Bedste for den civiliserede Verden. Men
naar vi betænke, hvad det vil sige, at vove saadanne
Brud paa Fortsættelsens (Continuitetens) Lov, som her
er skeet, da maae vi prise vor Lykke.

Det Mærkelige ved vor Forfatning af 1848 er netop
dette Dobbelte: at Forandringen er saa gjennemgribende,
og dens Indførelse dog saa lovlig. Det kaster
historisk Glands paa Frederik den Syvende og hans
Mænd, at Fortsættelse og Gjennembrud her vise sig
% --- p. 155 ---
\setcounter{footnote}{0}%
\opage{155}som to Sider af det Samme. Det er et Fodfæste for
dem, der holde paa det Bestaaende; den lovlige Vei,
den jævne, rolige Overgang er slaaet fast som Princip.
Men det er ogsaa en Støtte for dem, som kjækt og
freidigt arbeide for Fremskridt; thi af Principet sees
det jo, at man --- forstaaer sig, i sindig Alvor --- kan
tage en Sag fra Roden af, uden at afstedkomme en
voldsom, lovløs Omvæltning.

Saa have vi da, hvad der er en Velsignelse over
vort lille Folk og gjør „den tunge Eenhed“ let, med
Forsynets Bistand alle Betingelser for en driftig, men
rolig Udvikling af et sundt og kraftigt Folkeliv i
Danmark.

Theoretiske Overblik tør ikke forblandes med de
praktiske Opgaver; men saa vist som Sindelag og Handling
uafbrudt betinge hinanden, saa nødvendigt er det, at
Folket tilstrækkelig oplyses om dets egne Vilkaar ved
at oplyses om Menneskelivets Vilkaar i det Hele. Det
er Folkelærere, vel forberedte, vel udstyrede Folkelærere,
hvorpaa det i denne Henseende først og sidst
kommer an. Vi slutte med at minde om den store
Folkelærers egne Ord: \emph{Høsten er stor}, men \emph{Arbeiderne
ere faa; beder} Høstens Herre, at han
\emph{udsender Arbeidere i sin Høst} (Luk. 10, 2).

% --- p. 156 ---
\setcounter{footnote}{0}%
\lecture{Til Minde om Istedslaget.} {(i Aagaards Skov).}
\opage{156}Idag for fem og tyve Aar siden var der Høitid her i
Landet; da blev der holdt en Fest, --- en blodig Fest
paa Isteds brune Hede; det var en Offerfest i gammelnordisk
Stiil. Vi vide jo alle hvorom det gjaldt. Det
gjaldt en Dyst mod Overgreb af trodsigt Hovmod, af
Rænkespil, som bryder Tro og Love; det gjaldt at
værge om vort Grændseskjel, vor gamle Ret, vor unge
Frihed, vort Moderlands Selvstændighed og Ære. At
træde Oprørsøglen under Fod og løfte Dannebrogen
høit var Kampens Maal; og Maalet maatte naaes, thi
det var Kongens, det var Folkets og det var Soldatens
Villie. Saa gik det frem; trods Baghold, Kugleregn og
Damp og Blod, ja selv trods Lykkens Luner, gik vore
Tappre freidig frem til Maalet. Det var en Hædersdag
for danske Hjerter og for danske Løver. Schleppegrell
og Læssø! disse Navne har Saga ristet paa sit Skjold;
de straale der med ufordunklet Glands til Ihukommelse
hos et taknemmeligt Folk. Dog var det mere end
Navnkundighed, hvortil der beiledes. Den Helt, som
stred den store Strid, saae ei for høit, tog ikke Sigte
paa Berømmelsernes Stjerner; det var en djærv beskeden
% --- p. 157 ---
\setcounter{footnote}{0}%
\opage{157}Helt: Fanens Hæder, Pigens Smiil og Kongens Tak var
al hans Løn; det var den tappre Landsoldat, som
bragte her sit Høitidsoffer.

En herlig Seier uden Brask og Bram. Der var i
Glædens Jubel intet Overmod, ak nei, et stille Suk, et
taareblændet Veemod. Seiersefterretningen slog ned
som Lyn; men, saa fulgte med --- den dyre Regning,
den lange, den forfærdelige Liste. Enhver, som har
oplevet det, veed vel, hvad det vil sige: i Sjæleangst,
med bankende Hjerte og skjælvende Haand at gjennemløbe
Listen over Saarede og Faldne; det sortner En
for Øiet ved Tanken om at finde her --- ak, dyrebare
Navne. Af Glædens Farvespil og Sorgens Skygger har
Tiden med lemfældig Haand lidt efter lidt bragt Ligevægt
tilveie og malet os et Billede i Mindet.

Paa Istedseieren er desværre fulgt et Nederlag, et
Tab, som smerter; men Æren er uplettet. Overmagten
brød igjennem og sønderskar vort Land; men Dannebrog
er heel endnu, i denne Dug er ingen Rifter. At
binde an med Een, at staae for To, at vige med Foden
lidt for Tre og uden Skam at flye for % defect: „tor“ (the head of the f did not ink): the letter supplied
Fire, det var
Kong Frodes Lov om Mandemod, og denne Lov har
Landsoldaten trolig efterlevet. Betragte vi det Overtal,
mod hvilket vi i sidste Kamp have maattet stride, da
er vort Forsvar os til Hæder; thi her var mere end
To og Tre mod Een, og vore Fjender skulle, efter
Ridderskabets Love, just ikke rose sig af Seieren.

Endnu er Kampen staaende; endnu er der en
Seier for vor Selvstændighed at vinde, dog med andre
Vaaben og paa andre Vilkaar. Selvstændigheden er
jo af forskjellig Art, deels noget Udvortes, deels noget
Indvortes. Et stort Land med et stort Folk kan holde
% --- p. 158 ---
\setcounter{footnote}{0}%
\opage{158}sig selvstændigt udadtil, saa længe dets Vaabenmagt
formaaer at kue og beseire de udvortes Fjender; men
dermed er det ikke sagt, at det ogsaa indadtil er et
selvstændigt Folk; langt fra! Det er i Historien ofte
seet, at et saadant Folk kan være i megen Splid og
Tvedragt med sig selv, at det netop for Storhedens og
den ydre Selvstændigheds Skyld maa bære tunge Lænker.
Ogsaa et lille Land med et lille Folk kan udadtil bevare
sin Selvstændighed, dersom Naturen med sjælden
Gunst, ved steile Kyster, høie Fjelde utilgængelige Egne,
har sikkret det mod Overfald; men, selv omgivet af
en chinesisk Muur, kan et saa betrygget lille Folk dog
indadtil være høist uselvstændigt; det kan spalte sig i
uforsonlige Partier, det kan fortære sine Kræfter i et
tomt, ufrugtbart Røre, det kan oprive sig selv i smaaligt
Kiv. For den udvortes Selvstændighed tør ingen
menneskelig Magt indestaae; selv den Stærkeste kan
snuble: store Riger, store Throner falde. Over Folkenes
Skæbne i Tidernes Løb raader et styrende Forsyn, og,
det være vor Trøst, Forsynets Øie seer lige saa mildt
til de Smaa som til de Store.

Det gjælder da om at vinde den indre Selvstændighed;
men til indre Selvstændighed hører Enighed,
indbyrdes Enighed i Sandheds, Retfærdigheds, Friheds
og Kjærligheds Aand. Folkelivet er ikke et Liv i
Paradiis; den folkelige Jordbund maa renses, pløies
og dyrkes, det koster Arbeide. Omkring Frihedstræet
er der Roser og Kjærminder, ja! men ogsaa Torne og
Tidsler. Det er umuligt, at Selvstændighed og Frihed
kan udvikle sig hos et kraftigt opvoxende Folk uden
indvortes Spænding og Strid. Saa staaer Syn imod
Syn, Parti imod Parti, og modsatte Bestræbelser brydes
% --- p. 159 ---
\setcounter{footnote}{0}%
\opage{159}med hverandre. Det er nødvendigt, at der under slige
Brydninger maa være et stærkt Enheds og Enigheds
Grundlag. Brister Enighedens Grundlag, da er Folkets
Selvstændighed i Fare.

Hvo iblandt os har ikke erfaret, hvor vanskeligt
der er at bringe to Mennesker, om de end, hver for
sig, oprigtig ville det Rette, til at komme overeens,
naar der er Mistanke og hemmeligt Nag imellem dem.
I de offenlige Forhold, naar Partierne stride, kan det
ofte forekomme en Udenforstaaende, som om den klare
Sag, det rette Syn stod mæglende midt i Striden; men
med Eet er det igjen forsvundet, et Men, et Sideblik,
et enkelt Ord har jaget Forstaaelsen bort: det er ikke
uredelig Hensigt, ikke bevidst Underfundighed, det er
Nagets Orm, det er Mistankens Indskydelse. Hvad
skal rense Mistanken ud og kvæle Nagets Orm? En
inderlig stærk og levende Følelse af, hvad Folket er,
og hvad Fædrelandet kræver, naar store Interesser
staae paa Spil. Ogsaa i denne Henseende kan Mindet
om Istedslaget være os til Velsignelse. Der var Enighed,
ja der var der Eenhed, een Aand, een Sjæl, eet
Løsen. „Med Gud, for Konge og Fædreland!“ det var
et Ord, der gav Soldaten Mod; det boede i hans tause
Sind, det gav ham Vished om sit Kald, det tændte
Flammen i hans Bryst, det lød med Hurraraab i Skov
og Dal, det bruste hen over Heden. Kunde vi see
disse Offere, disse Døde, disse Døende, see Blodet risle
ud fra Hjertets Kilder, uden at glemme vor Hverdagskiv
og føle os opløftede over smaalig Egensindighed,
Mistanke og Nag, da vare vi et utaknemmeligt Folk;
Krigernes Minde, Aanden fra Heltenes Grave maatte
% --- p. 160 ---
\setcounter{footnote}{0}%
\opage{160}da vidne imod os, at vi gav vort bedste Eie, vor indre
Selvstændighed til Priis, og ikke vare Friheden værd.

Om Folkets og Fædrelandets store Anliggender er
Enighed i det Hele mere end en Følelsessag og mere
end et beregnende Forstandshensyn; det er en Villiessag;
der skal, der maa være Enighed om det Allervigtigste.
Den Mand, der ikke kan blive enig med sig
selv, bliver aldrig selvstændig, og det Rige, der er
splidagtigt med sig selv, maa falde. Til Enighed hører
Villie; til et enigt og selvstændigt Folk en klar og
moden Folkevillie. Foranstaltninger til Fædrelandets
Forsvar, de første Betingelser for Folkelivets alsidige
Udvikling i aandelig som i materiel Henseende høre til
de Ting, som ikke taale Omsvøb eller Forhaling; her
er Noget, som ikke kan vente til alle Krogveie ere
forsøgte, og alle Modstandsmidler hos de stridende
Parter udtømte.

Skal Selvstændighedsværket ikke gaae i Staa, maa
der holdes vaagent Øie med en dobbelt Fare. Det Ene
er: at give for meget efter og lade sig blænde af tomme
Løfter og intetsigende Overeenskomster. Det Andet er:
at trænge for stærkt paa og, koste, hvad koste vil,
sætte, som det hedder, sin Krig igjennem. Det Sidste
er lige saa fordærveligt som det Første. Ved at bruge
for stærke Midler kunne Folkets Ledere med al deres
Iver og gode Villie let komme til at bære sig ad som
de Kvaksalvere, der, for i Hast at kurere en enkelt
Skade, fordærve det hele Legeme og dræbe Patienten.
Ogsaa i denne Henseende have Fredens Mænd adskilligt
at lære af Krigerne. --- See til Høire! See til
Venstre! --- Lige ud! det hører til de unge Soldaters
første Øvelser. Dersom Sandhedens Vei laa aaben for
% --- p. 161 ---
\setcounter{footnote}{0}%
\opage{161}os, som en indgrøftet Landevei, behøvede vi ganske vist
hverken at see til Høire eller Venstre; men da vi bestandig
maae kæmpe os frem; da Veien kan bugte
sig gjennem Vildfarelsers Uføre, og Sandheden ligge
snart til Høire, snart til Venstre, saa gjælder det at
have aabent Øie for begge Muligheder: et simpelt travende
„Lige ud!“ kunde let bedrage os ved Indbildninger
om egen Ufeilbarhed. Og endnu Et. I
Slagets Tummel lyder et dobbelt Kommando-Ord, snart
et: Fremad, gaa paa! snart igjen et Holdt, staa! Det
sidste er for den til Kamp opflammede Soldat ofte
haardere end det første. At gaae paa, naar Alt er sat
ind, og Lidenskaben giver Fart, falder naturligt; men
at staae stille midt i en Kugleregn, see Kammeraterne storme
og falde, uden at turde ile dem til Hjælp, koster
Overvindelse. Ikke desto mindre er det, for Seierens
Skyld, lige saa vigtigt at parere den sidste Ordre som
den første. Til det sande Kampmod hører Selvbeherskelse;
skulde Nogen tvivle derom, lad ham raadføre sig med
Heltene fra Isted.

Naar et Folk i Alvor skal kæmpe for sin Selvstændighed,
er det ikke nok, at en Leder sætter sig
et Maal, samler en Flok omkring sig, et Fleertal, og
siger: vi ere Folket. Det er vel muligt, at en saadan
Fleertalsvillie kan være et første, begyndende Udtryk
for den sande og egenlige Folkevillie; ja, det er muligt,
men det maa dog prøves. Hvorpaa skal det da kjendes?
Det skal jo netop kjendes derpaa, at det Maal, der er
sat, ligelig omfatter Alt og Alle. At sigte opefter saaledes,
at det seer ud, som skulde det hele Folk særdeles
arbeide for nogle høiere Lærdomsanstalter, nogle
Vidende og Dannede, hvis Viden og Dannelse, det ikke
% --- p. 162 ---
\setcounter{footnote}{0}%
\opage{162}forstaaer, er usundt. At sigte nedefter, og det saaledes,
at det bestandig seer ud, som var det kun det Materielle,
det kjødelige Velvære, den timelige Fedme, der var det
ene Fornødne, er ogsaa usundt. I begge Tilfælde kunne
vi være forvissede om, at det er en Partivillie, vi have
for os, ingen sand Folkevillie. Hvor et Samfund lider
af den Grundskade, at et enkelt Parti, være sig et
Mindretals- eller et Fleertalsparti, raader eensidig, der
er Folkevillien lammet. Allermeest maae de smaa Folk
vogte sig for at forgude Fleertalsvillien; thi ved en
saadan Forgudelse underskrive de deres egen Dødsdom.
Har et Fleertal Ret, alene fordi det er Fleertal, hvad
følger saa heraf? Saa følger heraf, at, eftersom de
smaa Folk altid ere et Mindretal og de store et Fleertal,
saa har --- hvad en Stormagt gjerne vil paastaae
--- et lille Folk jo egenlig ingen Ret til at existere selvstændigt.


Som Maalet er, maae ogsaa Midlerne være; paa
Valg af Midler og paa Brugen af dem kjende vi igjen
Forskjellen mellem Partivillie og Folkevillie. Vi mindes
Kong Salomons Dom, der bød, at det levende Barn
skulde deles. Men den rette Moder vilde hellere overlade
sit Barn til den urette end taale, at det blev
deelt og dræbt. Viis os engang en Fleertalsvillie, som
uagtet den er sig sin Ret bevidst og kan bevidne det
høit, dog for Livets og Enighedens Skyld, lig Barnets
rette Moder, udviser den Selvfornegtelse at give efter
for et Øieblik, at intet Forhaandenværende skal krænkes
og intet nødvendigt Mellemled overspringes! Ja, naar
en Fleertalsvillie er saa høisindet, da er den sikker
paa sig selv; den vinder Modstandernes Agtelse, den
vinder mere end det, den vinder Folkets Hjerte: det
% --- p. 163 ---
\setcounter{footnote}{0}%
\opage{163}skal nok vise sig, at den er og bliver den rette, selvstændige
Folkevillie.

Idag for fem og tyve Aar siden holdt Folke-Aanden
Fest i Danmark. Forræderiet, som altid
speculerer i Parti-Aand, vilde splitte os ad og skære
et Stykke af vort eget Kjød; men Kongen hævede
sin Røst: \emph{det skal ei skee}! og den kongelige
Villie var en Folkevillie. Desværre, da denne Røst
forstummede i Døden, saa skete dog det, som vi længe
havde frygtet. Men Ordet fra de kongelige Læber er
ikke dødt; det maner os, at være paa vor Post, det
lyder end i vore Øren, en Aanderøst, et Spaadomsord,
stadfæstet af Guds Forsyn.

Der er en Selvstændighed, som ikke kan røves;
der er Noget i et ædelt Folk, som ikke skydes ned
med Kanoner. Stride vi endnu for Folkelivets Sag,
saa at vort Løsen er endnu som før: med Gud, for
\emph{Konge og Fædreland}! da lad kun dem, som stole
trygt paa egen Kraft, længe nok sætte Retten i Spydstagen;
lad dem kun mene, at Fleertalsovermagt er
nok til at tvinge et fremmed Sprog, et fremmed Væsen
og en fremmed Sjæl ind i et Folk, som ærlig har besluttet
at ville leve sit eget Liv: hver jordisk Overmagt
har dog sin Grændse; der er et Guddomsord,
som siger: hertil, ikke længere! Vi have ingen Ret
fremfor andre Folk, enten store eller smaa, til at kalde
os Guds eget Folk; men vi have baade Ret og Forpligtelse
til at føre vort Liv og kæmpe for vor Selvstændighed
saaledes, at vi med Grund kunne siges at
arbeide for en sindig og alsidig Udvikling af ædel og
sand Menneskelighed. Bliver vor Sag i dybere Forstand
en Menneskelighedens Sag, da tør vi troe og
% --- p. 164 ---
\setcounter{footnote}{0}%
\opage{164}haabe, at Sagen vil faae Fremgang, at ingen Overmagt
vil vove det Umenneskelige. Det er dog Menneskehedens
Sag og Menneskelighedens Seier, hvortil Historien
peger hen; men Menneskelighedens Seier er uadskillelig
fra et selvstændigt Folkeliv. Bort da med
bange Anelser og tunge Drømme! hver Trusel mod
vort Land og Folk skal veires hen; de skulle ei kue os,
„det skal ei skee“. Vi troe paa Gud og paa vor gode
Ret. Retfærdigheden har en mægtig Røst, og Ordet
giver Gjenlyd her i disse Dale; fra Mark og Skov, fra
Land og Hav, fra danske Hjerter og fra danske Løver,
fra Tonen i vor Sang, fra Stemmen i vort Bryst, fra
Himlens Hvælving og fra Istedheltes Grave gjenlyder
dette Ord: \emph{det skal ei skee}! Med denne Tro, i dette
Haab velsigne vi de Faldnes Minde, vi bringe de endnu
Levende vor Tak, og --- gamle Danmark leve!

\end{document}
